% Môn: Vật lý
% Lớp: 12
% Chương: Chương III. Từ trường
% Bài: L12C3 Bài 16. Từ thông. Hiện tượng cảm ứng điện từ
% Số câu: 324
% App đọc file này trực tiếp — sửa rồi Commit trên GitHub, bấm Đồng bộ GitHub .tex

% L12C3 Bài 16. Từ thông. Hiện tượng cảm ứng điện từ

% ===== Câu 1 =====
% ID: L12C3B16-01-DS
% Mức: TH
\dangbt{Khai thác đồ thị từ thông – thời gian và bài toán cảm ứng điện từ tổng hợp}
\begin{ex}
% ID: L12C3B16-01-DS
% Mức: TH
Xét Khai thác đồ thị từ thông – thời gian và bài toán cảm ứng điện từ tổng hợp (Tình huống 1). Hãy xác định tính đúng, sai của bốn phát biểu sau.
\choiceTF
{\True Độ lớn suất điện động cảm ứng bằng độ lớn hệ số góc của đồ thị liên kết từ thông theo thời gian.}
{Đồ thị từ thông nằm ngang ứng với suất điện động cảm ứng khác không.}
{\True Khi giải cần xác định đúng dữ kiện, phương chiều, điều kiện áp dụng và đơn vị.}
{Kết quả không phụ thuộc dữ kiện và có thể bỏ qua điều kiện.}
\loigiai{
\begin{itemchoice}
\item (Đúng) Độ lớn suất điện động cảm ứng bằng độ lớn hệ số góc của đồ thị liên kết từ thông theo thời gian. (Vì): ộ lớn suất điện động cảm ứng bằng độ lớn hệ số góc của đồ thị liên kết từ thông theo thời gian. b) Sai: Đồ thị từ thông nằm ngang ứng với suất điện động cảm ứng khác không. c) Đúng: Khi giải cần xác định đúng dữ kiện, phương chiều, điều kiện áp dụng và đơn vị. d) Sai: Kết quả không phụ thuộc dữ kiện và có thể bỏ qua điều kiện. Kiến thức cốt lõi: Độ lớn suất điện động cảm ứng bằng độ lớn hệ số góc của đồ thị liên kết từ thông theo thời gian
\item (Sai) Đồ thị từ thông nằm ngang ứng với suất điện động cảm ứng khác không.
\item (Đúng) Khi giải cần xác định đúng dữ kiện, phương chiều, điều kiện áp dụng và đơn vị.
\item (Sai) Kết quả không phụ thuộc dữ kiện và có thể bỏ qua điều kiện.
\end{itemchoice}
}
\end{ex}

% ===== Câu 2 =====
% ID: L12C3B16-01-TL
% Mức: TH
\begin{bt}
% ID: L12C3B16-01-TL
% Mức: TH
Trình bày kiến thức cốt lõi của Khai thác đồ thị từ thông – thời gian và bài toán cảm ứng điện từ tổng hợp và nêu điều kiện áp dụng.
\loigiai{
Cần nêu được: Độ lớn suất điện động cảm ứng bằng độ lớn hệ số góc của đồ thị liên kết từ thông theo thời gian. Khi vận dụng phải xác định đúng dữ kiện, phương chiều nếu có, điều kiện áp dụng, hệ đơn vị và kiểm tra tính hợp lí. Nhận định “Đồ thị từ thông nằm ngang ứng với suất điện động cảm ứng khác không.” sai.
}
\end{bt}

% ===== Câu 3 =====
% ID: L12C3B16-01-TLN
% Mức: TH
\begin{ex}
% ID: L12C3B16-01-TLN
% Mức: TH
Trong bốn phát biểu về Khai thác đồ thị từ thông – thời gian và bài toán cảm ứng điện từ tổng hợp, có bao nhiêu phát biểu đúng? (Chỉ ghi một số nguyên.) (1) Độ lớn suất điện động cảm ứng bằng độ lớn hệ số góc của đồ thị liên kết từ thông theo thời gian. (2) Đồ thị từ thông nằm ngang ứng với suất điện động cảm ứng khác không. (3) Cần kiểm tra điều kiện áp dụng. (4) Có thể bỏ qua đơn vị.
\shortans{2}
\loigiai{
Đối chiếu với kiến thức: Độ lớn suất điện động cảm ứng bằng độ lớn hệ số góc của đồ thị liên kết từ thông theo thời gian. Có 2 phát biểu đúng.
}
\end{ex}

% ===== Câu 4 =====
% ID: L12C3B16-01-TN
% Mức: NB
\begin{ex}
% ID: L12C3B16-01-TN
% Mức: NB
Phát biểu nào sau đây đúng về Khai thác đồ thị từ thông – thời gian và bài toán cảm ứng điện từ tổng hợp?
\choice
{Đồ thị từ thông nằm ngang ứng với suất điện động cảm ứng khác không.}
{Có thể bỏ qua phương, chiều và đơn vị trong mọi trường hợp.}
{Kết quả không phụ thuộc dữ kiện và điều kiện áp dụng.}
{\True Độ lớn suất điện động cảm ứng bằng độ lớn hệ số góc của đồ thị liên kết từ thông theo thời gian.}
\loigiai{
Độ lớn suất điện động cảm ứng bằng độ lớn hệ số góc của đồ thị liên kết từ thông theo thời gian. Vì vậy phương án $D$ đúng; các phương án còn lại sai.
}
\end{ex}

% ===== Câu 5 =====
% ID: L12C3B16-02-DS
% Mức: TH
\begin{ex}
% ID: L12C3B16-02-DS
% Mức: TH
Xét Khai thác đồ thị từ thông – thời gian và bài toán cảm ứng điện từ tổng hợp (Tình huống 2). Hãy xác định tính đúng, sai của bốn phát biểu sau.
\choiceTF
{Đồ thị từ thông nằm ngang ứng với suất điện động cảm ứng khác không.}
{\True Độ lớn suất điện động cảm ứng bằng độ lớn hệ số góc của đồ thị liên kết từ thông theo thời gian.}
{Kết quả không phụ thuộc dữ kiện và có thể bỏ qua điều kiện.}
{\True Khi giải cần xác định đúng dữ kiện, phương chiều, điều kiện áp dụng và đơn vị.}
\loigiai{
\begin{itemchoice}
\item (Sai) Đồ thị từ thông nằm ngang ứng với suất điện động cảm ứng khác không. (Vì): Đồ thị từ thông nằm ngang ứng với suất điện động cảm ứng khác không. b) Đúng: Độ lớn suất điện động cảm ứng bằng độ lớn hệ số góc của đồ thị liên kết từ thông theo thời gian. c) Sai: Kết quả không phụ thuộc dữ kiện và có thể bỏ qua điều kiện. d) Đúng: Khi giải cần xác định đúng dữ kiện, phương chiều, điều kiện áp dụng và đơn vị. Kiến thức cốt lõi: Độ lớn suất điện động cảm ứng bằng độ lớn hệ số góc của đồ thị liên kết từ thông theo thời gian
\item (Đúng) Độ lớn suất điện động cảm ứng bằng độ lớn hệ số góc của đồ thị liên kết từ thông theo thời gian.
\item (Sai) Kết quả không phụ thuộc dữ kiện và có thể bỏ qua điều kiện.
\item (Đúng) Khi giải cần xác định đúng dữ kiện, phương chiều, điều kiện áp dụng và đơn vị.
\end{itemchoice}
}
\end{ex}

% ===== Câu 6 =====
% ID: L12C3B16-02-TL
% Mức: TH
\begin{bt}
% ID: L12C3B16-02-TL
% Mức: TH
Giải thích vì sao nhận định sau đúng: “Độ lớn suất điện động cảm ứng bằng độ lớn hệ số góc của đồ thị liên kết từ thông theo thời gian.”
\loigiai{
Cần nêu được: Độ lớn suất điện động cảm ứng bằng độ lớn hệ số góc của đồ thị liên kết từ thông theo thời gian. Khi vận dụng phải xác định đúng dữ kiện, phương chiều nếu có, điều kiện áp dụng, hệ đơn vị và kiểm tra tính hợp lí. Nhận định “Đồ thị từ thông nằm ngang ứng với suất điện động cảm ứng khác không.” sai.
}
\end{bt}

% ===== Câu 7 =====
% ID: L12C3B16-02-TLN
% Mức: TH
\begin{ex}
% ID: L12C3B16-02-TLN
% Mức: TH
Trong bốn phát biểu về Khai thác đồ thị từ thông – thời gian và bài toán cảm ứng điện từ tổng hợp, có bao nhiêu phát biểu đúng? (Chỉ ghi một số nguyên.) (1) Đồ thị từ thông nằm ngang ứng với suất điện động cảm ứng khác không. (2) Độ lớn suất điện động cảm ứng bằng độ lớn hệ số góc của đồ thị liên kết từ thông theo thời gian. (3) Kết quả phải phù hợp ý nghĩa vật lí. (4) Mọi đại lượng đều là vô hướng.
\shortans{1}
\loigiai{
Đối chiếu với kiến thức: Độ lớn suất điện động cảm ứng bằng độ lớn hệ số góc của đồ thị liên kết từ thông theo thời gian. Có 1 phát biểu đúng.
}
\end{ex}

% ===== Câu 8 =====
% ID: L12C3B16-02-TN
% Mức: NB
\begin{ex}
% ID: L12C3B16-02-TN
% Mức: NB
Trong nội dung Khai thác đồ thị từ thông – thời gian và bài toán cảm ứng điện từ tổng hợp?
\choice
{\True Độ lớn suất điện động cảm ứng bằng độ lớn hệ số góc của đồ thị liên kết từ thông theo thời gian.}
{Đồ thị từ thông nằm ngang ứng với suất điện động cảm ứng khác không.}
{Có thể bỏ qua phương, chiều và đơn vị trong mọi trường hợp.}
{Kết quả không phụ thuộc dữ kiện và điều kiện áp dụng.}
\loigiai{
Độ lớn suất điện động cảm ứng bằng độ lớn hệ số góc của đồ thị liên kết từ thông theo thời gian. Vì vậy phương án $A$ đúng; các phương án còn lại sai.
}
\end{ex}

% ===== Câu 9 =====
% ID: L12C3B16-03-DS
% Mức: VD
\begin{ex}
% ID: L12C3B16-03-DS
% Mức: VD
Xét Khai thác đồ thị từ thông – thời gian và bài toán cảm ứng điện từ tổng hợp (Tình huống 3). Hãy xác định tính đúng, sai của bốn phát biểu sau.
\choiceTF
{\True Khi giải cần xác định đúng dữ kiện, phương chiều, điều kiện áp dụng và đơn vị.}
{Kết quả không phụ thuộc dữ kiện và có thể bỏ qua điều kiện.}
{\True Độ lớn suất điện động cảm ứng bằng độ lớn hệ số góc của đồ thị liên kết từ thông theo thời gian.}
{Đồ thị từ thông nằm ngang ứng với suất điện động cảm ứng khác không.}
\loigiai{
\begin{itemchoice}
\item (Đúng) Khi giải cần xác định đúng dữ kiện, phương chiều, điều kiện áp dụng và đơn vị. (Vì): Khi giải cần xác định đúng dữ kiện, phương chiều, điều kiện áp dụng và đơn vị. b) Sai: Kết quả không phụ thuộc dữ kiện và có thể bỏ qua điều kiện. c) Đúng: Độ lớn suất điện động cảm ứng bằng độ lớn hệ số góc của đồ thị liên kết từ thông theo thời gian. d) Sai: Đồ thị từ thông nằm ngang ứng với suất điện động cảm ứng khác không. Kiến thức cốt lõi: Độ lớn suất điện động cảm ứng bằng độ lớn hệ số góc của đồ thị liên kết từ thông theo thời gian
\item (Sai) Kết quả không phụ thuộc dữ kiện và có thể bỏ qua điều kiện.
\item (Đúng) Độ lớn suất điện động cảm ứng bằng độ lớn hệ số góc của đồ thị liên kết từ thông theo thời gian.
\item (Sai) Đồ thị từ thông nằm ngang ứng với suất điện động cảm ứng khác không.
\end{itemchoice}
}
\end{ex}

% ===== Câu 10 =====
% ID: L12C3B16-03-TL
% Mức: VD
\begin{bt}
% ID: L12C3B16-03-TL
% Mức: VD
Phân tích sai lầm trong nhận định: “Đồ thị từ thông nằm ngang ứng với suất điện động cảm ứng khác không.”
\loigiai{
Cần nêu được: Độ lớn suất điện động cảm ứng bằng độ lớn hệ số góc của đồ thị liên kết từ thông theo thời gian. Khi vận dụng phải xác định đúng dữ kiện, phương chiều nếu có, điều kiện áp dụng, hệ đơn vị và kiểm tra tính hợp lí. Nhận định “Đồ thị từ thông nằm ngang ứng với suất điện động cảm ứng khác không.” sai.
}
\end{bt}

% ===== Câu 11 =====
% ID: L12C3B16-03-TLN
% Mức: TH
\begin{ex}
% ID: L12C3B16-03-TLN
% Mức: TH
Trong bốn phát biểu về Khai thác đồ thị từ thông – thời gian và bài toán cảm ứng điện từ tổng hợp, có bao nhiêu phát biểu đúng? (Chỉ ghi một số nguyên.) (1) Độ lớn suất điện động cảm ứng bằng độ lớn hệ số góc của đồ thị liên kết từ thông theo thời gian. (2) Đồ thị từ thông nằm ngang ứng với suất điện động cảm ứng khác không. (3) Cần kiểm tra điều kiện áp dụng. (4) Có thể bỏ qua đơn vị.
\shortans{2}
\loigiai{
Đối chiếu với kiến thức: Độ lớn suất điện động cảm ứng bằng độ lớn hệ số góc của đồ thị liên kết từ thông theo thời gian. Có 2 phát biểu đúng.
}
\end{ex}

% ===== Câu 12 =====
% ID: L12C3B16-03-TN
% Mức: NB
\begin{ex}
% ID: L12C3B16-03-TN
% Mức: NB
Tình huống 3: chọn kết luận chính xác về Khai thác đồ thị từ thông – thời gian và bài toán cảm ứng điện từ tổng hợp?
\choice
{Kết quả không phụ thuộc dữ kiện và điều kiện áp dụng.}
{\True Độ lớn suất điện động cảm ứng bằng độ lớn hệ số góc của đồ thị liên kết từ thông theo thời gian.}
{Đồ thị từ thông nằm ngang ứng với suất điện động cảm ứng khác không.}
{Có thể bỏ qua phương, chiều và đơn vị trong mọi trường hợp.}
\loigiai{
Độ lớn suất điện động cảm ứng bằng độ lớn hệ số góc của đồ thị liên kết từ thông theo thời gian. Vì vậy phương án $B$ đúng; các phương án còn lại sai.
}
\end{ex}

% ===== Câu 13 =====
% ID: L12C3B16-04-DS
% Mức: VD
\begin{ex}
% ID: L12C3B16-04-DS
% Mức: VD
Xét Khai thác đồ thị từ thông – thời gian và bài toán cảm ứng điện từ tổng hợp (Tình huống 4). Hãy xác định tính đúng, sai của bốn phát biểu sau.
\choiceTF
{Kết quả không phụ thuộc dữ kiện và có thể bỏ qua điều kiện.}
{\True Khi giải cần xác định đúng dữ kiện, phương chiều, điều kiện áp dụng và đơn vị.}
{Đồ thị từ thông nằm ngang ứng với suất điện động cảm ứng khác không.}
{\True Độ lớn suất điện động cảm ứng bằng độ lớn hệ số góc của đồ thị liên kết từ thông theo thời gian.}
\loigiai{
\begin{itemchoice}
\item (Sai) Kết quả không phụ thuộc dữ kiện và có thể bỏ qua điều kiện. (Vì): Kết quả không phụ thuộc dữ kiện và có thể bỏ qua điều kiện. b) Đúng: Khi giải cần xác định đúng dữ kiện, phương chiều, điều kiện áp dụng và đơn vị. c) Sai: Đồ thị từ thông nằm ngang ứng với suất điện động cảm ứng khác không. d) Đúng: Độ lớn suất điện động cảm ứng bằng độ lớn hệ số góc của đồ thị liên kết từ thông theo thời gian. Kiến thức cốt lõi: Độ lớn suất điện động cảm ứng bằng độ lớn hệ số góc của đồ thị liên kết từ thông theo thời gian
\item (Đúng) Khi giải cần xác định đúng dữ kiện, phương chiều, điều kiện áp dụng và đơn vị.
\item (Sai) Đồ thị từ thông nằm ngang ứng với suất điện động cảm ứng khác không.
\item (Đúng) Độ lớn suất điện động cảm ứng bằng độ lớn hệ số góc của đồ thị liên kết từ thông theo thời gian.
\end{itemchoice}
}
\end{ex}

% ===== Câu 14 =====
% ID: L12C3B16-04-TL
% Mức: VD
\begin{bt}
% ID: L12C3B16-04-TL
% Mức: VD
Đề xuất các bước giải một bài toán thuộc dạng Khai thác đồ thị từ thông – thời gian và bài toán cảm ứng điện từ tổng hợp.
\loigiai{
Cần nêu được: Độ lớn suất điện động cảm ứng bằng độ lớn hệ số góc của đồ thị liên kết từ thông theo thời gian. Khi vận dụng phải xác định đúng dữ kiện, phương chiều nếu có, điều kiện áp dụng, hệ đơn vị và kiểm tra tính hợp lí. Nhận định “Đồ thị từ thông nằm ngang ứng với suất điện động cảm ứng khác không.” sai.
}
\end{bt}

% ===== Câu 15 =====
% ID: L12C3B16-04-TLN
% Mức: VD
\begin{ex}
% ID: L12C3B16-04-TLN
% Mức: VD
Trong bốn phát biểu về Khai thác đồ thị từ thông – thời gian và bài toán cảm ứng điện từ tổng hợp, có bao nhiêu phát biểu đúng? (Chỉ ghi một số nguyên.) (1) Đồ thị từ thông nằm ngang ứng với suất điện động cảm ứng khác không. (2) Độ lớn suất điện động cảm ứng bằng độ lớn hệ số góc của đồ thị liên kết từ thông theo thời gian. (3) Kết quả phải phù hợp ý nghĩa vật lí. (4) Mọi đại lượng đều là vô hướng.
\shortans{2}
\loigiai{
Đối chiếu với kiến thức: Độ lớn suất điện động cảm ứng bằng độ lớn hệ số góc của đồ thị liên kết từ thông theo thời gian. Có 2 phát biểu đúng.
}
\end{ex}

% ===== Câu 16 =====
% ID: L12C3B16-04-TN
% Mức: TH
\begin{ex}
% ID: L12C3B16-04-TN
% Mức: TH
Nội dung nào mô tả đúng nhất Khai thác đồ thị từ thông – thời gian và bài toán cảm ứng điện từ tổng hợp?
\choice
{Có thể bỏ qua phương, chiều và đơn vị trong mọi trường hợp.}
{Kết quả không phụ thuộc dữ kiện và điều kiện áp dụng.}
{\True Độ lớn suất điện động cảm ứng bằng độ lớn hệ số góc của đồ thị liên kết từ thông theo thời gian.}
{Đồ thị từ thông nằm ngang ứng với suất điện động cảm ứng khác không.}
\loigiai{
Độ lớn suất điện động cảm ứng bằng độ lớn hệ số góc của đồ thị liên kết từ thông theo thời gian. Vì vậy phương án $C$ đúng; các phương án còn lại sai.
}
\end{ex}

% ===== Câu 17 =====
% ID: L12C3B16-05-DS
% Mức: VD
\begin{ex}
% ID: L12C3B16-05-DS
% Mức: VD
Xét Khai thác đồ thị từ thông – thời gian và bài toán cảm ứng điện từ tổng hợp (Tình huống 5). Hãy xác định tính đúng, sai của bốn phát biểu sau.
\choiceTF
{\True Độ lớn suất điện động cảm ứng bằng độ lớn hệ số góc của đồ thị liên kết từ thông theo thời gian.}
{Kết quả không phụ thuộc dữ kiện và có thể bỏ qua điều kiện.}
{Đồ thị từ thông nằm ngang ứng với suất điện động cảm ứng khác không.}
{\True Khi giải cần xác định đúng dữ kiện, phương chiều, điều kiện áp dụng và đơn vị.}
\loigiai{
\begin{itemchoice}
\item (Đúng) Độ lớn suất điện động cảm ứng bằng độ lớn hệ số góc của đồ thị liên kết từ thông theo thời gian. (Vì): ộ lớn suất điện động cảm ứng bằng độ lớn hệ số góc của đồ thị liên kết từ thông theo thời gian. b) Sai: Kết quả không phụ thuộc dữ kiện và có thể bỏ qua điều kiện. c) Sai: Đồ thị từ thông nằm ngang ứng với suất điện động cảm ứng khác không. d) Đúng: Khi giải cần xác định đúng dữ kiện, phương chiều, điều kiện áp dụng và đơn vị. Kiến thức cốt lõi: Độ lớn suất điện động cảm ứng bằng độ lớn hệ số góc của đồ thị liên kết từ thông theo thời gian
\item (Sai) Kết quả không phụ thuộc dữ kiện và có thể bỏ qua điều kiện.
\item (Sai) Đồ thị từ thông nằm ngang ứng với suất điện động cảm ứng khác không.
\item (Đúng) Khi giải cần xác định đúng dữ kiện, phương chiều, điều kiện áp dụng và đơn vị.
\end{itemchoice}
}
\end{ex}

% ===== Câu 18 =====
% ID: L12C3B16-05-TL
% Mức: VD
\begin{bt}
% ID: L12C3B16-05-TL
% Mức: VD
Nêu một tình huống thực tiễn liên quan đến Khai thác đồ thị từ thông – thời gian và bài toán cảm ứng điện từ tổng hợp và giải thích bằng kiến thức Vật lí.
\loigiai{
Cần nêu được: Độ lớn suất điện động cảm ứng bằng độ lớn hệ số góc của đồ thị liên kết từ thông theo thời gian. Khi vận dụng phải xác định đúng dữ kiện, phương chiều nếu có, điều kiện áp dụng, hệ đơn vị và kiểm tra tính hợp lí. Nhận định “Đồ thị từ thông nằm ngang ứng với suất điện động cảm ứng khác không.” sai.
}
\end{bt}

% ===== Câu 19 =====
% ID: L12C3B16-05-TLN
% Mức: VD
\begin{ex}
% ID: L12C3B16-05-TLN
% Mức: VD
Trong bốn phát biểu về Khai thác đồ thị từ thông – thời gian và bài toán cảm ứng điện từ tổng hợp, có bao nhiêu phát biểu đúng? (Chỉ ghi một số nguyên.) (1) Độ lớn suất điện động cảm ứng bằng độ lớn hệ số góc của đồ thị liên kết từ thông theo thời gian. (2) Đồ thị từ thông nằm ngang ứng với suất điện động cảm ứng khác không. (3) Cần kiểm tra điều kiện áp dụng. (4) Có thể bỏ qua đơn vị.
\shortans{2}
\loigiai{
Đối chiếu với kiến thức: Độ lớn suất điện động cảm ứng bằng độ lớn hệ số góc của đồ thị liên kết từ thông theo thời gian. Có 2 phát biểu đúng.
}
\end{ex}

% ===== Câu 20 =====
% ID: L12C3B16-05-TN
% Mức: TH
\begin{ex}
% ID: L12C3B16-05-TN
% Mức: TH
Khi phân tích Khai thác đồ thị từ thông – thời gian và bài toán cảm ứng điện từ tổng hợp?
\choice
{Đồ thị từ thông nằm ngang ứng với suất điện động cảm ứng khác không.}
{Có thể bỏ qua phương, chiều và đơn vị trong mọi trường hợp.}
{Kết quả không phụ thuộc dữ kiện và điều kiện áp dụng.}
{\True Độ lớn suất điện động cảm ứng bằng độ lớn hệ số góc của đồ thị liên kết từ thông theo thời gian.}
\loigiai{
Độ lớn suất điện động cảm ứng bằng độ lớn hệ số góc của đồ thị liên kết từ thông theo thời gian. Vì vậy phương án $D$ đúng; các phương án còn lại sai.
}
\end{ex}

% ===== Câu 21 =====
% ID: L12C3B16-06-DS
% Mức: TH
\dangbt{Nhận biết hiện tượng cảm ứng điện từ và điều kiện xuất hiện dòng điện cảm ứng}
\begin{ex}
% ID: L12C3B16-06-DS
% Mức: TH
Xét Nhận biết hiện tượng cảm ứng điện từ và điều kiện xuất hiện dòng điện cảm ứng (Tình huống 1). Hãy xác định tính đúng, sai của bốn phát biểu sau.
\choiceTF
{\True Dòng điện cảm ứng xuất hiện trong mạch kín khi từ thông qua mạch biến thiên.}
{Chỉ cần đặt mạch kín đứng yên trong từ trường đều không đổi là luôn có dòng điện cảm ứng.}
{\True Khi giải cần xác định đúng dữ kiện, phương chiều, điều kiện áp dụng và đơn vị.}
{Kết quả không phụ thuộc dữ kiện và có thể bỏ qua điều kiện.}
\loigiai{
\begin{itemchoice}
\item (Đúng) Dòng điện cảm ứng xuất hiện trong mạch kín khi từ thông qua mạch biến thiên. (Vì): òng điện cảm ứng xuất hiện trong mạch kín khi từ thông qua mạch biến thiên. b) Sai: Chỉ cần đặt mạch kín đứng yên trong từ trường đều không đổi là luôn có dòng điện cảm ứng. c) Đúng: Khi giải cần xác định đúng dữ kiện, phương chiều, điều kiện áp dụng và đơn vị. d) Sai: Kết quả không phụ thuộc dữ kiện và có thể bỏ qua điều kiện. Kiến thức cốt lõi: Dòng điện cảm ứng xuất hiện trong mạch kín khi từ thông qua mạch biến thiên
\item (Sai) Chỉ cần đặt mạch kín đứng yên trong từ trường đều không đổi là luôn có dòng điện cảm ứng.
\item (Đúng) Khi giải cần xác định đúng dữ kiện, phương chiều, điều kiện áp dụng và đơn vị.
\item (Sai) Kết quả không phụ thuộc dữ kiện và có thể bỏ qua điều kiện.
\end{itemchoice}
}
\end{ex}

% ===== Câu 22 =====
% ID: L12C3B16-06-TL
% Mức: VDC
\dangbt{Khai thác đồ thị từ thông – thời gian và bài toán cảm ứng điện từ tổng hợp}
\begin{bt}
% ID: L12C3B16-06-TL
% Mức: VDC
So sánh nhận định đúng “Độ lớn suất điện động cảm ứng bằng độ lớn hệ số góc của đồ thị liên kết từ thông theo thời gian.” với nhận định sai “Đồ thị từ thông nằm ngang ứng với suất điện động cảm ứng khác không.”; chỉ rõ căn cứ Vật lí.
\loigiai{
Cần nêu được: Độ lớn suất điện động cảm ứng bằng độ lớn hệ số góc của đồ thị liên kết từ thông theo thời gian. Khi vận dụng phải xác định đúng dữ kiện, phương chiều nếu có, điều kiện áp dụng, hệ đơn vị và kiểm tra tính hợp lí. Nhận định “Đồ thị từ thông nằm ngang ứng với suất điện động cảm ứng khác không.” sai.
}
\end{bt}

% ===== Câu 23 =====
% ID: L12C3B16-06-TLN
% Mức: VDC
\begin{ex}
% ID: L12C3B16-06-TLN
% Mức: VDC
Trong bốn phát biểu về Khai thác đồ thị từ thông – thời gian và bài toán cảm ứng điện từ tổng hợp, có bao nhiêu phát biểu đúng? (Chỉ ghi một số nguyên.) (1) Khi giải cần xác định đúng dữ kiện, phương chiều, điều kiện áp dụng và đơn vị. (2) Đồ thị từ thông nằm ngang ứng với suất điện động cảm ứng khác không. (3) Phải dùng đúng định luật cho tình huống. (4) Độ lớn suất điện động cảm ứng bằng độ lớn hệ số góc của đồ thị liên kết từ thông theo thời gian.
\shortans{3}
\loigiai{
Đối chiếu với kiến thức: Độ lớn suất điện động cảm ứng bằng độ lớn hệ số góc của đồ thị liên kết từ thông theo thời gian. Có 3 phát biểu đúng.
}
\end{ex}

% ===== Câu 24 =====
% ID: L12C3B16-06-TN
% Mức: TH
\begin{ex}
% ID: L12C3B16-06-TN
% Mức: TH
Một học sinh ôn tập Khai thác đồ thị từ thông – thời gian và bài toán cảm ứng điện từ tổng hợp?
\choice
{\True Độ lớn suất điện động cảm ứng bằng độ lớn hệ số góc của đồ thị liên kết từ thông theo thời gian.}
{Đồ thị từ thông nằm ngang ứng với suất điện động cảm ứng khác không.}
{Có thể bỏ qua phương, chiều và đơn vị trong mọi trường hợp.}
{Kết quả không phụ thuộc dữ kiện và điều kiện áp dụng.}
\loigiai{
Độ lớn suất điện động cảm ứng bằng độ lớn hệ số góc của đồ thị liên kết từ thông theo thời gian. Vì vậy phương án $A$ đúng; các phương án còn lại sai.
}
\end{ex}

% ===== Câu 25 =====
% ID: L12C3B16-07-DS
% Mức: TH
\dangbt{Nhận biết hiện tượng cảm ứng điện từ và điều kiện xuất hiện dòng điện cảm ứng}
\begin{ex}
% ID: L12C3B16-07-DS
% Mức: TH
Xét Nhận biết hiện tượng cảm ứng điện từ và điều kiện xuất hiện dòng điện cảm ứng (Tình huống 2). Hãy xác định tính đúng, sai của bốn phát biểu sau.
\choiceTF
{Chỉ cần đặt mạch kín đứng yên trong từ trường đều không đổi là luôn có dòng điện cảm ứng.}
{\True Dòng điện cảm ứng xuất hiện trong mạch kín khi từ thông qua mạch biến thiên.}
{Kết quả không phụ thuộc dữ kiện và có thể bỏ qua điều kiện.}
{\True Khi giải cần xác định đúng dữ kiện, phương chiều, điều kiện áp dụng và đơn vị.}
\loigiai{
\begin{itemchoice}
\item (Sai) Chỉ cần đặt mạch kín đứng yên trong từ trường đều không đổi là luôn có dòng điện cảm ứng. (Vì): Chỉ cần đặt mạch kín đứng yên trong từ trường đều không đổi là luôn có dòng điện cảm ứng. b) Đúng: Dòng điện cảm ứng xuất hiện trong mạch kín khi từ thông qua mạch biến thiên. c) Sai: Kết quả không phụ thuộc dữ kiện và có thể bỏ qua điều kiện. d) Đúng: Khi giải cần xác định đúng dữ kiện, phương chiều, điều kiện áp dụng và đơn vị. Kiến thức cốt lõi: Dòng điện cảm ứng xuất hiện trong mạch kín khi từ thông qua mạch biến thiên
\item (Đúng) Dòng điện cảm ứng xuất hiện trong mạch kín khi từ thông qua mạch biến thiên.
\item (Sai) Kết quả không phụ thuộc dữ kiện và có thể bỏ qua điều kiện.
\item (Đúng) Khi giải cần xác định đúng dữ kiện, phương chiều, điều kiện áp dụng và đơn vị.
\end{itemchoice}
}
\end{ex}

% ===== Câu 26 =====
% ID: L12C3B16-07-TL
% Mức: TH
\begin{bt}
% ID: L12C3B16-07-TL
% Mức: TH
Trình bày kiến thức cốt lõi của Nhận biết hiện tượng cảm ứng điện từ và điều kiện xuất hiện dòng điện cảm ứng và nêu điều kiện áp dụng.
\loigiai{
Cần nêu được: Dòng điện cảm ứng xuất hiện trong mạch kín khi từ thông qua mạch biến thiên. Khi vận dụng phải xác định đúng dữ kiện, phương chiều nếu có, điều kiện áp dụng, hệ đơn vị và kiểm tra tính hợp lí. Nhận định “Chỉ cần đặt mạch kín đứng yên trong từ trường đều không đổi là luôn có dòng điện cảm ứng.” sai.
}
\end{bt}

% ===== Câu 27 =====
% ID: L12C3B16-07-TLN
% Mức: TH
\begin{ex}
% ID: L12C3B16-07-TLN
% Mức: TH
Trong bốn phát biểu về Nhận biết hiện tượng cảm ứng điện từ và điều kiện xuất hiện dòng điện cảm ứng, có bao nhiêu phát biểu đúng? (Chỉ ghi một số nguyên.) (1) Dòng điện cảm ứng xuất hiện trong mạch kín khi từ thông qua mạch biến thiên. (2) Chỉ cần đặt mạch kín đứng yên trong từ trường đều không đổi là luôn có dòng điện cảm ứng. (3) Cần kiểm tra điều kiện áp dụng. (4) Có thể bỏ qua đơn vị.
\shortans{2}
\loigiai{
Đối chiếu với kiến thức: Dòng điện cảm ứng xuất hiện trong mạch kín khi từ thông qua mạch biến thiên. Có 2 phát biểu đúng.
}
\end{ex}

% ===== Câu 28 =====
% ID: L12C3B16-07-TN
% Mức: TH
\dangbt{Khai thác đồ thị từ thông – thời gian và bài toán cảm ứng điện từ tổng hợp}
\begin{ex}
% ID: L12C3B16-07-TN
% Mức: TH
Chọn phát biểu không trái với kiến thức về Khai thác đồ thị từ thông – thời gian và bài toán cảm ứng điện từ tổng hợp?
\choice
{Kết quả không phụ thuộc dữ kiện và điều kiện áp dụng.}
{\True Độ lớn suất điện động cảm ứng bằng độ lớn hệ số góc của đồ thị liên kết từ thông theo thời gian.}
{Đồ thị từ thông nằm ngang ứng với suất điện động cảm ứng khác không.}
{Có thể bỏ qua phương, chiều và đơn vị trong mọi trường hợp.}
\loigiai{
Độ lớn suất điện động cảm ứng bằng độ lớn hệ số góc của đồ thị liên kết từ thông theo thời gian. Vì vậy phương án $B$ đúng; các phương án còn lại sai.
}
\end{ex}

% ===== Câu 29 =====
% ID: L12C3B16-08-DS
% Mức: VD
\dangbt{Nhận biết hiện tượng cảm ứng điện từ và điều kiện xuất hiện dòng điện cảm ứng}
\begin{ex}
% ID: L12C3B16-08-DS
% Mức: VD
Xét Nhận biết hiện tượng cảm ứng điện từ và điều kiện xuất hiện dòng điện cảm ứng (Tình huống 3). Hãy xác định tính đúng, sai của bốn phát biểu sau.
\choiceTF
{\True Khi giải cần xác định đúng dữ kiện, phương chiều, điều kiện áp dụng và đơn vị.}
{Kết quả không phụ thuộc dữ kiện và có thể bỏ qua điều kiện.}
{\True Dòng điện cảm ứng xuất hiện trong mạch kín khi từ thông qua mạch biến thiên.}
{Chỉ cần đặt mạch kín đứng yên trong từ trường đều không đổi là luôn có dòng điện cảm ứng.}
\loigiai{
\begin{itemchoice}
\item (Đúng) Khi giải cần xác định đúng dữ kiện, phương chiều, điều kiện áp dụng và đơn vị. (Vì): Khi giải cần xác định đúng dữ kiện, phương chiều, điều kiện áp dụng và đơn vị. b) Sai: Kết quả không phụ thuộc dữ kiện và có thể bỏ qua điều kiện. c) Đúng: Dòng điện cảm ứng xuất hiện trong mạch kín khi từ thông qua mạch biến thiên. d) Sai: Chỉ cần đặt mạch kín đứng yên trong từ trường đều không đổi là luôn có dòng điện cảm ứng. Kiến thức cốt lõi: Dòng điện cảm ứng xuất hiện trong mạch kín khi từ thông qua mạch biến thiên
\item (Sai) Kết quả không phụ thuộc dữ kiện và có thể bỏ qua điều kiện.
\item (Đúng) Dòng điện cảm ứng xuất hiện trong mạch kín khi từ thông qua mạch biến thiên.
\item (Sai) Chỉ cần đặt mạch kín đứng yên trong từ trường đều không đổi là luôn có dòng điện cảm ứng.
\end{itemchoice}
}
\end{ex}

% ===== Câu 30 =====
% ID: L12C3B16-08-TL
% Mức: TH
\begin{bt}
% ID: L12C3B16-08-TL
% Mức: TH
Giải thích vì sao nhận định sau đúng: “Dòng điện cảm ứng xuất hiện trong mạch kín khi từ thông qua mạch biến thiên.”
\loigiai{
Cần nêu được: Dòng điện cảm ứng xuất hiện trong mạch kín khi từ thông qua mạch biến thiên. Khi vận dụng phải xác định đúng dữ kiện, phương chiều nếu có, điều kiện áp dụng, hệ đơn vị và kiểm tra tính hợp lí. Nhận định “Chỉ cần đặt mạch kín đứng yên trong từ trường đều không đổi là luôn có dòng điện cảm ứng.” sai.
}
\end{bt}

% ===== Câu 31 =====
% ID: L12C3B16-08-TLN
% Mức: TH
\begin{ex}
% ID: L12C3B16-08-TLN
% Mức: TH
Trong bốn phát biểu về Nhận biết hiện tượng cảm ứng điện từ và điều kiện xuất hiện dòng điện cảm ứng, có bao nhiêu phát biểu đúng? (Chỉ ghi một số nguyên.) (1) Chỉ cần đặt mạch kín đứng yên trong từ trường đều không đổi là luôn có dòng điện cảm ứng. (2) Dòng điện cảm ứng xuất hiện trong mạch kín khi từ thông qua mạch biến thiên. (3) Kết quả phải phù hợp ý nghĩa vật lí. (4) Mọi đại lượng đều là vô hướng.
\shortans{1}
\loigiai{
Đối chiếu với kiến thức: Dòng điện cảm ứng xuất hiện trong mạch kín khi từ thông qua mạch biến thiên. Có 1 phát biểu đúng.
}
\end{ex}

% ===== Câu 32 =====
% ID: L12C3B16-08-TN
% Mức: VD
\dangbt{Khai thác đồ thị từ thông – thời gian và bài toán cảm ứng điện từ tổng hợp}
\begin{ex}
% ID: L12C3B16-08-TN
% Mức: VD
Cơ sở Vật lí đúng dùng để giải Khai thác đồ thị từ thông – thời gian và bài toán cảm ứng điện từ tổng hợp?
\choice
{Có thể bỏ qua phương, chiều và đơn vị trong mọi trường hợp.}
{Kết quả không phụ thuộc dữ kiện và điều kiện áp dụng.}
{\True Độ lớn suất điện động cảm ứng bằng độ lớn hệ số góc của đồ thị liên kết từ thông theo thời gian.}
{Đồ thị từ thông nằm ngang ứng với suất điện động cảm ứng khác không.}
\loigiai{
Độ lớn suất điện động cảm ứng bằng độ lớn hệ số góc của đồ thị liên kết từ thông theo thời gian. Vì vậy phương án $C$ đúng; các phương án còn lại sai.
}
\end{ex}

% ===== Câu 33 =====
% ID: L12C3B16-09-DS
% Mức: VD
\dangbt{Nhận biết hiện tượng cảm ứng điện từ và điều kiện xuất hiện dòng điện cảm ứng}
\begin{ex}
% ID: L12C3B16-09-DS
% Mức: VD
Xét Nhận biết hiện tượng cảm ứng điện từ và điều kiện xuất hiện dòng điện cảm ứng (Tình huống 4). Hãy xác định tính đúng, sai của bốn phát biểu sau.
\choiceTF
{Kết quả không phụ thuộc dữ kiện và có thể bỏ qua điều kiện.}
{\True Khi giải cần xác định đúng dữ kiện, phương chiều, điều kiện áp dụng và đơn vị.}
{Chỉ cần đặt mạch kín đứng yên trong từ trường đều không đổi là luôn có dòng điện cảm ứng.}
{\True Dòng điện cảm ứng xuất hiện trong mạch kín khi từ thông qua mạch biến thiên.}
\loigiai{
\begin{itemchoice}
\item (Sai) Kết quả không phụ thuộc dữ kiện và có thể bỏ qua điều kiện. (Vì): Kết quả không phụ thuộc dữ kiện và có thể bỏ qua điều kiện. b) Đúng: Khi giải cần xác định đúng dữ kiện, phương chiều, điều kiện áp dụng và đơn vị. c) Sai: Chỉ cần đặt mạch kín đứng yên trong từ trường đều không đổi là luôn có dòng điện cảm ứng. d) Đúng: Dòng điện cảm ứng xuất hiện trong mạch kín khi từ thông qua mạch biến thiên. Kiến thức cốt lõi: Dòng điện cảm ứng xuất hiện trong mạch kín khi từ thông qua mạch biến thiên
\item (Đúng) Khi giải cần xác định đúng dữ kiện, phương chiều, điều kiện áp dụng và đơn vị.
\item (Sai) Chỉ cần đặt mạch kín đứng yên trong từ trường đều không đổi là luôn có dòng điện cảm ứng.
\item (Đúng) Dòng điện cảm ứng xuất hiện trong mạch kín khi từ thông qua mạch biến thiên.
\end{itemchoice}
}
\end{ex}

% ===== Câu 34 =====
% ID: L12C3B16-09-TL
% Mức: VD
\begin{bt}
% ID: L12C3B16-09-TL
% Mức: VD
Phân tích sai lầm trong nhận định: “Chỉ cần đặt mạch kín đứng yên trong từ trường đều không đổi là luôn có dòng điện cảm ứng.”
\loigiai{
Cần nêu được: Dòng điện cảm ứng xuất hiện trong mạch kín khi từ thông qua mạch biến thiên. Khi vận dụng phải xác định đúng dữ kiện, phương chiều nếu có, điều kiện áp dụng, hệ đơn vị và kiểm tra tính hợp lí. Nhận định “Chỉ cần đặt mạch kín đứng yên trong từ trường đều không đổi là luôn có dòng điện cảm ứng.” sai.
}
\end{bt}

% ===== Câu 35 =====
% ID: L12C3B16-09-TLN
% Mức: TH
\begin{ex}
% ID: L12C3B16-09-TLN
% Mức: TH
Trong bốn phát biểu về Nhận biết hiện tượng cảm ứng điện từ và điều kiện xuất hiện dòng điện cảm ứng, có bao nhiêu phát biểu đúng? (Chỉ ghi một số nguyên.) (1) Dòng điện cảm ứng xuất hiện trong mạch kín khi từ thông qua mạch biến thiên. (2) Chỉ cần đặt mạch kín đứng yên trong từ trường đều không đổi là luôn có dòng điện cảm ứng. (3) Cần kiểm tra điều kiện áp dụng. (4) Có thể bỏ qua đơn vị.
\shortans{2}
\loigiai{
Đối chiếu với kiến thức: Dòng điện cảm ứng xuất hiện trong mạch kín khi từ thông qua mạch biến thiên. Có 2 phát biểu đúng.
}
\end{ex}

% ===== Câu 36 =====
% ID: L12C3B16-09-TN
% Mức: VD
\dangbt{Khai thác đồ thị từ thông – thời gian và bài toán cảm ứng điện từ tổng hợp}
\begin{ex}
% ID: L12C3B16-09-TN
% Mức: VD
Trong một bài thực hành về Khai thác đồ thị từ thông – thời gian và bài toán cảm ứng điện từ tổng hợp?
\choice
{Đồ thị từ thông nằm ngang ứng với suất điện động cảm ứng khác không.}
{Có thể bỏ qua phương, chiều và đơn vị trong mọi trường hợp.}
{Kết quả không phụ thuộc dữ kiện và điều kiện áp dụng.}
{\True Độ lớn suất điện động cảm ứng bằng độ lớn hệ số góc của đồ thị liên kết từ thông theo thời gian.}
\loigiai{
Độ lớn suất điện động cảm ứng bằng độ lớn hệ số góc của đồ thị liên kết từ thông theo thời gian. Vì vậy phương án $D$ đúng; các phương án còn lại sai.
}
\end{ex}

% ===== Câu 37 =====
% ID: L12C3B16-10-DS
% Mức: VD
\dangbt{Nhận biết hiện tượng cảm ứng điện từ và điều kiện xuất hiện dòng điện cảm ứng}
\begin{ex}
% ID: L12C3B16-10-DS
% Mức: VD
Xét Nhận biết hiện tượng cảm ứng điện từ và điều kiện xuất hiện dòng điện cảm ứng (Tình huống 5). Hãy xác định tính đúng, sai của bốn phát biểu sau.
\choiceTF
{\True Dòng điện cảm ứng xuất hiện trong mạch kín khi từ thông qua mạch biến thiên.}
{Kết quả không phụ thuộc dữ kiện và có thể bỏ qua điều kiện.}
{Chỉ cần đặt mạch kín đứng yên trong từ trường đều không đổi là luôn có dòng điện cảm ứng.}
{\True Khi giải cần xác định đúng dữ kiện, phương chiều, điều kiện áp dụng và đơn vị.}
\loigiai{
\begin{itemchoice}
\item (Đúng) Dòng điện cảm ứng xuất hiện trong mạch kín khi từ thông qua mạch biến thiên. (Vì): òng điện cảm ứng xuất hiện trong mạch kín khi từ thông qua mạch biến thiên. b) Sai: Kết quả không phụ thuộc dữ kiện và có thể bỏ qua điều kiện. c) Sai: Chỉ cần đặt mạch kín đứng yên trong từ trường đều không đổi là luôn có dòng điện cảm ứng. d) Đúng: Khi giải cần xác định đúng dữ kiện, phương chiều, điều kiện áp dụng và đơn vị. Kiến thức cốt lõi: Dòng điện cảm ứng xuất hiện trong mạch kín khi từ thông qua mạch biến thiên
\item (Sai) Kết quả không phụ thuộc dữ kiện và có thể bỏ qua điều kiện.
\item (Sai) Chỉ cần đặt mạch kín đứng yên trong từ trường đều không đổi là luôn có dòng điện cảm ứng.
\item (Đúng) Khi giải cần xác định đúng dữ kiện, phương chiều, điều kiện áp dụng và đơn vị.
\end{itemchoice}
}
\end{ex}

% ===== Câu 38 =====
% ID: L12C3B16-10-TL
% Mức: VD
\begin{bt}
% ID: L12C3B16-10-TL
% Mức: VD
Đề xuất các bước giải một bài toán thuộc dạng Nhận biết hiện tượng cảm ứng điện từ và điều kiện xuất hiện dòng điện cảm ứng.
\loigiai{
Cần nêu được: Dòng điện cảm ứng xuất hiện trong mạch kín khi từ thông qua mạch biến thiên. Khi vận dụng phải xác định đúng dữ kiện, phương chiều nếu có, điều kiện áp dụng, hệ đơn vị và kiểm tra tính hợp lí. Nhận định “Chỉ cần đặt mạch kín đứng yên trong từ trường đều không đổi là luôn có dòng điện cảm ứng.” sai.
}
\end{bt}

% ===== Câu 39 =====
% ID: L12C3B16-10-TLN
% Mức: VD
\begin{ex}
% ID: L12C3B16-10-TLN
% Mức: VD
Trong bốn phát biểu về Nhận biết hiện tượng cảm ứng điện từ và điều kiện xuất hiện dòng điện cảm ứng, có bao nhiêu phát biểu đúng? (Chỉ ghi một số nguyên.) (1) Chỉ cần đặt mạch kín đứng yên trong từ trường đều không đổi là luôn có dòng điện cảm ứng. (2) Dòng điện cảm ứng xuất hiện trong mạch kín khi từ thông qua mạch biến thiên. (3) Kết quả phải phù hợp ý nghĩa vật lí. (4) Mọi đại lượng đều là vô hướng.
\shortans{2}
\loigiai{
Đối chiếu với kiến thức: Dòng điện cảm ứng xuất hiện trong mạch kín khi từ thông qua mạch biến thiên. Có 2 phát biểu đúng.
}
\end{ex}

% ===== Câu 40 =====
% ID: L12C3B16-10-TN
% Mức: VD
\dangbt{Khai thác đồ thị từ thông – thời gian và bài toán cảm ứng điện từ tổng hợp}
\begin{ex}
% ID: L12C3B16-10-TN
% Mức: VD
Kết luận đúng khi vận dụng Khai thác đồ thị từ thông – thời gian và bài toán cảm ứng điện từ tổng hợp?
\choice
{\True Độ lớn suất điện động cảm ứng bằng độ lớn hệ số góc của đồ thị liên kết từ thông theo thời gian.}
{Đồ thị từ thông nằm ngang ứng với suất điện động cảm ứng khác không.}
{Có thể bỏ qua phương, chiều và đơn vị trong mọi trường hợp.}
{Kết quả không phụ thuộc dữ kiện và điều kiện áp dụng.}
\loigiai{
Độ lớn suất điện động cảm ứng bằng độ lớn hệ số góc của đồ thị liên kết từ thông theo thời gian. Vì vậy phương án $A$ đúng; các phương án còn lại sai.
}
\end{ex}

% ===== Câu 41 =====
% ID: L12C3B16-100-TN
% Mức: VD
\dangbt{Xác định chiều dòng điện cảm ứng bằng định luật Lenz}
\begin{ex}
% ID: L12C3B16-100-TN
% Mức: VD
Kết luận đúng khi vận dụng Xác định chiều dòng điện cảm ứng bằng định luật Lenz?
\choice
{\True Dòng điện cảm ứng có chiều sao cho từ trường do nó sinh ra chống lại sự biến thiên từ thông đã gây ra nó.}
{Dòng điện cảm ứng luôn làm tăng biến thiên từ thông ban đầu.}
{Có thể bỏ qua phương, chiều và đơn vị trong mọi trường hợp.}
{Kết quả không phụ thuộc dữ kiện và điều kiện áp dụng.}
\loigiai{
Dòng điện cảm ứng có chiều sao cho từ trường do nó sinh ra chống lại sự biến thiên từ thông đã gây ra nó. Vì vậy phương án $A$ đúng; các phương án còn lại sai.
}
\end{ex}

% ===== Câu 42 =====
% ID: L12C3B16-101-TN
% Mức: NB
\dangbt{Tính suất điện động cảm ứng theo định luật Faraday}
\begin{ex}
% ID: L12C3B16-101-TN
% Mức: NB
Phát biểu nào sau đây đúng về Tính suất điện động cảm ứng theo định luật Faraday?
\choice
{Suất điện động cảm ứng bằng $N$ Delta t/Delta Phi.}
{Có thể bỏ qua phương, chiều và đơn vị trong mọi trường hợp.}
{Kết quả không phụ thuộc dữ kiện và điều kiện áp dụng.}
{\True Độ lớn suất điện động cảm ứng trung bình của cuộn $N$ vòng là |e| = N|Delta Phi|/Delta t.}
\loigiai{
Độ lớn suất điện động cảm ứng trung bình của cuộn $N$ vòng là |e| = N|Delta Phi|/Delta t. Vì vậy phương án $D$ đúng; các phương án còn lại sai.
}
\end{ex}

% ===== Câu 43 =====
% ID: L12C3B16-102-TN
% Mức: NB
\begin{ex}
% ID: L12C3B16-102-TN
% Mức: NB
Trong nội dung Tính suất điện động cảm ứng theo định luật Faraday?
\choice
{\True Độ lớn suất điện động cảm ứng trung bình của cuộn $N$ vòng là |e| = N|Delta Phi|/Delta t.}
{Suất điện động cảm ứng bằng $N$ Delta t/Delta Phi.}
{Có thể bỏ qua phương, chiều và đơn vị trong mọi trường hợp.}
{Kết quả không phụ thuộc dữ kiện và điều kiện áp dụng.}
\loigiai{
Độ lớn suất điện động cảm ứng trung bình của cuộn $N$ vòng là |e| = N|Delta Phi|/Delta t. Vì vậy phương án $A$ đúng; các phương án còn lại sai.
}
\end{ex}

% ===== Câu 44 =====
% ID: L12C3B16-103-TN
% Mức: NB
\begin{ex}
% ID: L12C3B16-103-TN
% Mức: NB
Tình huống 3: chọn kết luận chính xác về Tính suất điện động cảm ứng theo định luật Faraday?
\choice
{Kết quả không phụ thuộc dữ kiện và điều kiện áp dụng.}
{\True Độ lớn suất điện động cảm ứng trung bình của cuộn $N$ vòng là |e| = N|Delta Phi|/Delta t.}
{Suất điện động cảm ứng bằng $N$ Delta t/Delta Phi.}
{Có thể bỏ qua phương, chiều và đơn vị trong mọi trường hợp.}
\loigiai{
Độ lớn suất điện động cảm ứng trung bình của cuộn $N$ vòng là |e| = N|Delta Phi|/Delta t. Vì vậy phương án $B$ đúng; các phương án còn lại sai.
}
\end{ex}

% ===== Câu 45 =====
% ID: L12C3B16-104-TN
% Mức: TH
\begin{ex}
% ID: L12C3B16-104-TN
% Mức: TH
Nội dung nào mô tả đúng nhất Tính suất điện động cảm ứng theo định luật Faraday?
\choice
{Có thể bỏ qua phương, chiều và đơn vị trong mọi trường hợp.}
{Kết quả không phụ thuộc dữ kiện và điều kiện áp dụng.}
{\True Độ lớn suất điện động cảm ứng trung bình của cuộn $N$ vòng là |e| = N|Delta Phi|/Delta t.}
{Suất điện động cảm ứng bằng $N$ Delta t/Delta Phi.}
\loigiai{
Độ lớn suất điện động cảm ứng trung bình của cuộn $N$ vòng là |e| = N|Delta Phi|/Delta t. Vì vậy phương án $C$ đúng; các phương án còn lại sai.
}
\end{ex}

% ===== Câu 46 =====
% ID: L12C3B16-105-TN
% Mức: TH
\begin{ex}
% ID: L12C3B16-105-TN
% Mức: TH
Khi phân tích Tính suất điện động cảm ứng theo định luật Faraday?
\choice
{Suất điện động cảm ứng bằng $N$ Delta t/Delta Phi.}
{Có thể bỏ qua phương, chiều và đơn vị trong mọi trường hợp.}
{Kết quả không phụ thuộc dữ kiện và điều kiện áp dụng.}
{\True Độ lớn suất điện động cảm ứng trung bình của cuộn $N$ vòng là |e| = N|Delta Phi|/Delta t.}
\loigiai{
Độ lớn suất điện động cảm ứng trung bình của cuộn $N$ vòng là |e| = N|Delta Phi|/Delta t. Vì vậy phương án $D$ đúng; các phương án còn lại sai.
}
\end{ex}

% ===== Câu 47 =====
% ID: L12C3B16-106-TN
% Mức: TH
\begin{ex}
% ID: L12C3B16-106-TN
% Mức: TH
Một học sinh ôn tập Tính suất điện động cảm ứng theo định luật Faraday?
\choice
{\True Độ lớn suất điện động cảm ứng trung bình của cuộn $N$ vòng là |e| = N|Delta Phi|/Delta t.}
{Suất điện động cảm ứng bằng $N$ Delta t/Delta Phi.}
{Có thể bỏ qua phương, chiều và đơn vị trong mọi trường hợp.}
{Kết quả không phụ thuộc dữ kiện và điều kiện áp dụng.}
\loigiai{
Độ lớn suất điện động cảm ứng trung bình của cuộn $N$ vòng là |e| = N|Delta Phi|/Delta t. Vì vậy phương án $A$ đúng; các phương án còn lại sai.
}
\end{ex}

% ===== Câu 48 =====
% ID: L12C3B16-107-TN
% Mức: TH
\begin{ex}
% ID: L12C3B16-107-TN
% Mức: TH
Chọn phát biểu không trái với kiến thức về Tính suất điện động cảm ứng theo định luật Faraday?
\choice
{Kết quả không phụ thuộc dữ kiện và điều kiện áp dụng.}
{\True Độ lớn suất điện động cảm ứng trung bình của cuộn $N$ vòng là |e| = N|Delta Phi|/Delta t.}
{Suất điện động cảm ứng bằng $N$ Delta t/Delta Phi.}
{Có thể bỏ qua phương, chiều và đơn vị trong mọi trường hợp.}
\loigiai{
Độ lớn suất điện động cảm ứng trung bình của cuộn $N$ vòng là |e| = N|Delta Phi|/Delta t. Vì vậy phương án $B$ đúng; các phương án còn lại sai.
}
\end{ex}

% ===== Câu 49 =====
% ID: L12C3B16-108-TN
% Mức: VD
\begin{ex}
% ID: L12C3B16-108-TN
% Mức: VD
Cơ sở Vật lí đúng dùng để giải Tính suất điện động cảm ứng theo định luật Faraday?
\choice
{Có thể bỏ qua phương, chiều và đơn vị trong mọi trường hợp.}
{Kết quả không phụ thuộc dữ kiện và điều kiện áp dụng.}
{\True Độ lớn suất điện động cảm ứng trung bình của cuộn $N$ vòng là |e| = N|Delta Phi|/Delta t.}
{Suất điện động cảm ứng bằng $N$ Delta t/Delta Phi.}
\loigiai{
Độ lớn suất điện động cảm ứng trung bình của cuộn $N$ vòng là |e| = N|Delta Phi|/Delta t. Vì vậy phương án $C$ đúng; các phương án còn lại sai.
}
\end{ex}

% ===== Câu 50 =====
% ID: L12C3B16-109-TN
% Mức: VD
\begin{ex}
% ID: L12C3B16-109-TN
% Mức: VD
Trong một bài thực hành về Tính suất điện động cảm ứng theo định luật Faraday?
\choice
{Suất điện động cảm ứng bằng $N$ Delta t/Delta Phi.}
{Có thể bỏ qua phương, chiều và đơn vị trong mọi trường hợp.}
{Kết quả không phụ thuộc dữ kiện và điều kiện áp dụng.}
{\True Độ lớn suất điện động cảm ứng trung bình của cuộn $N$ vòng là |e| = N|Delta Phi|/Delta t.}
\loigiai{
Độ lớn suất điện động cảm ứng trung bình của cuộn $N$ vòng là |e| = N|Delta Phi|/Delta t. Vì vậy phương án $D$ đúng; các phương án còn lại sai.
}
\end{ex}

% ===== Câu 51 =====
% ID: L12C3B16-11-DS
% Mức: TH
\dangbt{Nhận biết từ thông và tính từ thông qua một diện tích}
\begin{ex}
% ID: L12C3B16-11-DS
% Mức: TH
Xét Nhận biết từ thông và tính từ thông qua một diện tích (Tình huống 1). Hãy xác định tính đúng, sai của bốn phát biểu sau.
\choiceTF
{\True Từ thông qua diện tích $S$ trong từ trường đều là Phi = BS cos(alpha), với alpha là góc giữa $B$ và pháp tuyến.}
{Từ thông luôn bằng BS sin(alpha), với alpha là góc giữa $B$ và pháp tuyến.}
{\True Khi giải cần xác định đúng dữ kiện, phương chiều, điều kiện áp dụng và đơn vị.}
{Kết quả không phụ thuộc dữ kiện và có thể bỏ qua điều kiện.}
\loigiai{
\begin{itemchoice}
\item (Đúng) Từ thông qua diện tích $S$ trong từ trường đều là Phi = BS cos(alpha), với alpha là góc giữa $B$ và pháp tuyến. (Vì): lpha), với alpha là góc giữa $B$ và pháp tuyến
\item (Sai) Từ thông luôn bằng BS sin(alpha), với alpha là góc giữa $B$ và pháp tuyến.
\item (Đúng) Khi giải cần xác định đúng dữ kiện, phương chiều, điều kiện áp dụng và đơn vị.
\item (Sai) Kết quả không phụ thuộc dữ kiện và có thể bỏ qua điều kiện.
\end{itemchoice}
}
\end{ex}

% ===== Câu 52 =====
% ID: L12C3B16-11-TL
% Mức: VD
\dangbt{Nhận biết hiện tượng cảm ứng điện từ và điều kiện xuất hiện dòng điện cảm ứng}
\begin{bt}
% ID: L12C3B16-11-TL
% Mức: VD
Nêu một tình huống thực tiễn liên quan đến Nhận biết hiện tượng cảm ứng điện từ và điều kiện xuất hiện dòng điện cảm ứng và giải thích bằng kiến thức Vật lí.
\loigiai{
Cần nêu được: Dòng điện cảm ứng xuất hiện trong mạch kín khi từ thông qua mạch biến thiên. Khi vận dụng phải xác định đúng dữ kiện, phương chiều nếu có, điều kiện áp dụng, hệ đơn vị và kiểm tra tính hợp lí. Nhận định “Chỉ cần đặt mạch kín đứng yên trong từ trường đều không đổi là luôn có dòng điện cảm ứng.” sai.
}
\end{bt}

% ===== Câu 53 =====
% ID: L12C3B16-11-TLN
% Mức: VD
\begin{ex}
% ID: L12C3B16-11-TLN
% Mức: VD
Trong bốn phát biểu về Nhận biết hiện tượng cảm ứng điện từ và điều kiện xuất hiện dòng điện cảm ứng, có bao nhiêu phát biểu đúng? (Chỉ ghi một số nguyên.) (1) Dòng điện cảm ứng xuất hiện trong mạch kín khi từ thông qua mạch biến thiên. (2) Chỉ cần đặt mạch kín đứng yên trong từ trường đều không đổi là luôn có dòng điện cảm ứng. (3) Cần kiểm tra điều kiện áp dụng. (4) Có thể bỏ qua đơn vị.
\shortans{2}
\loigiai{
Đối chiếu với kiến thức: Dòng điện cảm ứng xuất hiện trong mạch kín khi từ thông qua mạch biến thiên. Có 2 phát biểu đúng.
}
\end{ex}

% ===== Câu 54 =====
% ID: L12C3B16-11-TN
% Mức: NB
\begin{ex}
% ID: L12C3B16-11-TN
% Mức: NB
Phát biểu nào sau đây đúng về Nhận biết hiện tượng cảm ứng điện từ và điều kiện xuất hiện dòng điện cảm ứng?
\choice
{Chỉ cần đặt mạch kín đứng yên trong từ trường đều không đổi là luôn có dòng điện cảm ứng.}
{Có thể bỏ qua phương, chiều và đơn vị trong mọi trường hợp.}
{Kết quả không phụ thuộc dữ kiện và điều kiện áp dụng.}
{\True Dòng điện cảm ứng xuất hiện trong mạch kín khi từ thông qua mạch biến thiên.}
\loigiai{
Dòng điện cảm ứng xuất hiện trong mạch kín khi từ thông qua mạch biến thiên. Vì vậy phương án $D$ đúng; các phương án còn lại sai.
}
\end{ex}

% ===== Câu 55 =====
% ID: L12C3B16-110-TN
% Mức: VD
\dangbt{Tính suất điện động cảm ứng theo định luật Faraday}
\begin{ex}
% ID: L12C3B16-110-TN
% Mức: VD
Kết luận đúng khi vận dụng Tính suất điện động cảm ứng theo định luật Faraday?
\choice
{\True Độ lớn suất điện động cảm ứng trung bình của cuộn $N$ vòng là |e| = N|Delta Phi|/Delta t.}
{Suất điện động cảm ứng bằng $N$ Delta t/Delta Phi.}
{Có thể bỏ qua phương, chiều và đơn vị trong mọi trường hợp.}
{Kết quả không phụ thuộc dữ kiện và điều kiện áp dụng.}
\loigiai{
Độ lớn suất điện động cảm ứng trung bình của cuộn $N$ vòng là |e| = N|Delta Phi|/Delta t. Vì vậy phương án $A$ đúng; các phương án còn lại sai.
}
\end{ex}

% ===== Câu 56 =====
% ID: L12C3B16-111-TN
% Mức: NB
\dangbt{Khai thác đồ thị từ thông – thời gian và bài toán cảm ứng điện từ tổng hợp}
\begin{ex}
% ID: L12C3B16-111-TN
% Mức: NB
Phát biểu nào sau đây đúng về Khai thác đồ thị từ thông – thời gian và bài toán cảm ứng điện từ tổng hợp?
\choice
{Đồ thị từ thông nằm ngang ứng với suất điện động cảm ứng khác không.}
{Có thể bỏ qua phương, chiều và đơn vị trong mọi trường hợp.}
{Kết quả không phụ thuộc dữ kiện và điều kiện áp dụng.}
{\True Độ lớn suất điện động cảm ứng bằng độ lớn hệ số góc của đồ thị liên kết từ thông theo thời gian.}
\loigiai{
Độ lớn suất điện động cảm ứng bằng độ lớn hệ số góc của đồ thị liên kết từ thông theo thời gian. Vì vậy phương án $D$ đúng; các phương án còn lại sai.
}
\end{ex}

% ===== Câu 57 =====
% ID: L12C3B16-112-TN
% Mức: NB
\begin{ex}
% ID: L12C3B16-112-TN
% Mức: NB
Trong nội dung Khai thác đồ thị từ thông – thời gian và bài toán cảm ứng điện từ tổng hợp?
\choice
{\True Độ lớn suất điện động cảm ứng bằng độ lớn hệ số góc của đồ thị liên kết từ thông theo thời gian.}
{Đồ thị từ thông nằm ngang ứng với suất điện động cảm ứng khác không.}
{Có thể bỏ qua phương, chiều và đơn vị trong mọi trường hợp.}
{Kết quả không phụ thuộc dữ kiện và điều kiện áp dụng.}
\loigiai{
Độ lớn suất điện động cảm ứng bằng độ lớn hệ số góc của đồ thị liên kết từ thông theo thời gian. Vì vậy phương án $A$ đúng; các phương án còn lại sai.
}
\end{ex}

% ===== Câu 58 =====
% ID: L12C3B16-113-TN
% Mức: NB
\begin{ex}
% ID: L12C3B16-113-TN
% Mức: NB
Tình huống 3: chọn kết luận chính xác về Khai thác đồ thị từ thông – thời gian và bài toán cảm ứng điện từ tổng hợp?
\choice
{Kết quả không phụ thuộc dữ kiện và điều kiện áp dụng.}
{\True Độ lớn suất điện động cảm ứng bằng độ lớn hệ số góc của đồ thị liên kết từ thông theo thời gian.}
{Đồ thị từ thông nằm ngang ứng với suất điện động cảm ứng khác không.}
{Có thể bỏ qua phương, chiều và đơn vị trong mọi trường hợp.}
\loigiai{
Độ lớn suất điện động cảm ứng bằng độ lớn hệ số góc của đồ thị liên kết từ thông theo thời gian. Vì vậy phương án $B$ đúng; các phương án còn lại sai.
}
\end{ex}

% ===== Câu 59 =====
% ID: L12C3B16-114-TN
% Mức: TH
\begin{ex}
% ID: L12C3B16-114-TN
% Mức: TH
Nội dung nào mô tả đúng nhất Khai thác đồ thị từ thông – thời gian và bài toán cảm ứng điện từ tổng hợp?
\choice
{Có thể bỏ qua phương, chiều và đơn vị trong mọi trường hợp.}
{Kết quả không phụ thuộc dữ kiện và điều kiện áp dụng.}
{\True Độ lớn suất điện động cảm ứng bằng độ lớn hệ số góc của đồ thị liên kết từ thông theo thời gian.}
{Đồ thị từ thông nằm ngang ứng với suất điện động cảm ứng khác không.}
\loigiai{
Độ lớn suất điện động cảm ứng bằng độ lớn hệ số góc của đồ thị liên kết từ thông theo thời gian. Vì vậy phương án $C$ đúng; các phương án còn lại sai.
}
\end{ex}

% ===== Câu 60 =====
% ID: L12C3B16-115-TN
% Mức: TH
\begin{ex}
% ID: L12C3B16-115-TN
% Mức: TH
Khi phân tích Khai thác đồ thị từ thông – thời gian và bài toán cảm ứng điện từ tổng hợp?
\choice
{Đồ thị từ thông nằm ngang ứng với suất điện động cảm ứng khác không.}
{Có thể bỏ qua phương, chiều và đơn vị trong mọi trường hợp.}
{Kết quả không phụ thuộc dữ kiện và điều kiện áp dụng.}
{\True Độ lớn suất điện động cảm ứng bằng độ lớn hệ số góc của đồ thị liên kết từ thông theo thời gian.}
\loigiai{
Độ lớn suất điện động cảm ứng bằng độ lớn hệ số góc của đồ thị liên kết từ thông theo thời gian. Vì vậy phương án $D$ đúng; các phương án còn lại sai.
}
\end{ex}

% ===== Câu 61 =====
% ID: L12C3B16-116-TN
% Mức: TH
\begin{ex}
% ID: L12C3B16-116-TN
% Mức: TH
Một học sinh ôn tập Khai thác đồ thị từ thông – thời gian và bài toán cảm ứng điện từ tổng hợp?
\choice
{\True Độ lớn suất điện động cảm ứng bằng độ lớn hệ số góc của đồ thị liên kết từ thông theo thời gian.}
{Đồ thị từ thông nằm ngang ứng với suất điện động cảm ứng khác không.}
{Có thể bỏ qua phương, chiều và đơn vị trong mọi trường hợp.}
{Kết quả không phụ thuộc dữ kiện và điều kiện áp dụng.}
\loigiai{
Độ lớn suất điện động cảm ứng bằng độ lớn hệ số góc của đồ thị liên kết từ thông theo thời gian. Vì vậy phương án $A$ đúng; các phương án còn lại sai.
}
\end{ex}

% ===== Câu 62 =====
% ID: L12C3B16-117-TN
% Mức: TH
\begin{ex}
% ID: L12C3B16-117-TN
% Mức: TH
Chọn phát biểu không trái với kiến thức về Khai thác đồ thị từ thông – thời gian và bài toán cảm ứng điện từ tổng hợp?
\choice
{Kết quả không phụ thuộc dữ kiện và điều kiện áp dụng.}
{\True Độ lớn suất điện động cảm ứng bằng độ lớn hệ số góc của đồ thị liên kết từ thông theo thời gian.}
{Đồ thị từ thông nằm ngang ứng với suất điện động cảm ứng khác không.}
{Có thể bỏ qua phương, chiều và đơn vị trong mọi trường hợp.}
\loigiai{
Độ lớn suất điện động cảm ứng bằng độ lớn hệ số góc của đồ thị liên kết từ thông theo thời gian. Vì vậy phương án $B$ đúng; các phương án còn lại sai.
}
\end{ex}

% ===== Câu 63 =====
% ID: L12C3B16-118-TN
% Mức: VD
\begin{ex}
% ID: L12C3B16-118-TN
% Mức: VD
Cơ sở Vật lí đúng dùng để giải Khai thác đồ thị từ thông – thời gian và bài toán cảm ứng điện từ tổng hợp?
\choice
{Có thể bỏ qua phương, chiều và đơn vị trong mọi trường hợp.}
{Kết quả không phụ thuộc dữ kiện và điều kiện áp dụng.}
{\True Độ lớn suất điện động cảm ứng bằng độ lớn hệ số góc của đồ thị liên kết từ thông theo thời gian.}
{Đồ thị từ thông nằm ngang ứng với suất điện động cảm ứng khác không.}
\loigiai{
Độ lớn suất điện động cảm ứng bằng độ lớn hệ số góc của đồ thị liên kết từ thông theo thời gian. Vì vậy phương án $C$ đúng; các phương án còn lại sai.
}
\end{ex}

% ===== Câu 64 =====
% ID: L12C3B16-119-TN
% Mức: VD
\begin{ex}
% ID: L12C3B16-119-TN
% Mức: VD
Trong một bài thực hành về Khai thác đồ thị từ thông – thời gian và bài toán cảm ứng điện từ tổng hợp?
\choice
{Đồ thị từ thông nằm ngang ứng với suất điện động cảm ứng khác không.}
{Có thể bỏ qua phương, chiều và đơn vị trong mọi trường hợp.}
{Kết quả không phụ thuộc dữ kiện và điều kiện áp dụng.}
{\True Độ lớn suất điện động cảm ứng bằng độ lớn hệ số góc của đồ thị liên kết từ thông theo thời gian.}
\loigiai{
Độ lớn suất điện động cảm ứng bằng độ lớn hệ số góc của đồ thị liên kết từ thông theo thời gian. Vì vậy phương án $D$ đúng; các phương án còn lại sai.
}
\end{ex}

% ===== Câu 65 =====
% ID: L12C3B16-12-DS
% Mức: TH
\dangbt{Nhận biết từ thông và tính từ thông qua một diện tích}
\begin{ex}
% ID: L12C3B16-12-DS
% Mức: TH
Xét Nhận biết từ thông và tính từ thông qua một diện tích (Tình huống 2). Hãy xác định tính đúng, sai của bốn phát biểu sau.
\choiceTF
{Từ thông luôn bằng BS sin(alpha), với alpha là góc giữa $B$ và pháp tuyến.}
{\True Từ thông qua diện tích $S$ trong từ trường đều là Phi = BS cos(alpha), với alpha là góc giữa $B$ và pháp tuyến.}
{Kết quả không phụ thuộc dữ kiện và có thể bỏ qua điều kiện.}
{\True Khi giải cần xác định đúng dữ kiện, phương chiều, điều kiện áp dụng và đơn vị.}
\loigiai{
\begin{itemchoice}
\item (Sai) Từ thông luôn bằng BS sin(alpha), với alpha là góc giữa $B$ và pháp tuyến. (Vì): lpha), với alpha là góc giữa $B$ và pháp tuyến
\item (Đúng) Từ thông qua diện tích $S$ trong từ trường đều là Phi = BS cos(alpha), với alpha là góc giữa $B$ và pháp tuyến.
\item (Sai) Kết quả không phụ thuộc dữ kiện và có thể bỏ qua điều kiện.
\item (Đúng) Khi giải cần xác định đúng dữ kiện, phương chiều, điều kiện áp dụng và đơn vị.
\end{itemchoice}
}
\end{ex}

% ===== Câu 66 =====
% ID: L12C3B16-12-TL
% Mức: VDC
\dangbt{Nhận biết hiện tượng cảm ứng điện từ và điều kiện xuất hiện dòng điện cảm ứng}
\begin{bt}
% ID: L12C3B16-12-TL
% Mức: VDC
So sánh nhận định đúng “Dòng điện cảm ứng xuất hiện trong mạch kín khi từ thông qua mạch biến thiên.” với nhận định sai “Chỉ cần đặt mạch kín đứng yên trong từ trường đều không đổi là luôn có dòng điện cảm ứng.”; chỉ rõ căn cứ Vật lí.
\loigiai{
Cần nêu được: Dòng điện cảm ứng xuất hiện trong mạch kín khi từ thông qua mạch biến thiên. Khi vận dụng phải xác định đúng dữ kiện, phương chiều nếu có, điều kiện áp dụng, hệ đơn vị và kiểm tra tính hợp lí. Nhận định “Chỉ cần đặt mạch kín đứng yên trong từ trường đều không đổi là luôn có dòng điện cảm ứng.” sai.
}
\end{bt}

% ===== Câu 67 =====
% ID: L12C3B16-12-TLN
% Mức: VDC
\begin{ex}
% ID: L12C3B16-12-TLN
% Mức: VDC
Trong bốn phát biểu về Nhận biết hiện tượng cảm ứng điện từ và điều kiện xuất hiện dòng điện cảm ứng, có bao nhiêu phát biểu đúng? (Chỉ ghi một số nguyên.) (1) Khi giải cần xác định đúng dữ kiện, phương chiều, điều kiện áp dụng và đơn vị. (2) Chỉ cần đặt mạch kín đứng yên trong từ trường đều không đổi là luôn có dòng điện cảm ứng. (3) Phải dùng đúng định luật cho tình huống. (4) Dòng điện cảm ứng xuất hiện trong mạch kín khi từ thông qua mạch biến thiên.
\shortans{3}
\loigiai{
Đối chiếu với kiến thức: Dòng điện cảm ứng xuất hiện trong mạch kín khi từ thông qua mạch biến thiên. Có 3 phát biểu đúng.
}
\end{ex}

% ===== Câu 68 =====
% ID: L12C3B16-12-TN
% Mức: NB
\begin{ex}
% ID: L12C3B16-12-TN
% Mức: NB
Trong nội dung Nhận biết hiện tượng cảm ứng điện từ và điều kiện xuất hiện dòng điện cảm ứng?
\choice
{\True Dòng điện cảm ứng xuất hiện trong mạch kín khi từ thông qua mạch biến thiên.}
{Chỉ cần đặt mạch kín đứng yên trong từ trường đều không đổi là luôn có dòng điện cảm ứng.}
{Có thể bỏ qua phương, chiều và đơn vị trong mọi trường hợp.}
{Kết quả không phụ thuộc dữ kiện và điều kiện áp dụng.}
\loigiai{
Dòng điện cảm ứng xuất hiện trong mạch kín khi từ thông qua mạch biến thiên. Vì vậy phương án $A$ đúng; các phương án còn lại sai.
}
\end{ex}

% ===== Câu 69 =====
% ID: L12C3B16-120-TN
% Mức: VD
\dangbt{Khai thác đồ thị từ thông – thời gian và bài toán cảm ứng điện từ tổng hợp}
\begin{ex}
% ID: L12C3B16-120-TN
% Mức: VD
Kết luận đúng khi vận dụng Khai thác đồ thị từ thông – thời gian và bài toán cảm ứng điện từ tổng hợp?
\choice
{\True Độ lớn suất điện động cảm ứng bằng độ lớn hệ số góc của đồ thị liên kết từ thông theo thời gian.}
{Đồ thị từ thông nằm ngang ứng với suất điện động cảm ứng khác không.}
{Có thể bỏ qua phương, chiều và đơn vị trong mọi trường hợp.}
{Kết quả không phụ thuộc dữ kiện và điều kiện áp dụng.}
\loigiai{
Độ lớn suất điện động cảm ứng bằng độ lớn hệ số góc của đồ thị liên kết từ thông theo thời gian. Vì vậy phương án $A$ đúng; các phương án còn lại sai.
}
\end{ex}

% ===== Câu 70 =====
% ID: L12C3B16-13-DS
% Mức: VD
\dangbt{Nhận biết từ thông và tính từ thông qua một diện tích}
\begin{ex}
% ID: L12C3B16-13-DS
% Mức: VD
Xét Nhận biết từ thông và tính từ thông qua một diện tích (Tình huống 3). Hãy xác định tính đúng, sai của bốn phát biểu sau.
\choiceTF
{\True Khi giải cần xác định đúng dữ kiện, phương chiều, điều kiện áp dụng và đơn vị.}
{Kết quả không phụ thuộc dữ kiện và có thể bỏ qua điều kiện.}
{\True Từ thông qua diện tích $S$ trong từ trường đều là Phi = BS cos(alpha), với alpha là góc giữa $B$ và pháp tuyến.}
{Từ thông luôn bằng BS sin(alpha), với alpha là góc giữa $B$ và pháp tuyến.}
\loigiai{
\begin{itemchoice}
\item (Đúng) Khi giải cần xác định đúng dữ kiện, phương chiều, điều kiện áp dụng và đơn vị. (Vì): lpha), với alpha là góc giữa $B$ và pháp tuyến
\item (Sai) Kết quả không phụ thuộc dữ kiện và có thể bỏ qua điều kiện.
\item (Đúng) Từ thông qua diện tích $S$ trong từ trường đều là Phi = BS cos(alpha), với alpha là góc giữa $B$ và pháp tuyến.
\item (Sai) Từ thông luôn bằng BS sin(alpha), với alpha là góc giữa $B$ và pháp tuyến.
\end{itemchoice}
}
\end{ex}

% ===== Câu 71 =====
% ID: L12C3B16-13-TL
% Mức: TH
\begin{bt}
% ID: L12C3B16-13-TL
% Mức: TH
Trình bày kiến thức cốt lõi của Nhận biết từ thông và tính từ thông qua một diện tích và nêu điều kiện áp dụng.
\loigiai{
Cần nêu được: Từ thông qua diện tích $S$ trong từ trường đều là Phi = BS cos(alpha), với alpha là góc giữa $B$ và pháp tuyến. Khi vận dụng phải xác định đúng dữ kiện, phương chiều nếu có, điều kiện áp dụng, hệ đơn vị và kiểm tra tính hợp lí. Nhận định “Từ thông luôn bằng BS sin(alpha), với alpha là góc giữa $B$ và pháp tuyến.” sai.
}
\end{bt}

% ===== Câu 72 =====
% ID: L12C3B16-13-TLN
% Mức: TH
\begin{ex}
% ID: L12C3B16-13-TLN
% Mức: TH
Trong bốn phát biểu về Nhận biết từ thông và tính từ thông qua một diện tích, có bao nhiêu phát biểu đúng? (Chỉ ghi một số nguyên.) (1) Từ thông qua diện tích $S$ trong từ trường đều là Phi = BS cos(alpha), với alpha là góc giữa $B$ và pháp tuyến. (2) Từ thông luôn bằng BS sin(alpha), với alpha là góc giữa $B$ và pháp tuyến. (3) Cần kiểm tra điều kiện áp dụng. (4) Có thể bỏ qua đơn vị.
\shortans{2}
\loigiai{
Đối chiếu với kiến thức: Từ thông qua diện tích $S$ trong từ trường đều là Phi = BS cos(alpha), với alpha là góc giữa $B$ và pháp tuyến. Có 2 phát biểu đúng.
}
\end{ex}

% ===== Câu 73 =====
% ID: L12C3B16-13-TN
% Mức: NB
\dangbt{Nhận biết hiện tượng cảm ứng điện từ và điều kiện xuất hiện dòng điện cảm ứng}
\begin{ex}
% ID: L12C3B16-13-TN
% Mức: NB
Tình huống 3: chọn kết luận chính xác về Nhận biết hiện tượng cảm ứng điện từ và điều kiện xuất hiện dòng điện cảm ứng?
\choice
{Kết quả không phụ thuộc dữ kiện và điều kiện áp dụng.}
{\True Dòng điện cảm ứng xuất hiện trong mạch kín khi từ thông qua mạch biến thiên.}
{Chỉ cần đặt mạch kín đứng yên trong từ trường đều không đổi là luôn có dòng điện cảm ứng.}
{Có thể bỏ qua phương, chiều và đơn vị trong mọi trường hợp.}
\loigiai{
Dòng điện cảm ứng xuất hiện trong mạch kín khi từ thông qua mạch biến thiên. Vì vậy phương án $B$ đúng; các phương án còn lại sai.
}
\end{ex}

% ===== Câu 74 =====
% ID: L12C3B16-14-DS
% Mức: VD
\dangbt{Nhận biết từ thông và tính từ thông qua một diện tích}
\begin{ex}
% ID: L12C3B16-14-DS
% Mức: VD
Xét Nhận biết từ thông và tính từ thông qua một diện tích (Tình huống 4). Hãy xác định tính đúng, sai của bốn phát biểu sau.
\choiceTF
{Kết quả không phụ thuộc dữ kiện và có thể bỏ qua điều kiện.}
{\True Khi giải cần xác định đúng dữ kiện, phương chiều, điều kiện áp dụng và đơn vị.}
{Từ thông luôn bằng BS sin(alpha), với alpha là góc giữa $B$ và pháp tuyến.}
{\True Từ thông qua diện tích $S$ trong từ trường đều là Phi = BS cos(alpha), với alpha là góc giữa $B$ và pháp tuyến.}
\loigiai{
\begin{itemchoice}
\item (Sai) Kết quả không phụ thuộc dữ kiện và có thể bỏ qua điều kiện. (Vì): lpha), với alpha là góc giữa $B$ và pháp tuyến
\item (Đúng) Khi giải cần xác định đúng dữ kiện, phương chiều, điều kiện áp dụng và đơn vị.
\item (Sai) Từ thông luôn bằng BS sin(alpha), với alpha là góc giữa $B$ và pháp tuyến.
\item (Đúng) Từ thông qua diện tích $S$ trong từ trường đều là Phi = BS cos(alpha), với alpha là góc giữa $B$ và pháp tuyến.
\end{itemchoice}
}
\end{ex}

% ===== Câu 75 =====
% ID: L12C3B16-14-TL
% Mức: TH
\begin{bt}
% ID: L12C3B16-14-TL
% Mức: TH
Giải thích vì sao nhận định sau đúng: “Từ thông qua diện tích $S$ trong từ trường đều là Phi = BS cos(alpha), với alpha là góc giữa $B$ và pháp tuyến.”
\loigiai{
Cần nêu được: Từ thông qua diện tích $S$ trong từ trường đều là Phi = BS cos(alpha), với alpha là góc giữa $B$ và pháp tuyến. Khi vận dụng phải xác định đúng dữ kiện, phương chiều nếu có, điều kiện áp dụng, hệ đơn vị và kiểm tra tính hợp lí. Nhận định “Từ thông luôn bằng BS sin(alpha), với alpha là góc giữa $B$ và pháp tuyến.” sai.
}
\end{bt}

% ===== Câu 76 =====
% ID: L12C3B16-14-TLN
% Mức: TH
\begin{ex}
% ID: L12C3B16-14-TLN
% Mức: TH
Trong bốn phát biểu về Nhận biết từ thông và tính từ thông qua một diện tích, có bao nhiêu phát biểu đúng? (Chỉ ghi một số nguyên.) (1) Từ thông luôn bằng BS sin(alpha), với alpha là góc giữa $B$ và pháp tuyến. (2) Từ thông qua diện tích $S$ trong từ trường đều là Phi = BS cos(alpha), với alpha là góc giữa $B$ và pháp tuyến. (3) Kết quả phải phù hợp ý nghĩa vật lí. (4) Mọi đại lượng đều là vô hướng.
\shortans{1}
\loigiai{
Đối chiếu với kiến thức: Từ thông qua diện tích $S$ trong từ trường đều là Phi = BS cos(alpha), với alpha là góc giữa $B$ và pháp tuyến. Có 1 phát biểu đúng.
}
\end{ex}

% ===== Câu 77 =====
% ID: L12C3B16-14-TN
% Mức: TH
\dangbt{Nhận biết hiện tượng cảm ứng điện từ và điều kiện xuất hiện dòng điện cảm ứng}
\begin{ex}
% ID: L12C3B16-14-TN
% Mức: TH
Nội dung nào mô tả đúng nhất Nhận biết hiện tượng cảm ứng điện từ và điều kiện xuất hiện dòng điện cảm ứng?
\choice
{Có thể bỏ qua phương, chiều và đơn vị trong mọi trường hợp.}
{Kết quả không phụ thuộc dữ kiện và điều kiện áp dụng.}
{\True Dòng điện cảm ứng xuất hiện trong mạch kín khi từ thông qua mạch biến thiên.}
{Chỉ cần đặt mạch kín đứng yên trong từ trường đều không đổi là luôn có dòng điện cảm ứng.}
\loigiai{
Dòng điện cảm ứng xuất hiện trong mạch kín khi từ thông qua mạch biến thiên. Vì vậy phương án $C$ đúng; các phương án còn lại sai.
}
\end{ex}

% ===== Câu 78 =====
% ID: L12C3B16-15-DS
% Mức: VD
\dangbt{Nhận biết từ thông và tính từ thông qua một diện tích}
\begin{ex}
% ID: L12C3B16-15-DS
% Mức: VD
Xét Nhận biết từ thông và tính từ thông qua một diện tích (Tình huống 5). Hãy xác định tính đúng, sai của bốn phát biểu sau.
\choiceTF
{\True Từ thông qua diện tích $S$ trong từ trường đều là Phi = BS cos(alpha), với alpha là góc giữa $B$ và pháp tuyến.}
{Kết quả không phụ thuộc dữ kiện và có thể bỏ qua điều kiện.}
{Từ thông luôn bằng BS sin(alpha), với alpha là góc giữa $B$ và pháp tuyến.}
{\True Khi giải cần xác định đúng dữ kiện, phương chiều, điều kiện áp dụng và đơn vị.}
\loigiai{
\begin{itemchoice}
\item (Đúng) Từ thông qua diện tích $S$ trong từ trường đều là Phi = BS cos(alpha), với alpha là góc giữa $B$ và pháp tuyến. (Vì): lpha), với alpha là góc giữa $B$ và pháp tuyến
\item (Sai) Kết quả không phụ thuộc dữ kiện và có thể bỏ qua điều kiện.
\item (Sai) Từ thông luôn bằng BS sin(alpha), với alpha là góc giữa $B$ và pháp tuyến.
\item (Đúng) Khi giải cần xác định đúng dữ kiện, phương chiều, điều kiện áp dụng và đơn vị.
\end{itemchoice}
}
\end{ex}

% ===== Câu 79 =====
% ID: L12C3B16-15-TL
% Mức: VD
\begin{bt}
% ID: L12C3B16-15-TL
% Mức: VD
Phân tích sai lầm trong nhận định: “Từ thông luôn bằng BS sin(alpha), với alpha là góc giữa $B$ và pháp tuyến.”
\loigiai{
Cần nêu được: Từ thông qua diện tích $S$ trong từ trường đều là Phi = BS cos(alpha), với alpha là góc giữa $B$ và pháp tuyến. Khi vận dụng phải xác định đúng dữ kiện, phương chiều nếu có, điều kiện áp dụng, hệ đơn vị và kiểm tra tính hợp lí. Nhận định “Từ thông luôn bằng BS sin(alpha), với alpha là góc giữa $B$ và pháp tuyến.” sai.
}
\end{bt}

% ===== Câu 80 =====
% ID: L12C3B16-15-TLN
% Mức: TH
\begin{ex}
% ID: L12C3B16-15-TLN
% Mức: TH
Trong bốn phát biểu về Nhận biết từ thông và tính từ thông qua một diện tích, có bao nhiêu phát biểu đúng? (Chỉ ghi một số nguyên.) (1) Từ thông qua diện tích $S$ trong từ trường đều là Phi = BS cos(alpha), với alpha là góc giữa $B$ và pháp tuyến. (2) Từ thông luôn bằng BS sin(alpha), với alpha là góc giữa $B$ và pháp tuyến. (3) Cần kiểm tra điều kiện áp dụng. (4) Có thể bỏ qua đơn vị.
\shortans{2}
\loigiai{
Đối chiếu với kiến thức: Từ thông qua diện tích $S$ trong từ trường đều là Phi = BS cos(alpha), với alpha là góc giữa $B$ và pháp tuyến. Có 2 phát biểu đúng.
}
\end{ex}

% ===== Câu 81 =====
% ID: L12C3B16-15-TN
% Mức: TH
\dangbt{Nhận biết hiện tượng cảm ứng điện từ và điều kiện xuất hiện dòng điện cảm ứng}
\begin{ex}
% ID: L12C3B16-15-TN
% Mức: TH
Khi phân tích Nhận biết hiện tượng cảm ứng điện từ và điều kiện xuất hiện dòng điện cảm ứng?
\choice
{Chỉ cần đặt mạch kín đứng yên trong từ trường đều không đổi là luôn có dòng điện cảm ứng.}
{Có thể bỏ qua phương, chiều và đơn vị trong mọi trường hợp.}
{Kết quả không phụ thuộc dữ kiện và điều kiện áp dụng.}
{\True Dòng điện cảm ứng xuất hiện trong mạch kín khi từ thông qua mạch biến thiên.}
\loigiai{
Dòng điện cảm ứng xuất hiện trong mạch kín khi từ thông qua mạch biến thiên. Vì vậy phương án $D$ đúng; các phương án còn lại sai.
}
\end{ex}

% ===== Câu 82 =====
% ID: L12C3B16-16-DS
% Mức: TH
\dangbt{Phân tích sự biến thiên từ thông qua khung dây}
\begin{ex}
% ID: L12C3B16-16-DS
% Mức: TH
Xét Phân tích sự biến thiên từ thông qua khung dây (Tình huống 1). Hãy xác định tính đúng, sai của bốn phát biểu sau.
\choiceTF
{\True Từ thông biến thiên khi $B$, diện tích mạch hoặc góc giữa $B$ và pháp tuyến thay đổi.}
{Từ thông không đổi trong mọi chuyển động của khung dây.}
{\True Khi giải cần xác định đúng dữ kiện, phương chiều, điều kiện áp dụng và đơn vị.}
{Kết quả không phụ thuộc dữ kiện và có thể bỏ qua điều kiện.}
\loigiai{
\begin{itemchoice}
\item (Đúng) Từ thông biến thiên khi $B$, diện tích mạch hoặc góc giữa $B$ và pháp tuyến thay đổi. (Vì): Từ thông biến thiên khi $B$, diện tích mạch hoặc góc giữa $B$ và pháp tuyến thay đổi. b) Sai: Từ thông không đổi trong mọi chuyển động của khung dây. c) Đúng: Khi giải cần xác định đúng dữ kiện, phương chiều, điều kiện áp dụng và đơn vị. d) Sai: Kết quả không phụ thuộc dữ kiện và có thể bỏ qua điều kiện. Kiến thức cốt lõi: Từ thông biến thiên khi $B$, diện tích mạch hoặc góc giữa $B$ và pháp tuyến thay đổi
\item (Sai) Từ thông không đổi trong mọi chuyển động của khung dây.
\item (Đúng) Khi giải cần xác định đúng dữ kiện, phương chiều, điều kiện áp dụng và đơn vị.
\item (Sai) Kết quả không phụ thuộc dữ kiện và có thể bỏ qua điều kiện.
\end{itemchoice}
}
\end{ex}

% ===== Câu 83 =====
% ID: L12C3B16-16-TL
% Mức: VD
\dangbt{Nhận biết từ thông và tính từ thông qua một diện tích}
\begin{bt}
% ID: L12C3B16-16-TL
% Mức: VD
Đề xuất các bước giải một bài toán thuộc dạng Nhận biết từ thông và tính từ thông qua một diện tích.
\loigiai{
Cần nêu được: Từ thông qua diện tích $S$ trong từ trường đều là Phi = BS cos(alpha), với alpha là góc giữa $B$ và pháp tuyến. Khi vận dụng phải xác định đúng dữ kiện, phương chiều nếu có, điều kiện áp dụng, hệ đơn vị và kiểm tra tính hợp lí. Nhận định “Từ thông luôn bằng BS sin(alpha), với alpha là góc giữa $B$ và pháp tuyến.” sai.
}
\end{bt}

% ===== Câu 84 =====
% ID: L12C3B16-16-TLN
% Mức: VD
\begin{ex}
% ID: L12C3B16-16-TLN
% Mức: VD
Trong bốn phát biểu về Nhận biết từ thông và tính từ thông qua một diện tích, có bao nhiêu phát biểu đúng? (Chỉ ghi một số nguyên.) (1) Từ thông luôn bằng BS sin(alpha), với alpha là góc giữa $B$ và pháp tuyến. (2) Từ thông qua diện tích $S$ trong từ trường đều là Phi = BS cos(alpha), với alpha là góc giữa $B$ và pháp tuyến. (3) Kết quả phải phù hợp ý nghĩa vật lí. (4) Mọi đại lượng đều là vô hướng.
\shortans{2}
\loigiai{
Đối chiếu với kiến thức: Từ thông qua diện tích $S$ trong từ trường đều là Phi = BS cos(alpha), với alpha là góc giữa $B$ và pháp tuyến. Có 2 phát biểu đúng.
}
\end{ex}

% ===== Câu 85 =====
% ID: L12C3B16-16-TN
% Mức: TH
\dangbt{Nhận biết hiện tượng cảm ứng điện từ và điều kiện xuất hiện dòng điện cảm ứng}
\begin{ex}
% ID: L12C3B16-16-TN
% Mức: TH
Một học sinh ôn tập Nhận biết hiện tượng cảm ứng điện từ và điều kiện xuất hiện dòng điện cảm ứng?
\choice
{\True Dòng điện cảm ứng xuất hiện trong mạch kín khi từ thông qua mạch biến thiên.}
{Chỉ cần đặt mạch kín đứng yên trong từ trường đều không đổi là luôn có dòng điện cảm ứng.}
{Có thể bỏ qua phương, chiều và đơn vị trong mọi trường hợp.}
{Kết quả không phụ thuộc dữ kiện và điều kiện áp dụng.}
\loigiai{
Dòng điện cảm ứng xuất hiện trong mạch kín khi từ thông qua mạch biến thiên. Vì vậy phương án $A$ đúng; các phương án còn lại sai.
}
\end{ex}

% ===== Câu 86 =====
% ID: L12C3B16-17-DS
% Mức: TH
\dangbt{Phân tích sự biến thiên từ thông qua khung dây}
\begin{ex}
% ID: L12C3B16-17-DS
% Mức: TH
Xét Phân tích sự biến thiên từ thông qua khung dây (Tình huống 2). Hãy xác định tính đúng, sai của bốn phát biểu sau.
\choiceTF
{Từ thông không đổi trong mọi chuyển động của khung dây.}
{\True Từ thông biến thiên khi $B$, diện tích mạch hoặc góc giữa $B$ và pháp tuyến thay đổi.}
{Kết quả không phụ thuộc dữ kiện và có thể bỏ qua điều kiện.}
{\True Khi giải cần xác định đúng dữ kiện, phương chiều, điều kiện áp dụng và đơn vị.}
\loigiai{
\begin{itemchoice}
\item (Sai) Từ thông không đổi trong mọi chuyển động của khung dây. (Vì): Từ thông không đổi trong mọi chuyển động của khung dây. b) Đúng: Từ thông biến thiên khi $B$, diện tích mạch hoặc góc giữa $B$ và pháp tuyến thay đổi. c) Sai: Kết quả không phụ thuộc dữ kiện và có thể bỏ qua điều kiện. d) Đúng: Khi giải cần xác định đúng dữ kiện, phương chiều, điều kiện áp dụng và đơn vị. Kiến thức cốt lõi: Từ thông biến thiên khi $B$, diện tích mạch hoặc góc giữa $B$ và pháp tuyến thay đổi
\item (Đúng) Từ thông biến thiên khi $B$, diện tích mạch hoặc góc giữa $B$ và pháp tuyến thay đổi.
\item (Sai) Kết quả không phụ thuộc dữ kiện và có thể bỏ qua điều kiện.
\item (Đúng) Khi giải cần xác định đúng dữ kiện, phương chiều, điều kiện áp dụng và đơn vị.
\end{itemchoice}
}
\end{ex}

% ===== Câu 87 =====
% ID: L12C3B16-17-TL
% Mức: VD
\dangbt{Nhận biết từ thông và tính từ thông qua một diện tích}
\begin{bt}
% ID: L12C3B16-17-TL
% Mức: VD
Nêu một tình huống thực tiễn liên quan đến Nhận biết từ thông và tính từ thông qua một diện tích và giải thích bằng kiến thức Vật lí.
\loigiai{
Cần nêu được: Từ thông qua diện tích $S$ trong từ trường đều là Phi = BS cos(alpha), với alpha là góc giữa $B$ và pháp tuyến. Khi vận dụng phải xác định đúng dữ kiện, phương chiều nếu có, điều kiện áp dụng, hệ đơn vị và kiểm tra tính hợp lí. Nhận định “Từ thông luôn bằng BS sin(alpha), với alpha là góc giữa $B$ và pháp tuyến.” sai.
}
\end{bt}

% ===== Câu 88 =====
% ID: L12C3B16-17-TLN
% Mức: VD
\begin{ex}
% ID: L12C3B16-17-TLN
% Mức: VD
Trong bốn phát biểu về Nhận biết từ thông và tính từ thông qua một diện tích, có bao nhiêu phát biểu đúng? (Chỉ ghi một số nguyên.) (1) Từ thông qua diện tích $S$ trong từ trường đều là Phi = BS cos(alpha), với alpha là góc giữa $B$ và pháp tuyến. (2) Từ thông luôn bằng BS sin(alpha), với alpha là góc giữa $B$ và pháp tuyến. (3) Cần kiểm tra điều kiện áp dụng. (4) Có thể bỏ qua đơn vị.
\shortans{2}
\loigiai{
Đối chiếu với kiến thức: Từ thông qua diện tích $S$ trong từ trường đều là Phi = BS cos(alpha), với alpha là góc giữa $B$ và pháp tuyến. Có 2 phát biểu đúng.
}
\end{ex}

% ===== Câu 89 =====
% ID: L12C3B16-17-TN
% Mức: TH
\dangbt{Nhận biết hiện tượng cảm ứng điện từ và điều kiện xuất hiện dòng điện cảm ứng}
\begin{ex}
% ID: L12C3B16-17-TN
% Mức: TH
Chọn phát biểu không trái với kiến thức về Nhận biết hiện tượng cảm ứng điện từ và điều kiện xuất hiện dòng điện cảm ứng?
\choice
{Kết quả không phụ thuộc dữ kiện và điều kiện áp dụng.}
{\True Dòng điện cảm ứng xuất hiện trong mạch kín khi từ thông qua mạch biến thiên.}
{Chỉ cần đặt mạch kín đứng yên trong từ trường đều không đổi là luôn có dòng điện cảm ứng.}
{Có thể bỏ qua phương, chiều và đơn vị trong mọi trường hợp.}
\loigiai{
Dòng điện cảm ứng xuất hiện trong mạch kín khi từ thông qua mạch biến thiên. Vì vậy phương án $B$ đúng; các phương án còn lại sai.
}
\end{ex}

% ===== Câu 90 =====
% ID: L12C3B16-18-DS
% Mức: VD
\dangbt{Phân tích sự biến thiên từ thông qua khung dây}
\begin{ex}
% ID: L12C3B16-18-DS
% Mức: VD
Xét Phân tích sự biến thiên từ thông qua khung dây (Tình huống 3). Hãy xác định tính đúng, sai của bốn phát biểu sau.
\choiceTF
{\True Khi giải cần xác định đúng dữ kiện, phương chiều, điều kiện áp dụng và đơn vị.}
{Kết quả không phụ thuộc dữ kiện và có thể bỏ qua điều kiện.}
{\True Từ thông biến thiên khi $B$, diện tích mạch hoặc góc giữa $B$ và pháp tuyến thay đổi.}
{Từ thông không đổi trong mọi chuyển động của khung dây.}
\loigiai{
\begin{itemchoice}
\item (Đúng) Khi giải cần xác định đúng dữ kiện, phương chiều, điều kiện áp dụng và đơn vị. (Vì): Khi giải cần xác định đúng dữ kiện, phương chiều, điều kiện áp dụng và đơn vị. b) Sai: Kết quả không phụ thuộc dữ kiện và có thể bỏ qua điều kiện. c) Đúng: Từ thông biến thiên khi $B$, diện tích mạch hoặc góc giữa $B$ và pháp tuyến thay đổi. d) Sai: Từ thông không đổi trong mọi chuyển động của khung dây. Kiến thức cốt lõi: Từ thông biến thiên khi $B$, diện tích mạch hoặc góc giữa $B$ và pháp tuyến thay đổi
\item (Sai) Kết quả không phụ thuộc dữ kiện và có thể bỏ qua điều kiện.
\item (Đúng) Từ thông biến thiên khi $B$, diện tích mạch hoặc góc giữa $B$ và pháp tuyến thay đổi.
\item (Sai) Từ thông không đổi trong mọi chuyển động của khung dây.
\end{itemchoice}
}
\end{ex}

% ===== Câu 91 =====
% ID: L12C3B16-18-TL
% Mức: VDC
\dangbt{Nhận biết từ thông và tính từ thông qua một diện tích}
\begin{bt}
% ID: L12C3B16-18-TL
% Mức: VDC
So sánh nhận định đúng “Từ thông qua diện tích $S$ trong từ trường đều là Phi = BS cos(alpha), với alpha là góc giữa $B$ và pháp tuyến.” với nhận định sai “Từ thông luôn bằng BS sin(alpha), với alpha là góc giữa $B$ và pháp tuyến.”; chỉ rõ căn cứ Vật lí.
\loigiai{
Cần nêu được: Từ thông qua diện tích $S$ trong từ trường đều là Phi = BS cos(alpha), với alpha là góc giữa $B$ và pháp tuyến. Khi vận dụng phải xác định đúng dữ kiện, phương chiều nếu có, điều kiện áp dụng, hệ đơn vị và kiểm tra tính hợp lí. Nhận định “Từ thông luôn bằng BS sin(alpha), với alpha là góc giữa $B$ và pháp tuyến.” sai.
}
\end{bt}

% ===== Câu 92 =====
% ID: L12C3B16-18-TLN
% Mức: VDC
\begin{ex}
% ID: L12C3B16-18-TLN
% Mức: VDC
Trong bốn phát biểu về Nhận biết từ thông và tính từ thông qua một diện tích, có bao nhiêu phát biểu đúng? (Chỉ ghi một số nguyên.) (1) Khi giải cần xác định đúng dữ kiện, phương chiều, điều kiện áp dụng và đơn vị. (2) Từ thông luôn bằng BS sin(alpha), với alpha là góc giữa $B$ và pháp tuyến. (3) Phải dùng đúng định luật cho tình huống. (4) Từ thông qua diện tích $S$ trong từ trường đều là Phi = BS cos(alpha), với alpha là góc giữa $B$ và pháp tuyến.
\shortans{3}
\loigiai{
Đối chiếu với kiến thức: Từ thông qua diện tích $S$ trong từ trường đều là Phi = BS cos(alpha), với alpha là góc giữa $B$ và pháp tuyến. Có 3 phát biểu đúng.
}
\end{ex}

% ===== Câu 93 =====
% ID: L12C3B16-18-TN
% Mức: VD
\dangbt{Nhận biết hiện tượng cảm ứng điện từ và điều kiện xuất hiện dòng điện cảm ứng}
\begin{ex}
% ID: L12C3B16-18-TN
% Mức: VD
Cơ sở Vật lí đúng dùng để giải Nhận biết hiện tượng cảm ứng điện từ và điều kiện xuất hiện dòng điện cảm ứng?
\choice
{Có thể bỏ qua phương, chiều và đơn vị trong mọi trường hợp.}
{Kết quả không phụ thuộc dữ kiện và điều kiện áp dụng.}
{\True Dòng điện cảm ứng xuất hiện trong mạch kín khi từ thông qua mạch biến thiên.}
{Chỉ cần đặt mạch kín đứng yên trong từ trường đều không đổi là luôn có dòng điện cảm ứng.}
\loigiai{
Dòng điện cảm ứng xuất hiện trong mạch kín khi từ thông qua mạch biến thiên. Vì vậy phương án $C$ đúng; các phương án còn lại sai.
}
\end{ex}

% ===== Câu 94 =====
% ID: L12C3B16-19-DS
% Mức: VD
\dangbt{Phân tích sự biến thiên từ thông qua khung dây}
\begin{ex}
% ID: L12C3B16-19-DS
% Mức: VD
Xét Phân tích sự biến thiên từ thông qua khung dây (Tình huống 4). Hãy xác định tính đúng, sai của bốn phát biểu sau.
\choiceTF
{Kết quả không phụ thuộc dữ kiện và có thể bỏ qua điều kiện.}
{\True Khi giải cần xác định đúng dữ kiện, phương chiều, điều kiện áp dụng và đơn vị.}
{Từ thông không đổi trong mọi chuyển động của khung dây.}
{\True Từ thông biến thiên khi $B$, diện tích mạch hoặc góc giữa $B$ và pháp tuyến thay đổi.}
\loigiai{
\begin{itemchoice}
\item (Sai) Kết quả không phụ thuộc dữ kiện và có thể bỏ qua điều kiện. (Vì): Kết quả không phụ thuộc dữ kiện và có thể bỏ qua điều kiện. b) Đúng: Khi giải cần xác định đúng dữ kiện, phương chiều, điều kiện áp dụng và đơn vị. c) Sai: Từ thông không đổi trong mọi chuyển động của khung dây. d) Đúng: Từ thông biến thiên khi $B$, diện tích mạch hoặc góc giữa $B$ và pháp tuyến thay đổi. Kiến thức cốt lõi: Từ thông biến thiên khi $B$, diện tích mạch hoặc góc giữa $B$ và pháp tuyến thay đổi
\item (Đúng) Khi giải cần xác định đúng dữ kiện, phương chiều, điều kiện áp dụng và đơn vị.
\item (Sai) Từ thông không đổi trong mọi chuyển động của khung dây.
\item (Đúng) Từ thông biến thiên khi $B$, diện tích mạch hoặc góc giữa $B$ và pháp tuyến thay đổi.
\end{itemchoice}
}
\end{ex}

% ===== Câu 95 =====
% ID: L12C3B16-19-TL
% Mức: TH
\begin{bt}
% ID: L12C3B16-19-TL
% Mức: TH
Trình bày kiến thức cốt lõi của Phân tích sự biến thiên từ thông qua khung dây và nêu điều kiện áp dụng.
\loigiai{
Cần nêu được: Từ thông biến thiên khi $B$, diện tích mạch hoặc góc giữa $B$ và pháp tuyến thay đổi. Khi vận dụng phải xác định đúng dữ kiện, phương chiều nếu có, điều kiện áp dụng, hệ đơn vị và kiểm tra tính hợp lí. Nhận định “Từ thông không đổi trong mọi chuyển động của khung dây.” sai.
}
\end{bt}

% ===== Câu 96 =====
% ID: L12C3B16-19-TLN
% Mức: TH
\begin{ex}
% ID: L12C3B16-19-TLN
% Mức: TH
Trong bốn phát biểu về Phân tích sự biến thiên từ thông qua khung dây, có bao nhiêu phát biểu đúng? (Chỉ ghi một số nguyên.) (1) Từ thông biến thiên khi $B$, diện tích mạch hoặc góc giữa $B$ và pháp tuyến thay đổi. (2) Từ thông không đổi trong mọi chuyển động của khung dây. (3) Cần kiểm tra điều kiện áp dụng. (4) Có thể bỏ qua đơn vị.
\shortans{2}
\loigiai{
Đối chiếu với kiến thức: Từ thông biến thiên khi $B$, diện tích mạch hoặc góc giữa $B$ và pháp tuyến thay đổi. Có 2 phát biểu đúng.
}
\end{ex}

% ===== Câu 97 =====
% ID: L12C3B16-19-TN
% Mức: VD
\dangbt{Nhận biết hiện tượng cảm ứng điện từ và điều kiện xuất hiện dòng điện cảm ứng}
\begin{ex}
% ID: L12C3B16-19-TN
% Mức: VD
Trong một bài thực hành về Nhận biết hiện tượng cảm ứng điện từ và điều kiện xuất hiện dòng điện cảm ứng?
\choice
{Chỉ cần đặt mạch kín đứng yên trong từ trường đều không đổi là luôn có dòng điện cảm ứng.}
{Có thể bỏ qua phương, chiều và đơn vị trong mọi trường hợp.}
{Kết quả không phụ thuộc dữ kiện và điều kiện áp dụng.}
{\True Dòng điện cảm ứng xuất hiện trong mạch kín khi từ thông qua mạch biến thiên.}
\loigiai{
Dòng điện cảm ứng xuất hiện trong mạch kín khi từ thông qua mạch biến thiên. Vì vậy phương án $D$ đúng; các phương án còn lại sai.
}
\end{ex}

% ===== Câu 98 =====
% ID: L12C3B16-20-DS
% Mức: VD
\dangbt{Phân tích sự biến thiên từ thông qua khung dây}
\begin{ex}
% ID: L12C3B16-20-DS
% Mức: VD
Xét Phân tích sự biến thiên từ thông qua khung dây (Tình huống 5). Hãy xác định tính đúng, sai của bốn phát biểu sau.
\choiceTF
{\True Từ thông biến thiên khi $B$, diện tích mạch hoặc góc giữa $B$ và pháp tuyến thay đổi.}
{Kết quả không phụ thuộc dữ kiện và có thể bỏ qua điều kiện.}
{Từ thông không đổi trong mọi chuyển động của khung dây.}
{\True Khi giải cần xác định đúng dữ kiện, phương chiều, điều kiện áp dụng và đơn vị.}
\loigiai{
\begin{itemchoice}
\item (Đúng) Từ thông biến thiên khi $B$, diện tích mạch hoặc góc giữa $B$ và pháp tuyến thay đổi. (Vì): Từ thông biến thiên khi $B$, diện tích mạch hoặc góc giữa $B$ và pháp tuyến thay đổi. b) Sai: Kết quả không phụ thuộc dữ kiện và có thể bỏ qua điều kiện. c) Sai: Từ thông không đổi trong mọi chuyển động của khung dây. d) Đúng: Khi giải cần xác định đúng dữ kiện, phương chiều, điều kiện áp dụng và đơn vị. Kiến thức cốt lõi: Từ thông biến thiên khi $B$, diện tích mạch hoặc góc giữa $B$ và pháp tuyến thay đổi
\item (Sai) Kết quả không phụ thuộc dữ kiện và có thể bỏ qua điều kiện.
\item (Sai) Từ thông không đổi trong mọi chuyển động của khung dây.
\item (Đúng) Khi giải cần xác định đúng dữ kiện, phương chiều, điều kiện áp dụng và đơn vị.
\end{itemchoice}
}
\end{ex}

% ===== Câu 99 =====
% ID: L12C3B16-20-TL
% Mức: TH
\begin{bt}
% ID: L12C3B16-20-TL
% Mức: TH
Giải thích vì sao nhận định sau đúng: “Từ thông biến thiên khi $B$, diện tích mạch hoặc góc giữa $B$ và pháp tuyến thay đổi.”
\loigiai{
Cần nêu được: Từ thông biến thiên khi $B$, diện tích mạch hoặc góc giữa $B$ và pháp tuyến thay đổi. Khi vận dụng phải xác định đúng dữ kiện, phương chiều nếu có, điều kiện áp dụng, hệ đơn vị và kiểm tra tính hợp lí. Nhận định “Từ thông không đổi trong mọi chuyển động của khung dây.” sai.
}
\end{bt}

% ===== Câu 100 =====
% ID: L12C3B16-20-TLN
% Mức: TH
\begin{ex}
% ID: L12C3B16-20-TLN
% Mức: TH
Trong bốn phát biểu về Phân tích sự biến thiên từ thông qua khung dây, có bao nhiêu phát biểu đúng? (Chỉ ghi một số nguyên.) (1) Từ thông không đổi trong mọi chuyển động của khung dây. (2) Từ thông biến thiên khi $B$, diện tích mạch hoặc góc giữa $B$ và pháp tuyến thay đổi. (3) Kết quả phải phù hợp ý nghĩa vật lí. (4) Mọi đại lượng đều là vô hướng.
\shortans{1}
\loigiai{
Đối chiếu với kiến thức: Từ thông biến thiên khi $B$, diện tích mạch hoặc góc giữa $B$ và pháp tuyến thay đổi. Có 1 phát biểu đúng.
}
\end{ex}

% ===== Câu 101 =====
% ID: L12C3B16-20-TN
% Mức: VD
\dangbt{Nhận biết hiện tượng cảm ứng điện từ và điều kiện xuất hiện dòng điện cảm ứng}
\begin{ex}
% ID: L12C3B16-20-TN
% Mức: VD
Kết luận đúng khi vận dụng Nhận biết hiện tượng cảm ứng điện từ và điều kiện xuất hiện dòng điện cảm ứng?
\choice
{\True Dòng điện cảm ứng xuất hiện trong mạch kín khi từ thông qua mạch biến thiên.}
{Chỉ cần đặt mạch kín đứng yên trong từ trường đều không đổi là luôn có dòng điện cảm ứng.}
{Có thể bỏ qua phương, chiều và đơn vị trong mọi trường hợp.}
{Kết quả không phụ thuộc dữ kiện và điều kiện áp dụng.}
\loigiai{
Dòng điện cảm ứng xuất hiện trong mạch kín khi từ thông qua mạch biến thiên. Vì vậy phương án $A$ đúng; các phương án còn lại sai.
}
\end{ex}

% ===== Câu 102 =====
% ID: L12C3B16-21-DS
% Mức: TH
\dangbt{Tính suất điện động cảm ứng theo định luật Faraday}
\begin{ex}
% ID: L12C3B16-21-DS
% Mức: TH
Xét Tính suất điện động cảm ứng theo định luật Faraday (Tình huống 1). Hãy xác định tính đúng, sai của bốn phát biểu sau.
\choiceTF
{\True Độ lớn suất điện động cảm ứng trung bình của cuộn $N$ vòng là |e| = N|Delta Phi|/Delta t.}
{Suất điện động cảm ứng bằng $N$ Delta t/Delta Phi.}
{\True Khi giải cần xác định đúng dữ kiện, phương chiều, điều kiện áp dụng và đơn vị.}
{Kết quả không phụ thuộc dữ kiện và có thể bỏ qua điều kiện.}
\loigiai{
\begin{itemchoice}
\item (Đúng) Độ lớn suất điện động cảm ứng trung bình của cuộn $N$ vòng là |e| = N|Delta Phi|/Delta t. (Vì): ộ lớn suất điện động cảm ứng trung bình của cuộn $N$ vòng là |e| = N|Delta Phi|/Delta t. b) Sai: Suất điện động cảm ứng bằng $N$ Delta t/Delta Phi. c) Đúng: Khi giải cần xác định đúng dữ kiện, phương chiều, điều kiện áp dụng và đơn vị. d) Sai: Kết quả không phụ thuộc dữ kiện và có thể bỏ qua điều kiện. Kiến thức cốt lõi: Độ lớn suất điện động cảm ứng trung bình của cuộn $N$ vòng là |e| = N|Delta Phi|/Delta t
\item (Sai) Suất điện động cảm ứng bằng $N$ Delta t/Delta Phi.
\item (Đúng) Khi giải cần xác định đúng dữ kiện, phương chiều, điều kiện áp dụng và đơn vị.
\item (Sai) Kết quả không phụ thuộc dữ kiện và có thể bỏ qua điều kiện.
\end{itemchoice}
}
\end{ex}

% ===== Câu 103 =====
% ID: L12C3B16-21-TL
% Mức: VD
\dangbt{Phân tích sự biến thiên từ thông qua khung dây}
\begin{bt}
% ID: L12C3B16-21-TL
% Mức: VD
Phân tích sai lầm trong nhận định: “Từ thông không đổi trong mọi chuyển động của khung dây.”
\loigiai{
Cần nêu được: Từ thông biến thiên khi $B$, diện tích mạch hoặc góc giữa $B$ và pháp tuyến thay đổi. Khi vận dụng phải xác định đúng dữ kiện, phương chiều nếu có, điều kiện áp dụng, hệ đơn vị và kiểm tra tính hợp lí. Nhận định “Từ thông không đổi trong mọi chuyển động của khung dây.” sai.
}
\end{bt}

% ===== Câu 104 =====
% ID: L12C3B16-21-TLN
% Mức: TH
\begin{ex}
% ID: L12C3B16-21-TLN
% Mức: TH
Trong bốn phát biểu về Phân tích sự biến thiên từ thông qua khung dây, có bao nhiêu phát biểu đúng? (Chỉ ghi một số nguyên.) (1) Từ thông biến thiên khi $B$, diện tích mạch hoặc góc giữa $B$ và pháp tuyến thay đổi. (2) Từ thông không đổi trong mọi chuyển động của khung dây. (3) Cần kiểm tra điều kiện áp dụng. (4) Có thể bỏ qua đơn vị.
\shortans{2}
\loigiai{
Đối chiếu với kiến thức: Từ thông biến thiên khi $B$, diện tích mạch hoặc góc giữa $B$ và pháp tuyến thay đổi. Có 2 phát biểu đúng.
}
\end{ex}

% ===== Câu 105 =====
% ID: L12C3B16-21-TN
% Mức: NB
\dangbt{Nhận biết từ thông và tính từ thông qua một diện tích}
\begin{ex}
% ID: L12C3B16-21-TN
% Mức: NB
Phát biểu nào sau đây đúng về Nhận biết từ thông và tính từ thông qua một diện tích?
\choice
{Từ thông luôn bằng BS sin(alpha), với alpha là góc giữa $B$ và pháp tuyến.}
{Có thể bỏ qua phương, chiều và đơn vị trong mọi trường hợp.}
{Kết quả không phụ thuộc dữ kiện và điều kiện áp dụng.}
{\True Từ thông qua diện tích $S$ trong từ trường đều là Phi = BS cos(alpha), với alpha là góc giữa $B$ và pháp tuyến.}
\loigiai{
Từ thông qua diện tích $S$ trong từ trường đều là Phi = BS cos(alpha), với alpha là góc giữa $B$ và pháp tuyến. Vì vậy phương án $D$ đúng; các phương án còn lại sai.
}
\end{ex}

% ===== Câu 106 =====
% ID: L12C3B16-22-DS
% Mức: TH
\dangbt{Tính suất điện động cảm ứng theo định luật Faraday}
\begin{ex}
% ID: L12C3B16-22-DS
% Mức: TH
Xét Tính suất điện động cảm ứng theo định luật Faraday (Tình huống 2). Hãy xác định tính đúng, sai của bốn phát biểu sau.
\choiceTF
{Suất điện động cảm ứng bằng $N$ Delta t/Delta Phi.}
{\True Độ lớn suất điện động cảm ứng trung bình của cuộn $N$ vòng là |e| = N|Delta Phi|/Delta t.}
{Kết quả không phụ thuộc dữ kiện và có thể bỏ qua điều kiện.}
{\True Khi giải cần xác định đúng dữ kiện, phương chiều, điều kiện áp dụng và đơn vị.}
\loigiai{
\begin{itemchoice}
\item (Sai) Suất điện động cảm ứng bằng $N$ Delta t/Delta Phi. (Vì): uất điện động cảm ứng bằng $N$ Delta t/Delta Phi. b) Đúng: Độ lớn suất điện động cảm ứng trung bình của cuộn $N$ vòng là |e| = N|Delta Phi|/Delta t. c) Sai: Kết quả không phụ thuộc dữ kiện và có thể bỏ qua điều kiện. d) Đúng: Khi giải cần xác định đúng dữ kiện, phương chiều, điều kiện áp dụng và đơn vị. Kiến thức cốt lõi: Độ lớn suất điện động cảm ứng trung bình của cuộn $N$ vòng là |e| = N|Delta Phi|/Delta t
\item (Đúng) Độ lớn suất điện động cảm ứng trung bình của cuộn $N$ vòng là |e| = N|Delta Phi|/Delta t.
\item (Sai) Kết quả không phụ thuộc dữ kiện và có thể bỏ qua điều kiện.
\item (Đúng) Khi giải cần xác định đúng dữ kiện, phương chiều, điều kiện áp dụng và đơn vị.
\end{itemchoice}
}
\end{ex}

% ===== Câu 107 =====
% ID: L12C3B16-22-TL
% Mức: VD
\dangbt{Phân tích sự biến thiên từ thông qua khung dây}
\begin{bt}
% ID: L12C3B16-22-TL
% Mức: VD
Đề xuất các bước giải một bài toán thuộc dạng Phân tích sự biến thiên từ thông qua khung dây.
\loigiai{
Cần nêu được: Từ thông biến thiên khi $B$, diện tích mạch hoặc góc giữa $B$ và pháp tuyến thay đổi. Khi vận dụng phải xác định đúng dữ kiện, phương chiều nếu có, điều kiện áp dụng, hệ đơn vị và kiểm tra tính hợp lí. Nhận định “Từ thông không đổi trong mọi chuyển động của khung dây.” sai.
}
\end{bt}

% ===== Câu 108 =====
% ID: L12C3B16-22-TLN
% Mức: VD
\begin{ex}
% ID: L12C3B16-22-TLN
% Mức: VD
Trong bốn phát biểu về Phân tích sự biến thiên từ thông qua khung dây, có bao nhiêu phát biểu đúng? (Chỉ ghi một số nguyên.) (1) Từ thông không đổi trong mọi chuyển động của khung dây. (2) Từ thông biến thiên khi $B$, diện tích mạch hoặc góc giữa $B$ và pháp tuyến thay đổi. (3) Kết quả phải phù hợp ý nghĩa vật lí. (4) Mọi đại lượng đều là vô hướng.
\shortans{2}
\loigiai{
Đối chiếu với kiến thức: Từ thông biến thiên khi $B$, diện tích mạch hoặc góc giữa $B$ và pháp tuyến thay đổi. Có 2 phát biểu đúng.
}
\end{ex}

% ===== Câu 109 =====
% ID: L12C3B16-22-TN
% Mức: NB
\dangbt{Nhận biết từ thông và tính từ thông qua một diện tích}
\begin{ex}
% ID: L12C3B16-22-TN
% Mức: NB
Trong nội dung Nhận biết từ thông và tính từ thông qua một diện tích?
\choice
{\True Từ thông qua diện tích $S$ trong từ trường đều là Phi = BS cos(alpha), với alpha là góc giữa $B$ và pháp tuyến.}
{Từ thông luôn bằng BS sin(alpha), với alpha là góc giữa $B$ và pháp tuyến.}
{Có thể bỏ qua phương, chiều và đơn vị trong mọi trường hợp.}
{Kết quả không phụ thuộc dữ kiện và điều kiện áp dụng.}
\loigiai{
Từ thông qua diện tích $S$ trong từ trường đều là Phi = BS cos(alpha), với alpha là góc giữa $B$ và pháp tuyến. Vì vậy phương án $A$ đúng; các phương án còn lại sai.
}
\end{ex}

% ===== Câu 110 =====
% ID: L12C3B16-23-DS
% Mức: VD
\dangbt{Tính suất điện động cảm ứng theo định luật Faraday}
\begin{ex}
% ID: L12C3B16-23-DS
% Mức: VD
Xét Tính suất điện động cảm ứng theo định luật Faraday (Tình huống 3). Hãy xác định tính đúng, sai của bốn phát biểu sau.
\choiceTF
{\True Khi giải cần xác định đúng dữ kiện, phương chiều, điều kiện áp dụng và đơn vị.}
{Kết quả không phụ thuộc dữ kiện và có thể bỏ qua điều kiện.}
{\True Độ lớn suất điện động cảm ứng trung bình của cuộn $N$ vòng là |e| = N|Delta Phi|/Delta t.}
{Suất điện động cảm ứng bằng $N$ Delta t/Delta Phi.}
\loigiai{
\begin{itemchoice}
\item (Đúng) Khi giải cần xác định đúng dữ kiện, phương chiều, điều kiện áp dụng và đơn vị. (Vì): Khi giải cần xác định đúng dữ kiện, phương chiều, điều kiện áp dụng và đơn vị. b) Sai: Kết quả không phụ thuộc dữ kiện và có thể bỏ qua điều kiện. c) Đúng: Độ lớn suất điện động cảm ứng trung bình của cuộn $N$ vòng là |e| = N|Delta Phi|/Delta t. d) Sai: Suất điện động cảm ứng bằng $N$ Delta t/Delta Phi. Kiến thức cốt lõi: Độ lớn suất điện động cảm ứng trung bình của cuộn $N$ vòng là |e| = N|Delta Phi|/Delta t
\item (Sai) Kết quả không phụ thuộc dữ kiện và có thể bỏ qua điều kiện.
\item (Đúng) Độ lớn suất điện động cảm ứng trung bình của cuộn $N$ vòng là |e| = N|Delta Phi|/Delta t.
\item (Sai) Suất điện động cảm ứng bằng $N$ Delta t/Delta Phi.
\end{itemchoice}
}
\end{ex}

% ===== Câu 111 =====
% ID: L12C3B16-23-TL
% Mức: VD
\dangbt{Phân tích sự biến thiên từ thông qua khung dây}
\begin{bt}
% ID: L12C3B16-23-TL
% Mức: VD
Nêu một tình huống thực tiễn liên quan đến Phân tích sự biến thiên từ thông qua khung dây và giải thích bằng kiến thức Vật lí.
\loigiai{
Cần nêu được: Từ thông biến thiên khi $B$, diện tích mạch hoặc góc giữa $B$ và pháp tuyến thay đổi. Khi vận dụng phải xác định đúng dữ kiện, phương chiều nếu có, điều kiện áp dụng, hệ đơn vị và kiểm tra tính hợp lí. Nhận định “Từ thông không đổi trong mọi chuyển động của khung dây.” sai.
}
\end{bt}

% ===== Câu 112 =====
% ID: L12C3B16-23-TLN
% Mức: VD
\begin{ex}
% ID: L12C3B16-23-TLN
% Mức: VD
Trong bốn phát biểu về Phân tích sự biến thiên từ thông qua khung dây, có bao nhiêu phát biểu đúng? (Chỉ ghi một số nguyên.) (1) Từ thông biến thiên khi $B$, diện tích mạch hoặc góc giữa $B$ và pháp tuyến thay đổi. (2) Từ thông không đổi trong mọi chuyển động của khung dây. (3) Cần kiểm tra điều kiện áp dụng. (4) Có thể bỏ qua đơn vị.
\shortans{2}
\loigiai{
Đối chiếu với kiến thức: Từ thông biến thiên khi $B$, diện tích mạch hoặc góc giữa $B$ và pháp tuyến thay đổi. Có 2 phát biểu đúng.
}
\end{ex}

% ===== Câu 113 =====
% ID: L12C3B16-23-TN
% Mức: NB
\dangbt{Nhận biết từ thông và tính từ thông qua một diện tích}
\begin{ex}
% ID: L12C3B16-23-TN
% Mức: NB
Tình huống 3: chọn kết luận chính xác về Nhận biết từ thông và tính từ thông qua một diện tích?
\choice
{Kết quả không phụ thuộc dữ kiện và điều kiện áp dụng.}
{\True Từ thông qua diện tích $S$ trong từ trường đều là Phi = BS cos(alpha), với alpha là góc giữa $B$ và pháp tuyến.}
{Từ thông luôn bằng BS sin(alpha), với alpha là góc giữa $B$ và pháp tuyến.}
{Có thể bỏ qua phương, chiều và đơn vị trong mọi trường hợp.}
\loigiai{
Từ thông qua diện tích $S$ trong từ trường đều là Phi = BS cos(alpha), với alpha là góc giữa $B$ và pháp tuyến. Vì vậy phương án $B$ đúng; các phương án còn lại sai.
}
\end{ex}

% ===== Câu 114 =====
% ID: L12C3B16-24-DS
% Mức: VD
\dangbt{Tính suất điện động cảm ứng theo định luật Faraday}
\begin{ex}
% ID: L12C3B16-24-DS
% Mức: VD
Xét Tính suất điện động cảm ứng theo định luật Faraday (Tình huống 4). Hãy xác định tính đúng, sai của bốn phát biểu sau.
\choiceTF
{Kết quả không phụ thuộc dữ kiện và có thể bỏ qua điều kiện.}
{\True Khi giải cần xác định đúng dữ kiện, phương chiều, điều kiện áp dụng và đơn vị.}
{Suất điện động cảm ứng bằng $N$ Delta t/Delta Phi.}
{\True Độ lớn suất điện động cảm ứng trung bình của cuộn $N$ vòng là |e| = N|Delta Phi|/Delta t.}
\loigiai{
\begin{itemchoice}
\item (Sai) Kết quả không phụ thuộc dữ kiện và có thể bỏ qua điều kiện. (Vì): Kết quả không phụ thuộc dữ kiện và có thể bỏ qua điều kiện. b) Đúng: Khi giải cần xác định đúng dữ kiện, phương chiều, điều kiện áp dụng và đơn vị. c) Sai: Suất điện động cảm ứng bằng $N$ Delta t/Delta Phi. d) Đúng: Độ lớn suất điện động cảm ứng trung bình của cuộn $N$ vòng là |e| = N|Delta Phi|/Delta t. Kiến thức cốt lõi: Độ lớn suất điện động cảm ứng trung bình của cuộn $N$ vòng là |e| = N|Delta Phi|/Delta t
\item (Đúng) Khi giải cần xác định đúng dữ kiện, phương chiều, điều kiện áp dụng và đơn vị.
\item (Sai) Suất điện động cảm ứng bằng $N$ Delta t/Delta Phi.
\item (Đúng) Độ lớn suất điện động cảm ứng trung bình của cuộn $N$ vòng là |e| = N|Delta Phi|/Delta t.
\end{itemchoice}
}
\end{ex}

% ===== Câu 115 =====
% ID: L12C3B16-24-TL
% Mức: VDC
\dangbt{Phân tích sự biến thiên từ thông qua khung dây}
\begin{bt}
% ID: L12C3B16-24-TL
% Mức: VDC
So sánh nhận định đúng “Từ thông biến thiên khi $B$, diện tích mạch hoặc góc giữa $B$ và pháp tuyến thay đổi.” với nhận định sai “Từ thông không đổi trong mọi chuyển động của khung dây.”; chỉ rõ căn cứ Vật lí.
\loigiai{
Cần nêu được: Từ thông biến thiên khi $B$, diện tích mạch hoặc góc giữa $B$ và pháp tuyến thay đổi. Khi vận dụng phải xác định đúng dữ kiện, phương chiều nếu có, điều kiện áp dụng, hệ đơn vị và kiểm tra tính hợp lí. Nhận định “Từ thông không đổi trong mọi chuyển động của khung dây.” sai.
}
\end{bt}

% ===== Câu 116 =====
% ID: L12C3B16-24-TLN
% Mức: VDC
\begin{ex}
% ID: L12C3B16-24-TLN
% Mức: VDC
Trong bốn phát biểu về Phân tích sự biến thiên từ thông qua khung dây, có bao nhiêu phát biểu đúng? (Chỉ ghi một số nguyên.) (1) Khi giải cần xác định đúng dữ kiện, phương chiều, điều kiện áp dụng và đơn vị. (2) Từ thông không đổi trong mọi chuyển động của khung dây. (3) Phải dùng đúng định luật cho tình huống. (4) Từ thông biến thiên khi $B$, diện tích mạch hoặc góc giữa $B$ và pháp tuyến thay đổi.
\shortans{3}
\loigiai{
Đối chiếu với kiến thức: Từ thông biến thiên khi $B$, diện tích mạch hoặc góc giữa $B$ và pháp tuyến thay đổi. Có 3 phát biểu đúng.
}
\end{ex}

% ===== Câu 117 =====
% ID: L12C3B16-24-TN
% Mức: TH
\dangbt{Nhận biết từ thông và tính từ thông qua một diện tích}
\begin{ex}
% ID: L12C3B16-24-TN
% Mức: TH
Nội dung nào mô tả đúng nhất Nhận biết từ thông và tính từ thông qua một diện tích?
\choice
{Có thể bỏ qua phương, chiều và đơn vị trong mọi trường hợp.}
{Kết quả không phụ thuộc dữ kiện và điều kiện áp dụng.}
{\True Từ thông qua diện tích $S$ trong từ trường đều là Phi = BS cos(alpha), với alpha là góc giữa $B$ và pháp tuyến.}
{Từ thông luôn bằng BS sin(alpha), với alpha là góc giữa $B$ và pháp tuyến.}
\loigiai{
Từ thông qua diện tích $S$ trong từ trường đều là Phi = BS cos(alpha), với alpha là góc giữa $B$ và pháp tuyến. Vì vậy phương án $C$ đúng; các phương án còn lại sai.
}
\end{ex}

% ===== Câu 118 =====
% ID: L12C3B16-25-DS
% Mức: VD
\dangbt{Tính suất điện động cảm ứng theo định luật Faraday}
\begin{ex}
% ID: L12C3B16-25-DS
% Mức: VD
Xét Tính suất điện động cảm ứng theo định luật Faraday (Tình huống 5). Hãy xác định tính đúng, sai của bốn phát biểu sau.
\choiceTF
{\True Độ lớn suất điện động cảm ứng trung bình của cuộn $N$ vòng là |e| = N|Delta Phi|/Delta t.}
{Kết quả không phụ thuộc dữ kiện và có thể bỏ qua điều kiện.}
{Suất điện động cảm ứng bằng $N$ Delta t/Delta Phi.}
{\True Khi giải cần xác định đúng dữ kiện, phương chiều, điều kiện áp dụng và đơn vị.}
\loigiai{
\begin{itemchoice}
\item (Đúng) Độ lớn suất điện động cảm ứng trung bình của cuộn $N$ vòng là |e| = N|Delta Phi|/Delta t. (Vì): ộ lớn suất điện động cảm ứng trung bình của cuộn $N$ vòng là |e| = N|Delta Phi|/Delta t. b) Sai: Kết quả không phụ thuộc dữ kiện và có thể bỏ qua điều kiện. c) Sai: Suất điện động cảm ứng bằng $N$ Delta t/Delta Phi. d) Đúng: Khi giải cần xác định đúng dữ kiện, phương chiều, điều kiện áp dụng và đơn vị. Kiến thức cốt lõi: Độ lớn suất điện động cảm ứng trung bình của cuộn $N$ vòng là |e| = N|Delta Phi|/Delta t
\item (Sai) Kết quả không phụ thuộc dữ kiện và có thể bỏ qua điều kiện.
\item (Sai) Suất điện động cảm ứng bằng $N$ Delta t/Delta Phi.
\item (Đúng) Khi giải cần xác định đúng dữ kiện, phương chiều, điều kiện áp dụng và đơn vị.
\end{itemchoice}
}
\end{ex}

% ===== Câu 119 =====
% ID: L12C3B16-25-TL
% Mức: TH
\begin{bt}
% ID: L12C3B16-25-TL
% Mức: TH
Trình bày kiến thức cốt lõi của Tính suất điện động cảm ứng theo định luật Faraday và nêu điều kiện áp dụng.
\loigiai{
Cần nêu được: Độ lớn suất điện động cảm ứng trung bình của cuộn $N$ vòng là |e| = N|Delta Phi|/Delta t. Khi vận dụng phải xác định đúng dữ kiện, phương chiều nếu có, điều kiện áp dụng, hệ đơn vị và kiểm tra tính hợp lí. Nhận định “Suất điện động cảm ứng bằng $N$ Delta t/Delta Phi.” sai.
}
\end{bt}

% ===== Câu 120 =====
% ID: L12C3B16-25-TLN
% Mức: TH
\begin{ex}
% ID: L12C3B16-25-TLN
% Mức: TH
Trong bốn phát biểu về Tính suất điện động cảm ứng theo định luật Faraday, có bao nhiêu phát biểu đúng? (Chỉ ghi một số nguyên.) (1) Độ lớn suất điện động cảm ứng trung bình của cuộn $N$ vòng là |e| = N|Delta Phi|/Delta t. (2) Suất điện động cảm ứng bằng $N$ Delta t/Delta Phi. (3) Cần kiểm tra điều kiện áp dụng. (4) Có thể bỏ qua đơn vị.
\shortans{2}
\loigiai{
Đối chiếu với kiến thức: Độ lớn suất điện động cảm ứng trung bình của cuộn $N$ vòng là |e| = N|Delta Phi|/Delta t. Có 2 phát biểu đúng.
}
\end{ex}

% ===== Câu 121 =====
% ID: L12C3B16-25-TN
% Mức: TH
\dangbt{Nhận biết từ thông và tính từ thông qua một diện tích}
\begin{ex}
% ID: L12C3B16-25-TN
% Mức: TH
Khi phân tích Nhận biết từ thông và tính từ thông qua một diện tích?
\choice
{Từ thông luôn bằng BS sin(alpha), với alpha là góc giữa $B$ và pháp tuyến.}
{Có thể bỏ qua phương, chiều và đơn vị trong mọi trường hợp.}
{Kết quả không phụ thuộc dữ kiện và điều kiện áp dụng.}
{\True Từ thông qua diện tích $S$ trong từ trường đều là Phi = BS cos(alpha), với alpha là góc giữa $B$ và pháp tuyến.}
\loigiai{
Từ thông qua diện tích $S$ trong từ trường đều là Phi = BS cos(alpha), với alpha là góc giữa $B$ và pháp tuyến. Vì vậy phương án $D$ đúng; các phương án còn lại sai.
}
\end{ex}

% ===== Câu 122 =====
% ID: L12C3B16-26-DS
% Mức: TH
\dangbt{Xác định chiều dòng điện cảm ứng bằng định luật Lenz}
\begin{ex}
% ID: L12C3B16-26-DS
% Mức: TH
Xét Xác định chiều dòng điện cảm ứng bằng định luật Lenz (Tình huống 1). Hãy xác định tính đúng, sai của bốn phát biểu sau.
\choiceTF
{\True Dòng điện cảm ứng có chiều sao cho từ trường do nó sinh ra chống lại sự biến thiên từ thông đã gây ra nó.}
{Dòng điện cảm ứng luôn làm tăng biến thiên từ thông ban đầu.}
{\True Khi giải cần xác định đúng dữ kiện, phương chiều, điều kiện áp dụng và đơn vị.}
{Kết quả không phụ thuộc dữ kiện và có thể bỏ qua điều kiện.}
\loigiai{
\begin{itemchoice}
\item (Đúng) Dòng điện cảm ứng có chiều sao cho từ trường do nó sinh ra chống lại sự biến thiên từ thông đã gây ra nó. (Vì): òng điện cảm ứng có chiều sao cho từ trường do nó sinh ra chống lại sự biến thiên từ thông đã gây ra nó. b) Sai: Dòng điện cảm ứng luôn làm tăng biến thiên từ thông ban đầu. c) Đúng: Khi giải cần xác định đúng dữ kiện, phương chiều, điều kiện áp dụng và đơn vị. d) Sai: Kết quả không phụ thuộc dữ kiện và có thể bỏ qua điều kiện. Kiến thức cốt lõi: Dòng điện cảm ứng có chiều sao cho từ trường do nó sinh ra chống lại sự biến thiên từ thông đã gây ra nó
\item (Sai) Dòng điện cảm ứng luôn làm tăng biến thiên từ thông ban đầu.
\item (Đúng) Khi giải cần xác định đúng dữ kiện, phương chiều, điều kiện áp dụng và đơn vị.
\item (Sai) Kết quả không phụ thuộc dữ kiện và có thể bỏ qua điều kiện.
\end{itemchoice}
}
\end{ex}

% ===== Câu 123 =====
% ID: L12C3B16-26-TL
% Mức: TH
\dangbt{Tính suất điện động cảm ứng theo định luật Faraday}
\begin{bt}
% ID: L12C3B16-26-TL
% Mức: TH
Giải thích vì sao nhận định sau đúng: “Độ lớn suất điện động cảm ứng trung bình của cuộn $N$ vòng là |e| = N|Delta Phi|/Delta t.”
\loigiai{
Cần nêu được: Độ lớn suất điện động cảm ứng trung bình của cuộn $N$ vòng là |e| = N|Delta Phi|/Delta t. Khi vận dụng phải xác định đúng dữ kiện, phương chiều nếu có, điều kiện áp dụng, hệ đơn vị và kiểm tra tính hợp lí. Nhận định “Suất điện động cảm ứng bằng $N$ Delta t/Delta Phi.” sai.
}
\end{bt}

% ===== Câu 124 =====
% ID: L12C3B16-26-TLN
% Mức: TH
\begin{ex}
% ID: L12C3B16-26-TLN
% Mức: TH
Trong bốn phát biểu về Tính suất điện động cảm ứng theo định luật Faraday, có bao nhiêu phát biểu đúng? (Chỉ ghi một số nguyên.) (1) Suất điện động cảm ứng bằng $N$ Delta t/Delta Phi. (2) Độ lớn suất điện động cảm ứng trung bình của cuộn $N$ vòng là |e| = N|Delta Phi|/Delta t. (3) Kết quả phải phù hợp ý nghĩa vật lí. (4) Mọi đại lượng đều là vô hướng.
\shortans{1}
\loigiai{
Đối chiếu với kiến thức: Độ lớn suất điện động cảm ứng trung bình của cuộn $N$ vòng là |e| = N|Delta Phi|/Delta t. Có 1 phát biểu đúng.
}
\end{ex}

% ===== Câu 125 =====
% ID: L12C3B16-26-TN
% Mức: TH
\dangbt{Nhận biết từ thông và tính từ thông qua một diện tích}
\begin{ex}
% ID: L12C3B16-26-TN
% Mức: TH
Một học sinh ôn tập Nhận biết từ thông và tính từ thông qua một diện tích?
\choice
{\True Từ thông qua diện tích $S$ trong từ trường đều là Phi = BS cos(alpha), với alpha là góc giữa $B$ và pháp tuyến.}
{Từ thông luôn bằng BS sin(alpha), với alpha là góc giữa $B$ và pháp tuyến.}
{Có thể bỏ qua phương, chiều và đơn vị trong mọi trường hợp.}
{Kết quả không phụ thuộc dữ kiện và điều kiện áp dụng.}
\loigiai{
Từ thông qua diện tích $S$ trong từ trường đều là Phi = BS cos(alpha), với alpha là góc giữa $B$ và pháp tuyến. Vì vậy phương án $A$ đúng; các phương án còn lại sai.
}
\end{ex}

% ===== Câu 126 =====
% ID: L12C3B16-27-DS
% Mức: TH
\dangbt{Xác định chiều dòng điện cảm ứng bằng định luật Lenz}
\begin{ex}
% ID: L12C3B16-27-DS
% Mức: TH
Xét Xác định chiều dòng điện cảm ứng bằng định luật Lenz (Tình huống 2). Hãy xác định tính đúng, sai của bốn phát biểu sau.
\choiceTF
{Dòng điện cảm ứng luôn làm tăng biến thiên từ thông ban đầu.}
{\True Dòng điện cảm ứng có chiều sao cho từ trường do nó sinh ra chống lại sự biến thiên từ thông đã gây ra nó.}
{Kết quả không phụ thuộc dữ kiện và có thể bỏ qua điều kiện.}
{\True Khi giải cần xác định đúng dữ kiện, phương chiều, điều kiện áp dụng và đơn vị.}
\loigiai{
\begin{itemchoice}
\item (Sai) Dòng điện cảm ứng luôn làm tăng biến thiên từ thông ban đầu. (Vì): Dòng điện cảm ứng luôn làm tăng biến thiên từ thông ban đầu. b) Đúng: Dòng điện cảm ứng có chiều sao cho từ trường do nó sinh ra chống lại sự biến thiên từ thông đã gây ra nó. c) Sai: Kết quả không phụ thuộc dữ kiện và có thể bỏ qua điều kiện. d) Đúng: Khi giải cần xác định đúng dữ kiện, phương chiều, điều kiện áp dụng và đơn vị. Kiến thức cốt lõi: Dòng điện cảm ứng có chiều sao cho từ trường do nó sinh ra chống lại sự biến thiên từ thông đã gây ra nó
\item (Đúng) Dòng điện cảm ứng có chiều sao cho từ trường do nó sinh ra chống lại sự biến thiên từ thông đã gây ra nó.
\item (Sai) Kết quả không phụ thuộc dữ kiện và có thể bỏ qua điều kiện.
\item (Đúng) Khi giải cần xác định đúng dữ kiện, phương chiều, điều kiện áp dụng và đơn vị.
\end{itemchoice}
}
\end{ex}

% ===== Câu 127 =====
% ID: L12C3B16-27-TL
% Mức: VD
\dangbt{Tính suất điện động cảm ứng theo định luật Faraday}
\begin{bt}
% ID: L12C3B16-27-TL
% Mức: VD
Phân tích sai lầm trong nhận định: “Suất điện động cảm ứng bằng $N$ Delta t/Delta Phi.”
\loigiai{
Cần nêu được: Độ lớn suất điện động cảm ứng trung bình của cuộn $N$ vòng là |e| = N|Delta Phi|/Delta t. Khi vận dụng phải xác định đúng dữ kiện, phương chiều nếu có, điều kiện áp dụng, hệ đơn vị và kiểm tra tính hợp lí. Nhận định “Suất điện động cảm ứng bằng $N$ Delta t/Delta Phi.” sai.
}
\end{bt}

% ===== Câu 128 =====
% ID: L12C3B16-27-TLN
% Mức: TH
\begin{ex}
% ID: L12C3B16-27-TLN
% Mức: TH
Trong bốn phát biểu về Tính suất điện động cảm ứng theo định luật Faraday, có bao nhiêu phát biểu đúng? (Chỉ ghi một số nguyên.) (1) Độ lớn suất điện động cảm ứng trung bình của cuộn $N$ vòng là |e| = N|Delta Phi|/Delta t. (2) Suất điện động cảm ứng bằng $N$ Delta t/Delta Phi. (3) Cần kiểm tra điều kiện áp dụng. (4) Có thể bỏ qua đơn vị.
\shortans{2}
\loigiai{
Đối chiếu với kiến thức: Độ lớn suất điện động cảm ứng trung bình của cuộn $N$ vòng là |e| = N|Delta Phi|/Delta t. Có 2 phát biểu đúng.
}
\end{ex}

% ===== Câu 129 =====
% ID: L12C3B16-27-TN
% Mức: TH
\dangbt{Nhận biết từ thông và tính từ thông qua một diện tích}
\begin{ex}
% ID: L12C3B16-27-TN
% Mức: TH
Chọn phát biểu không trái với kiến thức về Nhận biết từ thông và tính từ thông qua một diện tích?
\choice
{Kết quả không phụ thuộc dữ kiện và điều kiện áp dụng.}
{\True Từ thông qua diện tích $S$ trong từ trường đều là Phi = BS cos(alpha), với alpha là góc giữa $B$ và pháp tuyến.}
{Từ thông luôn bằng BS sin(alpha), với alpha là góc giữa $B$ và pháp tuyến.}
{Có thể bỏ qua phương, chiều và đơn vị trong mọi trường hợp.}
\loigiai{
Từ thông qua diện tích $S$ trong từ trường đều là Phi = BS cos(alpha), với alpha là góc giữa $B$ và pháp tuyến. Vì vậy phương án $B$ đúng; các phương án còn lại sai.
}
\end{ex}

% ===== Câu 130 =====
% ID: L12C3B16-28-DS
% Mức: VD
\dangbt{Xác định chiều dòng điện cảm ứng bằng định luật Lenz}
\begin{ex}
% ID: L12C3B16-28-DS
% Mức: VD
Xét Xác định chiều dòng điện cảm ứng bằng định luật Lenz (Tình huống 3). Hãy xác định tính đúng, sai của bốn phát biểu sau.
\choiceTF
{\True Khi giải cần xác định đúng dữ kiện, phương chiều, điều kiện áp dụng và đơn vị.}
{Kết quả không phụ thuộc dữ kiện và có thể bỏ qua điều kiện.}
{\True Dòng điện cảm ứng có chiều sao cho từ trường do nó sinh ra chống lại sự biến thiên từ thông đã gây ra nó.}
{Dòng điện cảm ứng luôn làm tăng biến thiên từ thông ban đầu.}
\loigiai{
\begin{itemchoice}
\item (Đúng) Khi giải cần xác định đúng dữ kiện, phương chiều, điều kiện áp dụng và đơn vị. (Vì): Khi giải cần xác định đúng dữ kiện, phương chiều, điều kiện áp dụng và đơn vị. b) Sai: Kết quả không phụ thuộc dữ kiện và có thể bỏ qua điều kiện. c) Đúng: Dòng điện cảm ứng có chiều sao cho từ trường do nó sinh ra chống lại sự biến thiên từ thông đã gây ra nó. d) Sai: Dòng điện cảm ứng luôn làm tăng biến thiên từ thông ban đầu. Kiến thức cốt lõi: Dòng điện cảm ứng có chiều sao cho từ trường do nó sinh ra chống lại sự biến thiên từ thông đã gây ra nó
\item (Sai) Kết quả không phụ thuộc dữ kiện và có thể bỏ qua điều kiện.
\item (Đúng) Dòng điện cảm ứng có chiều sao cho từ trường do nó sinh ra chống lại sự biến thiên từ thông đã gây ra nó.
\item (Sai) Dòng điện cảm ứng luôn làm tăng biến thiên từ thông ban đầu.
\end{itemchoice}
}
\end{ex}

% ===== Câu 131 =====
% ID: L12C3B16-28-TL
% Mức: VD
\dangbt{Tính suất điện động cảm ứng theo định luật Faraday}
\begin{bt}
% ID: L12C3B16-28-TL
% Mức: VD
Đề xuất các bước giải một bài toán thuộc dạng Tính suất điện động cảm ứng theo định luật Faraday.
\loigiai{
Cần nêu được: Độ lớn suất điện động cảm ứng trung bình của cuộn $N$ vòng là |e| = N|Delta Phi|/Delta t. Khi vận dụng phải xác định đúng dữ kiện, phương chiều nếu có, điều kiện áp dụng, hệ đơn vị và kiểm tra tính hợp lí. Nhận định “Suất điện động cảm ứng bằng $N$ Delta t/Delta Phi.” sai.
}
\end{bt}

% ===== Câu 132 =====
% ID: L12C3B16-28-TLN
% Mức: VD
\begin{ex}
% ID: L12C3B16-28-TLN
% Mức: VD
Trong bốn phát biểu về Tính suất điện động cảm ứng theo định luật Faraday, có bao nhiêu phát biểu đúng? (Chỉ ghi một số nguyên.) (1) Suất điện động cảm ứng bằng $N$ Delta t/Delta Phi. (2) Độ lớn suất điện động cảm ứng trung bình của cuộn $N$ vòng là |e| = N|Delta Phi|/Delta t. (3) Kết quả phải phù hợp ý nghĩa vật lí. (4) Mọi đại lượng đều là vô hướng.
\shortans{2}
\loigiai{
Đối chiếu với kiến thức: Độ lớn suất điện động cảm ứng trung bình của cuộn $N$ vòng là |e| = N|Delta Phi|/Delta t. Có 2 phát biểu đúng.
}
\end{ex}

% ===== Câu 133 =====
% ID: L12C3B16-28-TN
% Mức: VD
\dangbt{Nhận biết từ thông và tính từ thông qua một diện tích}
\begin{ex}
% ID: L12C3B16-28-TN
% Mức: VD
Cơ sở Vật lí đúng dùng để giải Nhận biết từ thông và tính từ thông qua một diện tích?
\choice
{Có thể bỏ qua phương, chiều và đơn vị trong mọi trường hợp.}
{Kết quả không phụ thuộc dữ kiện và điều kiện áp dụng.}
{\True Từ thông qua diện tích $S$ trong từ trường đều là Phi = BS cos(alpha), với alpha là góc giữa $B$ và pháp tuyến.}
{Từ thông luôn bằng BS sin(alpha), với alpha là góc giữa $B$ và pháp tuyến.}
\loigiai{
Từ thông qua diện tích $S$ trong từ trường đều là Phi = BS cos(alpha), với alpha là góc giữa $B$ và pháp tuyến. Vì vậy phương án $C$ đúng; các phương án còn lại sai.
}
\end{ex}

% ===== Câu 134 =====
% ID: L12C3B16-29-DS
% Mức: VD
\dangbt{Xác định chiều dòng điện cảm ứng bằng định luật Lenz}
\begin{ex}
% ID: L12C3B16-29-DS
% Mức: VD
Xét Xác định chiều dòng điện cảm ứng bằng định luật Lenz (Tình huống 4). Hãy xác định tính đúng, sai của bốn phát biểu sau.
\choiceTF
{Kết quả không phụ thuộc dữ kiện và có thể bỏ qua điều kiện.}
{\True Khi giải cần xác định đúng dữ kiện, phương chiều, điều kiện áp dụng và đơn vị.}
{Dòng điện cảm ứng luôn làm tăng biến thiên từ thông ban đầu.}
{\True Dòng điện cảm ứng có chiều sao cho từ trường do nó sinh ra chống lại sự biến thiên từ thông đã gây ra nó.}
\loigiai{
\begin{itemchoice}
\item (Sai) Kết quả không phụ thuộc dữ kiện và có thể bỏ qua điều kiện. (Vì): Kết quả không phụ thuộc dữ kiện và có thể bỏ qua điều kiện. b) Đúng: Khi giải cần xác định đúng dữ kiện, phương chiều, điều kiện áp dụng và đơn vị. c) Sai: Dòng điện cảm ứng luôn làm tăng biến thiên từ thông ban đầu. d) Đúng: Dòng điện cảm ứng có chiều sao cho từ trường do nó sinh ra chống lại sự biến thiên từ thông đã gây ra nó. Kiến thức cốt lõi: Dòng điện cảm ứng có chiều sao cho từ trường do nó sinh ra chống lại sự biến thiên từ thông đã gây ra nó
\item (Đúng) Khi giải cần xác định đúng dữ kiện, phương chiều, điều kiện áp dụng và đơn vị.
\item (Sai) Dòng điện cảm ứng luôn làm tăng biến thiên từ thông ban đầu.
\item (Đúng) Dòng điện cảm ứng có chiều sao cho từ trường do nó sinh ra chống lại sự biến thiên từ thông đã gây ra nó.
\end{itemchoice}
}
\end{ex}

% ===== Câu 135 =====
% ID: L12C3B16-29-TL
% Mức: VD
\dangbt{Tính suất điện động cảm ứng theo định luật Faraday}
\begin{bt}
% ID: L12C3B16-29-TL
% Mức: VD
Nêu một tình huống thực tiễn liên quan đến Tính suất điện động cảm ứng theo định luật Faraday và giải thích bằng kiến thức Vật lí.
\loigiai{
Cần nêu được: Độ lớn suất điện động cảm ứng trung bình của cuộn $N$ vòng là |e| = N|Delta Phi|/Delta t. Khi vận dụng phải xác định đúng dữ kiện, phương chiều nếu có, điều kiện áp dụng, hệ đơn vị và kiểm tra tính hợp lí. Nhận định “Suất điện động cảm ứng bằng $N$ Delta t/Delta Phi.” sai.
}
\end{bt}

% ===== Câu 136 =====
% ID: L12C3B16-29-TLN
% Mức: VD
\begin{ex}
% ID: L12C3B16-29-TLN
% Mức: VD
Trong bốn phát biểu về Tính suất điện động cảm ứng theo định luật Faraday, có bao nhiêu phát biểu đúng? (Chỉ ghi một số nguyên.) (1) Độ lớn suất điện động cảm ứng trung bình của cuộn $N$ vòng là |e| = N|Delta Phi|/Delta t. (2) Suất điện động cảm ứng bằng $N$ Delta t/Delta Phi. (3) Cần kiểm tra điều kiện áp dụng. (4) Có thể bỏ qua đơn vị.
\shortans{2}
\loigiai{
Đối chiếu với kiến thức: Độ lớn suất điện động cảm ứng trung bình của cuộn $N$ vòng là |e| = N|Delta Phi|/Delta t. Có 2 phát biểu đúng.
}
\end{ex}

% ===== Câu 137 =====
% ID: L12C3B16-29-TN
% Mức: VD
\dangbt{Nhận biết từ thông và tính từ thông qua một diện tích}
\begin{ex}
% ID: L12C3B16-29-TN
% Mức: VD
Trong một bài thực hành về Nhận biết từ thông và tính từ thông qua một diện tích?
\choice
{Từ thông luôn bằng BS sin(alpha), với alpha là góc giữa $B$ và pháp tuyến.}
{Có thể bỏ qua phương, chiều và đơn vị trong mọi trường hợp.}
{Kết quả không phụ thuộc dữ kiện và điều kiện áp dụng.}
{\True Từ thông qua diện tích $S$ trong từ trường đều là Phi = BS cos(alpha), với alpha là góc giữa $B$ và pháp tuyến.}
\loigiai{
Từ thông qua diện tích $S$ trong từ trường đều là Phi = BS cos(alpha), với alpha là góc giữa $B$ và pháp tuyến. Vì vậy phương án $D$ đúng; các phương án còn lại sai.
}
\end{ex}

% ===== Câu 138 =====
% ID: L12C3B16-30-DS
% Mức: VD
\dangbt{Xác định chiều dòng điện cảm ứng bằng định luật Lenz}
\begin{ex}
% ID: L12C3B16-30-DS
% Mức: VD
Xét Xác định chiều dòng điện cảm ứng bằng định luật Lenz (Tình huống 5). Hãy xác định tính đúng, sai của bốn phát biểu sau.
\choiceTF
{\True Dòng điện cảm ứng có chiều sao cho từ trường do nó sinh ra chống lại sự biến thiên từ thông đã gây ra nó.}
{Kết quả không phụ thuộc dữ kiện và có thể bỏ qua điều kiện.}
{Dòng điện cảm ứng luôn làm tăng biến thiên từ thông ban đầu.}
{\True Khi giải cần xác định đúng dữ kiện, phương chiều, điều kiện áp dụng và đơn vị.}
\loigiai{
\begin{itemchoice}
\item (Đúng) Dòng điện cảm ứng có chiều sao cho từ trường do nó sinh ra chống lại sự biến thiên từ thông đã gây ra nó. (Vì): òng điện cảm ứng có chiều sao cho từ trường do nó sinh ra chống lại sự biến thiên từ thông đã gây ra nó. b) Sai: Kết quả không phụ thuộc dữ kiện và có thể bỏ qua điều kiện. c) Sai: Dòng điện cảm ứng luôn làm tăng biến thiên từ thông ban đầu. d) Đúng: Khi giải cần xác định đúng dữ kiện, phương chiều, điều kiện áp dụng và đơn vị. Kiến thức cốt lõi: Dòng điện cảm ứng có chiều sao cho từ trường do nó sinh ra chống lại sự biến thiên từ thông đã gây ra nó
\item (Sai) Kết quả không phụ thuộc dữ kiện và có thể bỏ qua điều kiện.
\item (Sai) Dòng điện cảm ứng luôn làm tăng biến thiên từ thông ban đầu.
\item (Đúng) Khi giải cần xác định đúng dữ kiện, phương chiều, điều kiện áp dụng và đơn vị.
\end{itemchoice}
}
\end{ex}

% ===== Câu 139 =====
% ID: L12C3B16-30-TL
% Mức: VDC
\dangbt{Tính suất điện động cảm ứng theo định luật Faraday}
\begin{bt}
% ID: L12C3B16-30-TL
% Mức: VDC
So sánh nhận định đúng “Độ lớn suất điện động cảm ứng trung bình của cuộn $N$ vòng là |e| = N|Delta Phi|/Delta t.” với nhận định sai “Suất điện động cảm ứng bằng $N$ Delta t/Delta Phi.”; chỉ rõ căn cứ Vật lí.
\loigiai{
Cần nêu được: Độ lớn suất điện động cảm ứng trung bình của cuộn $N$ vòng là |e| = N|Delta Phi|/Delta t. Khi vận dụng phải xác định đúng dữ kiện, phương chiều nếu có, điều kiện áp dụng, hệ đơn vị và kiểm tra tính hợp lí. Nhận định “Suất điện động cảm ứng bằng $N$ Delta t/Delta Phi.” sai.
}
\end{bt}

% ===== Câu 140 =====
% ID: L12C3B16-30-TLN
% Mức: VDC
\begin{ex}
% ID: L12C3B16-30-TLN
% Mức: VDC
Trong bốn phát biểu về Tính suất điện động cảm ứng theo định luật Faraday, có bao nhiêu phát biểu đúng? (Chỉ ghi một số nguyên.) (1) Khi giải cần xác định đúng dữ kiện, phương chiều, điều kiện áp dụng và đơn vị. (2) Suất điện động cảm ứng bằng $N$ Delta t/Delta Phi. (3) Phải dùng đúng định luật cho tình huống. (4) Độ lớn suất điện động cảm ứng trung bình của cuộn $N$ vòng là |e| = N|Delta Phi|/Delta t.
\shortans{3}
\loigiai{
Đối chiếu với kiến thức: Độ lớn suất điện động cảm ứng trung bình của cuộn $N$ vòng là |e| = N|Delta Phi|/Delta t. Có 3 phát biểu đúng.
}
\end{ex}

% ===== Câu 141 =====
% ID: L12C3B16-30-TN
% Mức: VD
\dangbt{Nhận biết từ thông và tính từ thông qua một diện tích}
\begin{ex}
% ID: L12C3B16-30-TN
% Mức: VD
Kết luận đúng khi vận dụng Nhận biết từ thông và tính từ thông qua một diện tích?
\choice
{\True Từ thông qua diện tích $S$ trong từ trường đều là Phi = BS cos(alpha), với alpha là góc giữa $B$ và pháp tuyến.}
{Từ thông luôn bằng BS sin(alpha), với alpha là góc giữa $B$ và pháp tuyến.}
{Có thể bỏ qua phương, chiều và đơn vị trong mọi trường hợp.}
{Kết quả không phụ thuộc dữ kiện và điều kiện áp dụng.}
\loigiai{
Từ thông qua diện tích $S$ trong từ trường đều là Phi = BS cos(alpha), với alpha là góc giữa $B$ và pháp tuyến. Vì vậy phương án $A$ đúng; các phương án còn lại sai.
}
\end{ex}

% ===== Câu 142 =====
% ID: L12C3B16-31-DS
% Mức: TH
\begin{ex}
% ID: L12C3B16-31-DS
% Mức: TH
Xét Nhận biết từ thông và tính từ thông qua một diện tích (Tình huống 1). Hãy xác định tính đúng, sai của bốn phát biểu sau.
\choiceTF
{\True Từ thông qua diện tích $S$ trong từ trường đều là Phi = BS cos(alpha), với alpha là góc giữa $B$ và pháp tuyến.}
{Từ thông luôn bằng BS sin(alpha), với alpha là góc giữa $B$ và pháp tuyến.}
{\True Khi giải cần xác định đúng dữ kiện, phương chiều, điều kiện áp dụng và đơn vị.}
{Kết quả không phụ thuộc dữ kiện và có thể bỏ qua điều kiện.}
\loigiai{
\begin{itemchoice}
\item (Đúng) Từ thông qua diện tích $S$ trong từ trường đều là Phi = BS cos(alpha), với alpha là góc giữa $B$ và pháp tuyến. (Vì): lpha), với alpha là góc giữa $B$ và pháp tuyến
\item (Sai) Từ thông luôn bằng BS sin(alpha), với alpha là góc giữa $B$ và pháp tuyến.
\item (Đúng) Khi giải cần xác định đúng dữ kiện, phương chiều, điều kiện áp dụng và đơn vị.
\item (Sai) Kết quả không phụ thuộc dữ kiện và có thể bỏ qua điều kiện.
\end{itemchoice}
}
\end{ex}

% ===== Câu 143 =====
% ID: L12C3B16-31-TL
% Mức: TH
\dangbt{Xác định chiều dòng điện cảm ứng bằng định luật Lenz}
\begin{bt}
% ID: L12C3B16-31-TL
% Mức: TH
Trình bày kiến thức cốt lõi của Xác định chiều dòng điện cảm ứng bằng định luật Lenz và nêu điều kiện áp dụng.
\loigiai{
Cần nêu được: Dòng điện cảm ứng có chiều sao cho từ trường do nó sinh ra chống lại sự biến thiên từ thông đã gây ra nó. Khi vận dụng phải xác định đúng dữ kiện, phương chiều nếu có, điều kiện áp dụng, hệ đơn vị và kiểm tra tính hợp lí. Nhận định “Dòng điện cảm ứng luôn làm tăng biến thiên từ thông ban đầu.” sai.
}
\end{bt}

% ===== Câu 144 =====
% ID: L12C3B16-31-TLN
% Mức: TH
\begin{ex}
% ID: L12C3B16-31-TLN
% Mức: TH
Trong bốn phát biểu về Xác định chiều dòng điện cảm ứng bằng định luật Lenz, có bao nhiêu phát biểu đúng? (Chỉ ghi một số nguyên.) (1) Dòng điện cảm ứng có chiều sao cho từ trường do nó sinh ra chống lại sự biến thiên từ thông đã gây ra nó. (2) Dòng điện cảm ứng luôn làm tăng biến thiên từ thông ban đầu. (3) Cần kiểm tra điều kiện áp dụng. (4) Có thể bỏ qua đơn vị.
\shortans{2}
\loigiai{
Đối chiếu với kiến thức: Dòng điện cảm ứng có chiều sao cho từ trường do nó sinh ra chống lại sự biến thiên từ thông đã gây ra nó. Có 2 phát biểu đúng.
}
\end{ex}

% ===== Câu 145 =====
% ID: L12C3B16-31-TN
% Mức: NB
\dangbt{Phân tích sự biến thiên từ thông qua khung dây}
\begin{ex}
% ID: L12C3B16-31-TN
% Mức: NB
Phát biểu nào sau đây đúng về Phân tích sự biến thiên từ thông qua khung dây?
\choice
{Từ thông không đổi trong mọi chuyển động của khung dây.}
{Có thể bỏ qua phương, chiều và đơn vị trong mọi trường hợp.}
{Kết quả không phụ thuộc dữ kiện và điều kiện áp dụng.}
{\True Từ thông biến thiên khi $B$, diện tích mạch hoặc góc giữa $B$ và pháp tuyến thay đổi.}
\loigiai{
Từ thông biến thiên khi $B$, diện tích mạch hoặc góc giữa $B$ và pháp tuyến thay đổi. Vì vậy phương án $D$ đúng; các phương án còn lại sai.
}
\end{ex}

% ===== Câu 146 =====
% ID: L12C3B16-32-DS
% Mức: TH
\dangbt{Nhận biết từ thông và tính từ thông qua một diện tích}
\begin{ex}
% ID: L12C3B16-32-DS
% Mức: TH
Xét Nhận biết từ thông và tính từ thông qua một diện tích (Tình huống 2). Hãy xác định tính đúng, sai của bốn phát biểu sau.
\choiceTF
{Từ thông luôn bằng BS sin(alpha), với alpha là góc giữa $B$ và pháp tuyến.}
{\True Từ thông qua diện tích $S$ trong từ trường đều là Phi = BS cos(alpha), với alpha là góc giữa $B$ và pháp tuyến.}
{Kết quả không phụ thuộc dữ kiện và có thể bỏ qua điều kiện.}
{\True Khi giải cần xác định đúng dữ kiện, phương chiều, điều kiện áp dụng và đơn vị.}
\loigiai{
\begin{itemchoice}
\item (Sai) Từ thông luôn bằng BS sin(alpha), với alpha là góc giữa $B$ và pháp tuyến. (Vì): lpha), với alpha là góc giữa $B$ và pháp tuyến
\item (Đúng) Từ thông qua diện tích $S$ trong từ trường đều là Phi = BS cos(alpha), với alpha là góc giữa $B$ và pháp tuyến.
\item (Sai) Kết quả không phụ thuộc dữ kiện và có thể bỏ qua điều kiện.
\item (Đúng) Khi giải cần xác định đúng dữ kiện, phương chiều, điều kiện áp dụng và đơn vị.
\end{itemchoice}
}
\end{ex}

% ===== Câu 147 =====
% ID: L12C3B16-32-TL
% Mức: TH
\dangbt{Xác định chiều dòng điện cảm ứng bằng định luật Lenz}
\begin{bt}
% ID: L12C3B16-32-TL
% Mức: TH
Giải thích vì sao nhận định sau đúng: “Dòng điện cảm ứng có chiều sao cho từ trường do nó sinh ra chống lại sự biến thiên từ thông đã gây ra nó.”
\loigiai{
Cần nêu được: Dòng điện cảm ứng có chiều sao cho từ trường do nó sinh ra chống lại sự biến thiên từ thông đã gây ra nó. Khi vận dụng phải xác định đúng dữ kiện, phương chiều nếu có, điều kiện áp dụng, hệ đơn vị và kiểm tra tính hợp lí. Nhận định “Dòng điện cảm ứng luôn làm tăng biến thiên từ thông ban đầu.” sai.
}
\end{bt}

% ===== Câu 148 =====
% ID: L12C3B16-32-TLN
% Mức: TH
\begin{ex}
% ID: L12C3B16-32-TLN
% Mức: TH
Trong bốn phát biểu về Xác định chiều dòng điện cảm ứng bằng định luật Lenz, có bao nhiêu phát biểu đúng? (Chỉ ghi một số nguyên.) (1) Dòng điện cảm ứng luôn làm tăng biến thiên từ thông ban đầu. (2) Dòng điện cảm ứng có chiều sao cho từ trường do nó sinh ra chống lại sự biến thiên từ thông đã gây ra nó. (3) Kết quả phải phù hợp ý nghĩa vật lí. (4) Mọi đại lượng đều là vô hướng.
\shortans{1}
\loigiai{
Đối chiếu với kiến thức: Dòng điện cảm ứng có chiều sao cho từ trường do nó sinh ra chống lại sự biến thiên từ thông đã gây ra nó. Có 1 phát biểu đúng.
}
\end{ex}

% ===== Câu 149 =====
% ID: L12C3B16-32-TN
% Mức: NB
\dangbt{Phân tích sự biến thiên từ thông qua khung dây}
\begin{ex}
% ID: L12C3B16-32-TN
% Mức: NB
Trong nội dung Phân tích sự biến thiên từ thông qua khung dây?
\choice
{\True Từ thông biến thiên khi $B$, diện tích mạch hoặc góc giữa $B$ và pháp tuyến thay đổi.}
{Từ thông không đổi trong mọi chuyển động của khung dây.}
{Có thể bỏ qua phương, chiều và đơn vị trong mọi trường hợp.}
{Kết quả không phụ thuộc dữ kiện và điều kiện áp dụng.}
\loigiai{
Từ thông biến thiên khi $B$, diện tích mạch hoặc góc giữa $B$ và pháp tuyến thay đổi. Vì vậy phương án $A$ đúng; các phương án còn lại sai.
}
\end{ex}

% ===== Câu 150 =====
% ID: L12C3B16-33-DS
% Mức: VD
\dangbt{Nhận biết từ thông và tính từ thông qua một diện tích}
\begin{ex}
% ID: L12C3B16-33-DS
% Mức: VD
Xét Nhận biết từ thông và tính từ thông qua một diện tích (Tình huống 3). Hãy xác định tính đúng, sai của bốn phát biểu sau.
\choiceTF
{\True Khi giải cần xác định đúng dữ kiện, phương chiều, điều kiện áp dụng và đơn vị.}
{Kết quả không phụ thuộc dữ kiện và có thể bỏ qua điều kiện.}
{\True Từ thông qua diện tích $S$ trong từ trường đều là Phi = BS cos(alpha), với alpha là góc giữa $B$ và pháp tuyến.}
{Từ thông luôn bằng BS sin(alpha), với alpha là góc giữa $B$ và pháp tuyến.}
\loigiai{
\begin{itemchoice}
\item (Đúng) Khi giải cần xác định đúng dữ kiện, phương chiều, điều kiện áp dụng và đơn vị. (Vì): lpha), với alpha là góc giữa $B$ và pháp tuyến
\item (Sai) Kết quả không phụ thuộc dữ kiện và có thể bỏ qua điều kiện.
\item (Đúng) Từ thông qua diện tích $S$ trong từ trường đều là Phi = BS cos(alpha), với alpha là góc giữa $B$ và pháp tuyến.
\item (Sai) Từ thông luôn bằng BS sin(alpha), với alpha là góc giữa $B$ và pháp tuyến.
\end{itemchoice}
}
\end{ex}

% ===== Câu 151 =====
% ID: L12C3B16-33-TL
% Mức: VD
\dangbt{Xác định chiều dòng điện cảm ứng bằng định luật Lenz}
\begin{bt}
% ID: L12C3B16-33-TL
% Mức: VD
Phân tích sai lầm trong nhận định: “Dòng điện cảm ứng luôn làm tăng biến thiên từ thông ban đầu.”
\loigiai{
Cần nêu được: Dòng điện cảm ứng có chiều sao cho từ trường do nó sinh ra chống lại sự biến thiên từ thông đã gây ra nó. Khi vận dụng phải xác định đúng dữ kiện, phương chiều nếu có, điều kiện áp dụng, hệ đơn vị và kiểm tra tính hợp lí. Nhận định “Dòng điện cảm ứng luôn làm tăng biến thiên từ thông ban đầu.” sai.
}
\end{bt}

% ===== Câu 152 =====
% ID: L12C3B16-33-TLN
% Mức: TH
\begin{ex}
% ID: L12C3B16-33-TLN
% Mức: TH
Trong bốn phát biểu về Xác định chiều dòng điện cảm ứng bằng định luật Lenz, có bao nhiêu phát biểu đúng? (Chỉ ghi một số nguyên.) (1) Dòng điện cảm ứng có chiều sao cho từ trường do nó sinh ra chống lại sự biến thiên từ thông đã gây ra nó. (2) Dòng điện cảm ứng luôn làm tăng biến thiên từ thông ban đầu. (3) Cần kiểm tra điều kiện áp dụng. (4) Có thể bỏ qua đơn vị.
\shortans{2}
\loigiai{
Đối chiếu với kiến thức: Dòng điện cảm ứng có chiều sao cho từ trường do nó sinh ra chống lại sự biến thiên từ thông đã gây ra nó. Có 2 phát biểu đúng.
}
\end{ex}

% ===== Câu 153 =====
% ID: L12C3B16-33-TN
% Mức: NB
\dangbt{Phân tích sự biến thiên từ thông qua khung dây}
\begin{ex}
% ID: L12C3B16-33-TN
% Mức: NB
Tình huống 3: chọn kết luận chính xác về Phân tích sự biến thiên từ thông qua khung dây?
\choice
{Kết quả không phụ thuộc dữ kiện và điều kiện áp dụng.}
{\True Từ thông biến thiên khi $B$, diện tích mạch hoặc góc giữa $B$ và pháp tuyến thay đổi.}
{Từ thông không đổi trong mọi chuyển động của khung dây.}
{Có thể bỏ qua phương, chiều và đơn vị trong mọi trường hợp.}
\loigiai{
Từ thông biến thiên khi $B$, diện tích mạch hoặc góc giữa $B$ và pháp tuyến thay đổi. Vì vậy phương án $B$ đúng; các phương án còn lại sai.
}
\end{ex}

% ===== Câu 154 =====
% ID: L12C3B16-34-DS
% Mức: VD
\dangbt{Nhận biết từ thông và tính từ thông qua một diện tích}
\begin{ex}
% ID: L12C3B16-34-DS
% Mức: VD
Xét Nhận biết từ thông và tính từ thông qua một diện tích (Tình huống 4). Hãy xác định tính đúng, sai của bốn phát biểu sau.
\choiceTF
{Kết quả không phụ thuộc dữ kiện và có thể bỏ qua điều kiện.}
{\True Khi giải cần xác định đúng dữ kiện, phương chiều, điều kiện áp dụng và đơn vị.}
{Từ thông luôn bằng BS sin(alpha), với alpha là góc giữa $B$ và pháp tuyến.}
{\True Từ thông qua diện tích $S$ trong từ trường đều là Phi = BS cos(alpha), với alpha là góc giữa $B$ và pháp tuyến.}
\loigiai{
\begin{itemchoice}
\item (Sai) Kết quả không phụ thuộc dữ kiện và có thể bỏ qua điều kiện. (Vì): lpha), với alpha là góc giữa $B$ và pháp tuyến
\item (Đúng) Khi giải cần xác định đúng dữ kiện, phương chiều, điều kiện áp dụng và đơn vị.
\item (Sai) Từ thông luôn bằng BS sin(alpha), với alpha là góc giữa $B$ và pháp tuyến.
\item (Đúng) Từ thông qua diện tích $S$ trong từ trường đều là Phi = BS cos(alpha), với alpha là góc giữa $B$ và pháp tuyến.
\end{itemchoice}
}
\end{ex}

% ===== Câu 155 =====
% ID: L12C3B16-34-TL
% Mức: VD
\dangbt{Xác định chiều dòng điện cảm ứng bằng định luật Lenz}
\begin{bt}
% ID: L12C3B16-34-TL
% Mức: VD
Đề xuất các bước giải một bài toán thuộc dạng Xác định chiều dòng điện cảm ứng bằng định luật Lenz.
\loigiai{
Cần nêu được: Dòng điện cảm ứng có chiều sao cho từ trường do nó sinh ra chống lại sự biến thiên từ thông đã gây ra nó. Khi vận dụng phải xác định đúng dữ kiện, phương chiều nếu có, điều kiện áp dụng, hệ đơn vị và kiểm tra tính hợp lí. Nhận định “Dòng điện cảm ứng luôn làm tăng biến thiên từ thông ban đầu.” sai.
}
\end{bt}

% ===== Câu 156 =====
% ID: L12C3B16-34-TLN
% Mức: VD
\begin{ex}
% ID: L12C3B16-34-TLN
% Mức: VD
Trong bốn phát biểu về Xác định chiều dòng điện cảm ứng bằng định luật Lenz, có bao nhiêu phát biểu đúng? (Chỉ ghi một số nguyên.) (1) Dòng điện cảm ứng luôn làm tăng biến thiên từ thông ban đầu. (2) Dòng điện cảm ứng có chiều sao cho từ trường do nó sinh ra chống lại sự biến thiên từ thông đã gây ra nó. (3) Kết quả phải phù hợp ý nghĩa vật lí. (4) Mọi đại lượng đều là vô hướng.
\shortans{2}
\loigiai{
Đối chiếu với kiến thức: Dòng điện cảm ứng có chiều sao cho từ trường do nó sinh ra chống lại sự biến thiên từ thông đã gây ra nó. Có 2 phát biểu đúng.
}
\end{ex}

% ===== Câu 157 =====
% ID: L12C3B16-34-TN
% Mức: TH
\dangbt{Phân tích sự biến thiên từ thông qua khung dây}
\begin{ex}
% ID: L12C3B16-34-TN
% Mức: TH
Nội dung nào mô tả đúng nhất Phân tích sự biến thiên từ thông qua khung dây?
\choice
{Có thể bỏ qua phương, chiều và đơn vị trong mọi trường hợp.}
{Kết quả không phụ thuộc dữ kiện và điều kiện áp dụng.}
{\True Từ thông biến thiên khi $B$, diện tích mạch hoặc góc giữa $B$ và pháp tuyến thay đổi.}
{Từ thông không đổi trong mọi chuyển động của khung dây.}
\loigiai{
Từ thông biến thiên khi $B$, diện tích mạch hoặc góc giữa $B$ và pháp tuyến thay đổi. Vì vậy phương án $C$ đúng; các phương án còn lại sai.
}
\end{ex}

% ===== Câu 158 =====
% ID: L12C3B16-35-DS
% Mức: VD
\dangbt{Nhận biết từ thông và tính từ thông qua một diện tích}
\begin{ex}
% ID: L12C3B16-35-DS
% Mức: VD
Xét Nhận biết từ thông và tính từ thông qua một diện tích (Tình huống 5). Hãy xác định tính đúng, sai của bốn phát biểu sau.
\choiceTF
{\True Từ thông qua diện tích $S$ trong từ trường đều là Phi = BS cos(alpha), với alpha là góc giữa $B$ và pháp tuyến.}
{Kết quả không phụ thuộc dữ kiện và có thể bỏ qua điều kiện.}
{Từ thông luôn bằng BS sin(alpha), với alpha là góc giữa $B$ và pháp tuyến.}
{\True Khi giải cần xác định đúng dữ kiện, phương chiều, điều kiện áp dụng và đơn vị.}
\loigiai{
\begin{itemchoice}
\item (Đúng) Từ thông qua diện tích $S$ trong từ trường đều là Phi = BS cos(alpha), với alpha là góc giữa $B$ và pháp tuyến. (Vì): lpha), với alpha là góc giữa $B$ và pháp tuyến
\item (Sai) Kết quả không phụ thuộc dữ kiện và có thể bỏ qua điều kiện.
\item (Sai) Từ thông luôn bằng BS sin(alpha), với alpha là góc giữa $B$ và pháp tuyến.
\item (Đúng) Khi giải cần xác định đúng dữ kiện, phương chiều, điều kiện áp dụng và đơn vị.
\end{itemchoice}
}
\end{ex}

% ===== Câu 159 =====
% ID: L12C3B16-35-TL
% Mức: VD
\dangbt{Xác định chiều dòng điện cảm ứng bằng định luật Lenz}
\begin{bt}
% ID: L12C3B16-35-TL
% Mức: VD
Nêu một tình huống thực tiễn liên quan đến Xác định chiều dòng điện cảm ứng bằng định luật Lenz và giải thích bằng kiến thức Vật lí.
\loigiai{
Cần nêu được: Dòng điện cảm ứng có chiều sao cho từ trường do nó sinh ra chống lại sự biến thiên từ thông đã gây ra nó. Khi vận dụng phải xác định đúng dữ kiện, phương chiều nếu có, điều kiện áp dụng, hệ đơn vị và kiểm tra tính hợp lí. Nhận định “Dòng điện cảm ứng luôn làm tăng biến thiên từ thông ban đầu.” sai.
}
\end{bt}

% ===== Câu 160 =====
% ID: L12C3B16-35-TLN
% Mức: VD
\begin{ex}
% ID: L12C3B16-35-TLN
% Mức: VD
Trong bốn phát biểu về Xác định chiều dòng điện cảm ứng bằng định luật Lenz, có bao nhiêu phát biểu đúng? (Chỉ ghi một số nguyên.) (1) Dòng điện cảm ứng có chiều sao cho từ trường do nó sinh ra chống lại sự biến thiên từ thông đã gây ra nó. (2) Dòng điện cảm ứng luôn làm tăng biến thiên từ thông ban đầu. (3) Cần kiểm tra điều kiện áp dụng. (4) Có thể bỏ qua đơn vị.
\shortans{2}
\loigiai{
Đối chiếu với kiến thức: Dòng điện cảm ứng có chiều sao cho từ trường do nó sinh ra chống lại sự biến thiên từ thông đã gây ra nó. Có 2 phát biểu đúng.
}
\end{ex}

% ===== Câu 161 =====
% ID: L12C3B16-35-TN
% Mức: TH
\dangbt{Phân tích sự biến thiên từ thông qua khung dây}
\begin{ex}
% ID: L12C3B16-35-TN
% Mức: TH
Khi phân tích Phân tích sự biến thiên từ thông qua khung dây?
\choice
{Từ thông không đổi trong mọi chuyển động của khung dây.}
{Có thể bỏ qua phương, chiều và đơn vị trong mọi trường hợp.}
{Kết quả không phụ thuộc dữ kiện và điều kiện áp dụng.}
{\True Từ thông biến thiên khi $B$, diện tích mạch hoặc góc giữa $B$ và pháp tuyến thay đổi.}
\loigiai{
Từ thông biến thiên khi $B$, diện tích mạch hoặc góc giữa $B$ và pháp tuyến thay đổi. Vì vậy phương án $D$ đúng; các phương án còn lại sai.
}
\end{ex}

% ===== Câu 162 =====
% ID: L12C3B16-36-DS
% Mức: TH
\begin{ex}
% ID: L12C3B16-36-DS
% Mức: TH
Xét Phân tích sự biến thiên từ thông qua khung dây (Tình huống 1). Hãy xác định tính đúng, sai của bốn phát biểu sau.
\choiceTF
{\True Từ thông biến thiên khi $B$, diện tích mạch hoặc góc giữa $B$ và pháp tuyến thay đổi.}
{Từ thông không đổi trong mọi chuyển động của khung dây.}
{\True Khi giải cần xác định đúng dữ kiện, phương chiều, điều kiện áp dụng và đơn vị.}
{Kết quả không phụ thuộc dữ kiện và có thể bỏ qua điều kiện.}
\loigiai{
\begin{itemchoice}
\item (Đúng) Từ thông biến thiên khi $B$, diện tích mạch hoặc góc giữa $B$ và pháp tuyến thay đổi. (Vì): Từ thông biến thiên khi $B$, diện tích mạch hoặc góc giữa $B$ và pháp tuyến thay đổi. b) Sai: Từ thông không đổi trong mọi chuyển động của khung dây. c) Đúng: Khi giải cần xác định đúng dữ kiện, phương chiều, điều kiện áp dụng và đơn vị. d) Sai: Kết quả không phụ thuộc dữ kiện và có thể bỏ qua điều kiện. Kiến thức cốt lõi: Từ thông biến thiên khi $B$, diện tích mạch hoặc góc giữa $B$ và pháp tuyến thay đổi
\item (Sai) Từ thông không đổi trong mọi chuyển động của khung dây.
\item (Đúng) Khi giải cần xác định đúng dữ kiện, phương chiều, điều kiện áp dụng và đơn vị.
\item (Sai) Kết quả không phụ thuộc dữ kiện và có thể bỏ qua điều kiện.
\end{itemchoice}
}
\end{ex}

% ===== Câu 163 =====
% ID: L12C3B16-36-TL
% Mức: VDC
\dangbt{Xác định chiều dòng điện cảm ứng bằng định luật Lenz}
\begin{bt}
% ID: L12C3B16-36-TL
% Mức: VDC
So sánh nhận định đúng “Dòng điện cảm ứng có chiều sao cho từ trường do nó sinh ra chống lại sự biến thiên từ thông đã gây ra nó.” với nhận định sai “Dòng điện cảm ứng luôn làm tăng biến thiên từ thông ban đầu.”; chỉ rõ căn cứ Vật lí.
\loigiai{
Cần nêu được: Dòng điện cảm ứng có chiều sao cho từ trường do nó sinh ra chống lại sự biến thiên từ thông đã gây ra nó. Khi vận dụng phải xác định đúng dữ kiện, phương chiều nếu có, điều kiện áp dụng, hệ đơn vị và kiểm tra tính hợp lí. Nhận định “Dòng điện cảm ứng luôn làm tăng biến thiên từ thông ban đầu.” sai.
}
\end{bt}

% ===== Câu 164 =====
% ID: L12C3B16-36-TLN
% Mức: VDC
\begin{ex}
% ID: L12C3B16-36-TLN
% Mức: VDC
Trong bốn phát biểu về Xác định chiều dòng điện cảm ứng bằng định luật Lenz, có bao nhiêu phát biểu đúng? (Chỉ ghi một số nguyên.) (1) Khi giải cần xác định đúng dữ kiện, phương chiều, điều kiện áp dụng và đơn vị. (2) Dòng điện cảm ứng luôn làm tăng biến thiên từ thông ban đầu. (3) Phải dùng đúng định luật cho tình huống. (4) Dòng điện cảm ứng có chiều sao cho từ trường do nó sinh ra chống lại sự biến thiên từ thông đã gây ra nó.
\shortans{3}
\loigiai{
Đối chiếu với kiến thức: Dòng điện cảm ứng có chiều sao cho từ trường do nó sinh ra chống lại sự biến thiên từ thông đã gây ra nó. Có 3 phát biểu đúng.
}
\end{ex}

% ===== Câu 165 =====
% ID: L12C3B16-36-TN
% Mức: TH
\dangbt{Phân tích sự biến thiên từ thông qua khung dây}
\begin{ex}
% ID: L12C3B16-36-TN
% Mức: TH
Một học sinh ôn tập Phân tích sự biến thiên từ thông qua khung dây?
\choice
{\True Từ thông biến thiên khi $B$, diện tích mạch hoặc góc giữa $B$ và pháp tuyến thay đổi.}
{Từ thông không đổi trong mọi chuyển động của khung dây.}
{Có thể bỏ qua phương, chiều và đơn vị trong mọi trường hợp.}
{Kết quả không phụ thuộc dữ kiện và điều kiện áp dụng.}
\loigiai{
Từ thông biến thiên khi $B$, diện tích mạch hoặc góc giữa $B$ và pháp tuyến thay đổi. Vì vậy phương án $A$ đúng; các phương án còn lại sai.
}
\end{ex}

% ===== Câu 166 =====
% ID: L12C3B16-37-DS
% Mức: TH
\begin{ex}
% ID: L12C3B16-37-DS
% Mức: TH
Xét Phân tích sự biến thiên từ thông qua khung dây (Tình huống 2). Hãy xác định tính đúng, sai của bốn phát biểu sau.
\choiceTF
{Từ thông không đổi trong mọi chuyển động của khung dây.}
{\True Từ thông biến thiên khi $B$, diện tích mạch hoặc góc giữa $B$ và pháp tuyến thay đổi.}
{Kết quả không phụ thuộc dữ kiện và có thể bỏ qua điều kiện.}
{\True Khi giải cần xác định đúng dữ kiện, phương chiều, điều kiện áp dụng và đơn vị.}
\loigiai{
\begin{itemchoice}
\item (Sai) Từ thông không đổi trong mọi chuyển động của khung dây. (Vì): Từ thông không đổi trong mọi chuyển động của khung dây. b) Đúng: Từ thông biến thiên khi $B$, diện tích mạch hoặc góc giữa $B$ và pháp tuyến thay đổi. c) Sai: Kết quả không phụ thuộc dữ kiện và có thể bỏ qua điều kiện. d) Đúng: Khi giải cần xác định đúng dữ kiện, phương chiều, điều kiện áp dụng và đơn vị. Kiến thức cốt lõi: Từ thông biến thiên khi $B$, diện tích mạch hoặc góc giữa $B$ và pháp tuyến thay đổi
\item (Đúng) Từ thông biến thiên khi $B$, diện tích mạch hoặc góc giữa $B$ và pháp tuyến thay đổi.
\item (Sai) Kết quả không phụ thuộc dữ kiện và có thể bỏ qua điều kiện.
\item (Đúng) Khi giải cần xác định đúng dữ kiện, phương chiều, điều kiện áp dụng và đơn vị.
\end{itemchoice}
}
\end{ex}

% ===== Câu 167 =====
% ID: L12C3B16-37-TL
% Mức: TH
\dangbt{Nhận biết từ thông và tính từ thông qua một diện tích}
\begin{bt}
% ID: L12C3B16-37-TL
% Mức: TH
Trình bày kiến thức cốt lõi của Nhận biết từ thông và tính từ thông qua một diện tích và nêu điều kiện áp dụng.
\loigiai{
Cần nêu được: Từ thông qua diện tích $S$ trong từ trường đều là Phi = BS cos(alpha), với alpha là góc giữa $B$ và pháp tuyến. Khi vận dụng phải xác định đúng dữ kiện, phương chiều nếu có, điều kiện áp dụng, hệ đơn vị và kiểm tra tính hợp lí. Nhận định “Từ thông luôn bằng BS sin(alpha), với alpha là góc giữa $B$ và pháp tuyến.” sai.
}
\end{bt}

% ===== Câu 168 =====
% ID: L12C3B16-37-TLN
% Mức: TH
\begin{ex}
% ID: L12C3B16-37-TLN
% Mức: TH
Trong bốn phát biểu về Nhận biết từ thông và tính từ thông qua một diện tích, có bao nhiêu phát biểu đúng? (Chỉ ghi một số nguyên.) (1) Từ thông qua diện tích $S$ trong từ trường đều là Phi = BS cos(alpha), với alpha là góc giữa $B$ và pháp tuyến. (2) Từ thông luôn bằng BS sin(alpha), với alpha là góc giữa $B$ và pháp tuyến. (3) Cần kiểm tra điều kiện áp dụng. (4) Có thể bỏ qua đơn vị.
\shortans{2}
\loigiai{
Đối chiếu với kiến thức: Từ thông qua diện tích $S$ trong từ trường đều là Phi = BS cos(alpha), với alpha là góc giữa $B$ và pháp tuyến. Có 2 phát biểu đúng.
}
\end{ex}

% ===== Câu 169 =====
% ID: L12C3B16-37-TN
% Mức: TH
\dangbt{Phân tích sự biến thiên từ thông qua khung dây}
\begin{ex}
% ID: L12C3B16-37-TN
% Mức: TH
Chọn phát biểu không trái với kiến thức về Phân tích sự biến thiên từ thông qua khung dây?
\choice
{Kết quả không phụ thuộc dữ kiện và điều kiện áp dụng.}
{\True Từ thông biến thiên khi $B$, diện tích mạch hoặc góc giữa $B$ và pháp tuyến thay đổi.}
{Từ thông không đổi trong mọi chuyển động của khung dây.}
{Có thể bỏ qua phương, chiều và đơn vị trong mọi trường hợp.}
\loigiai{
Từ thông biến thiên khi $B$, diện tích mạch hoặc góc giữa $B$ và pháp tuyến thay đổi. Vì vậy phương án $B$ đúng; các phương án còn lại sai.
}
\end{ex}

% ===== Câu 170 =====
% ID: L12C3B16-38-DS
% Mức: VD
\begin{ex}
% ID: L12C3B16-38-DS
% Mức: VD
Xét Phân tích sự biến thiên từ thông qua khung dây (Tình huống 3). Hãy xác định tính đúng, sai của bốn phát biểu sau.
\choiceTF
{\True Khi giải cần xác định đúng dữ kiện, phương chiều, điều kiện áp dụng và đơn vị.}
{Kết quả không phụ thuộc dữ kiện và có thể bỏ qua điều kiện.}
{\True Từ thông biến thiên khi $B$, diện tích mạch hoặc góc giữa $B$ và pháp tuyến thay đổi.}
{Từ thông không đổi trong mọi chuyển động của khung dây.}
\loigiai{
\begin{itemchoice}
\item (Đúng) Khi giải cần xác định đúng dữ kiện, phương chiều, điều kiện áp dụng và đơn vị. (Vì): Khi giải cần xác định đúng dữ kiện, phương chiều, điều kiện áp dụng và đơn vị. b) Sai: Kết quả không phụ thuộc dữ kiện và có thể bỏ qua điều kiện. c) Đúng: Từ thông biến thiên khi $B$, diện tích mạch hoặc góc giữa $B$ và pháp tuyến thay đổi. d) Sai: Từ thông không đổi trong mọi chuyển động của khung dây. Kiến thức cốt lõi: Từ thông biến thiên khi $B$, diện tích mạch hoặc góc giữa $B$ và pháp tuyến thay đổi
\item (Sai) Kết quả không phụ thuộc dữ kiện và có thể bỏ qua điều kiện.
\item (Đúng) Từ thông biến thiên khi $B$, diện tích mạch hoặc góc giữa $B$ và pháp tuyến thay đổi.
\item (Sai) Từ thông không đổi trong mọi chuyển động của khung dây.
\end{itemchoice}
}
\end{ex}

% ===== Câu 171 =====
% ID: L12C3B16-38-TL
% Mức: TH
\dangbt{Nhận biết từ thông và tính từ thông qua một diện tích}
\begin{bt}
% ID: L12C3B16-38-TL
% Mức: TH
Giải thích vì sao nhận định sau đúng: “Từ thông qua diện tích $S$ trong từ trường đều là Phi = BS cos(alpha), với alpha là góc giữa $B$ và pháp tuyến.”
\loigiai{
Cần nêu được: Từ thông qua diện tích $S$ trong từ trường đều là Phi = BS cos(alpha), với alpha là góc giữa $B$ và pháp tuyến. Khi vận dụng phải xác định đúng dữ kiện, phương chiều nếu có, điều kiện áp dụng, hệ đơn vị và kiểm tra tính hợp lí. Nhận định “Từ thông luôn bằng BS sin(alpha), với alpha là góc giữa $B$ và pháp tuyến.” sai.
}
\end{bt}

% ===== Câu 172 =====
% ID: L12C3B16-38-TLN
% Mức: TH
\begin{ex}
% ID: L12C3B16-38-TLN
% Mức: TH
Trong bốn phát biểu về Nhận biết từ thông và tính từ thông qua một diện tích, có bao nhiêu phát biểu đúng? (Chỉ ghi một số nguyên.) (1) Từ thông luôn bằng BS sin(alpha), với alpha là góc giữa $B$ và pháp tuyến. (2) Từ thông qua diện tích $S$ trong từ trường đều là Phi = BS cos(alpha), với alpha là góc giữa $B$ và pháp tuyến. (3) Kết quả phải phù hợp ý nghĩa vật lí. (4) Mọi đại lượng đều là vô hướng.
\shortans{1}
\loigiai{
Đối chiếu với kiến thức: Từ thông qua diện tích $S$ trong từ trường đều là Phi = BS cos(alpha), với alpha là góc giữa $B$ và pháp tuyến. Có 1 phát biểu đúng.
}
\end{ex}

% ===== Câu 173 =====
% ID: L12C3B16-38-TN
% Mức: VD
\dangbt{Phân tích sự biến thiên từ thông qua khung dây}
\begin{ex}
% ID: L12C3B16-38-TN
% Mức: VD
Cơ sở Vật lí đúng dùng để giải Phân tích sự biến thiên từ thông qua khung dây?
\choice
{Có thể bỏ qua phương, chiều và đơn vị trong mọi trường hợp.}
{Kết quả không phụ thuộc dữ kiện và điều kiện áp dụng.}
{\True Từ thông biến thiên khi $B$, diện tích mạch hoặc góc giữa $B$ và pháp tuyến thay đổi.}
{Từ thông không đổi trong mọi chuyển động của khung dây.}
\loigiai{
Từ thông biến thiên khi $B$, diện tích mạch hoặc góc giữa $B$ và pháp tuyến thay đổi. Vì vậy phương án $C$ đúng; các phương án còn lại sai.
}
\end{ex}

% ===== Câu 174 =====
% ID: L12C3B16-39-DS
% Mức: VD
\begin{ex}
% ID: L12C3B16-39-DS
% Mức: VD
Xét Phân tích sự biến thiên từ thông qua khung dây (Tình huống 4). Hãy xác định tính đúng, sai của bốn phát biểu sau.
\choiceTF
{Kết quả không phụ thuộc dữ kiện và có thể bỏ qua điều kiện.}
{\True Khi giải cần xác định đúng dữ kiện, phương chiều, điều kiện áp dụng và đơn vị.}
{Từ thông không đổi trong mọi chuyển động của khung dây.}
{\True Từ thông biến thiên khi $B$, diện tích mạch hoặc góc giữa $B$ và pháp tuyến thay đổi.}
\loigiai{
\begin{itemchoice}
\item (Sai) Kết quả không phụ thuộc dữ kiện và có thể bỏ qua điều kiện. (Vì): Kết quả không phụ thuộc dữ kiện và có thể bỏ qua điều kiện. b) Đúng: Khi giải cần xác định đúng dữ kiện, phương chiều, điều kiện áp dụng và đơn vị. c) Sai: Từ thông không đổi trong mọi chuyển động của khung dây. d) Đúng: Từ thông biến thiên khi $B$, diện tích mạch hoặc góc giữa $B$ và pháp tuyến thay đổi. Kiến thức cốt lõi: Từ thông biến thiên khi $B$, diện tích mạch hoặc góc giữa $B$ và pháp tuyến thay đổi
\item (Đúng) Khi giải cần xác định đúng dữ kiện, phương chiều, điều kiện áp dụng và đơn vị.
\item (Sai) Từ thông không đổi trong mọi chuyển động của khung dây.
\item (Đúng) Từ thông biến thiên khi $B$, diện tích mạch hoặc góc giữa $B$ và pháp tuyến thay đổi.
\end{itemchoice}
}
\end{ex}

% ===== Câu 175 =====
% ID: L12C3B16-39-TL
% Mức: VD
\dangbt{Nhận biết từ thông và tính từ thông qua một diện tích}
\begin{bt}
% ID: L12C3B16-39-TL
% Mức: VD
Phân tích sai lầm trong nhận định: “Từ thông luôn bằng BS sin(alpha), với alpha là góc giữa $B$ và pháp tuyến.”
\loigiai{
Cần nêu được: Từ thông qua diện tích $S$ trong từ trường đều là Phi = BS cos(alpha), với alpha là góc giữa $B$ và pháp tuyến. Khi vận dụng phải xác định đúng dữ kiện, phương chiều nếu có, điều kiện áp dụng, hệ đơn vị và kiểm tra tính hợp lí. Nhận định “Từ thông luôn bằng BS sin(alpha), với alpha là góc giữa $B$ và pháp tuyến.” sai.
}
\end{bt}

% ===== Câu 176 =====
% ID: L12C3B16-39-TLN
% Mức: TH
\begin{ex}
% ID: L12C3B16-39-TLN
% Mức: TH
Trong bốn phát biểu về Nhận biết từ thông và tính từ thông qua một diện tích, có bao nhiêu phát biểu đúng? (Chỉ ghi một số nguyên.) (1) Từ thông qua diện tích $S$ trong từ trường đều là Phi = BS cos(alpha), với alpha là góc giữa $B$ và pháp tuyến. (2) Từ thông luôn bằng BS sin(alpha), với alpha là góc giữa $B$ và pháp tuyến. (3) Cần kiểm tra điều kiện áp dụng. (4) Có thể bỏ qua đơn vị.
\shortans{2}
\loigiai{
Đối chiếu với kiến thức: Từ thông qua diện tích $S$ trong từ trường đều là Phi = BS cos(alpha), với alpha là góc giữa $B$ và pháp tuyến. Có 2 phát biểu đúng.
}
\end{ex}

% ===== Câu 177 =====
% ID: L12C3B16-39-TN
% Mức: VD
\dangbt{Phân tích sự biến thiên từ thông qua khung dây}
\begin{ex}
% ID: L12C3B16-39-TN
% Mức: VD
Trong một bài thực hành về Phân tích sự biến thiên từ thông qua khung dây?
\choice
{Từ thông không đổi trong mọi chuyển động của khung dây.}
{Có thể bỏ qua phương, chiều và đơn vị trong mọi trường hợp.}
{Kết quả không phụ thuộc dữ kiện và điều kiện áp dụng.}
{\True Từ thông biến thiên khi $B$, diện tích mạch hoặc góc giữa $B$ và pháp tuyến thay đổi.}
\loigiai{
Từ thông biến thiên khi $B$, diện tích mạch hoặc góc giữa $B$ và pháp tuyến thay đổi. Vì vậy phương án $D$ đúng; các phương án còn lại sai.
}
\end{ex}

% ===== Câu 178 =====
% ID: L12C3B16-40-DS
% Mức: VD
\begin{ex}
% ID: L12C3B16-40-DS
% Mức: VD
Xét Phân tích sự biến thiên từ thông qua khung dây (Tình huống 5). Hãy xác định tính đúng, sai của bốn phát biểu sau.
\choiceTF
{\True Từ thông biến thiên khi $B$, diện tích mạch hoặc góc giữa $B$ và pháp tuyến thay đổi.}
{Kết quả không phụ thuộc dữ kiện và có thể bỏ qua điều kiện.}
{Từ thông không đổi trong mọi chuyển động của khung dây.}
{\True Khi giải cần xác định đúng dữ kiện, phương chiều, điều kiện áp dụng và đơn vị.}
\loigiai{
\begin{itemchoice}
\item (Đúng) Từ thông biến thiên khi $B$, diện tích mạch hoặc góc giữa $B$ và pháp tuyến thay đổi. (Vì): Từ thông biến thiên khi $B$, diện tích mạch hoặc góc giữa $B$ và pháp tuyến thay đổi. b) Sai: Kết quả không phụ thuộc dữ kiện và có thể bỏ qua điều kiện. c) Sai: Từ thông không đổi trong mọi chuyển động của khung dây. d) Đúng: Khi giải cần xác định đúng dữ kiện, phương chiều, điều kiện áp dụng và đơn vị. Kiến thức cốt lõi: Từ thông biến thiên khi $B$, diện tích mạch hoặc góc giữa $B$ và pháp tuyến thay đổi
\item (Sai) Kết quả không phụ thuộc dữ kiện và có thể bỏ qua điều kiện.
\item (Sai) Từ thông không đổi trong mọi chuyển động của khung dây.
\item (Đúng) Khi giải cần xác định đúng dữ kiện, phương chiều, điều kiện áp dụng và đơn vị.
\end{itemchoice}
}
\end{ex}

% ===== Câu 179 =====
% ID: L12C3B16-40-TL
% Mức: VD
\dangbt{Nhận biết từ thông và tính từ thông qua một diện tích}
\begin{bt}
% ID: L12C3B16-40-TL
% Mức: VD
Đề xuất các bước giải một bài toán thuộc dạng Nhận biết từ thông và tính từ thông qua một diện tích.
\loigiai{
Cần nêu được: Từ thông qua diện tích $S$ trong từ trường đều là Phi = BS cos(alpha), với alpha là góc giữa $B$ và pháp tuyến. Khi vận dụng phải xác định đúng dữ kiện, phương chiều nếu có, điều kiện áp dụng, hệ đơn vị và kiểm tra tính hợp lí. Nhận định “Từ thông luôn bằng BS sin(alpha), với alpha là góc giữa $B$ và pháp tuyến.” sai.
}
\end{bt}

% ===== Câu 180 =====
% ID: L12C3B16-40-TLN
% Mức: VD
\begin{ex}
% ID: L12C3B16-40-TLN
% Mức: VD
Trong bốn phát biểu về Nhận biết từ thông và tính từ thông qua một diện tích, có bao nhiêu phát biểu đúng? (Chỉ ghi một số nguyên.) (1) Từ thông luôn bằng BS sin(alpha), với alpha là góc giữa $B$ và pháp tuyến. (2) Từ thông qua diện tích $S$ trong từ trường đều là Phi = BS cos(alpha), với alpha là góc giữa $B$ và pháp tuyến. (3) Kết quả phải phù hợp ý nghĩa vật lí. (4) Mọi đại lượng đều là vô hướng.
\shortans{2}
\loigiai{
Đối chiếu với kiến thức: Từ thông qua diện tích $S$ trong từ trường đều là Phi = BS cos(alpha), với alpha là góc giữa $B$ và pháp tuyến. Có 2 phát biểu đúng.
}
\end{ex}

% ===== Câu 181 =====
% ID: L12C3B16-40-TN
% Mức: VD
\dangbt{Phân tích sự biến thiên từ thông qua khung dây}
\begin{ex}
% ID: L12C3B16-40-TN
% Mức: VD
Kết luận đúng khi vận dụng Phân tích sự biến thiên từ thông qua khung dây?
\choice
{\True Từ thông biến thiên khi $B$, diện tích mạch hoặc góc giữa $B$ và pháp tuyến thay đổi.}
{Từ thông không đổi trong mọi chuyển động của khung dây.}
{Có thể bỏ qua phương, chiều và đơn vị trong mọi trường hợp.}
{Kết quả không phụ thuộc dữ kiện và điều kiện áp dụng.}
\loigiai{
Từ thông biến thiên khi $B$, diện tích mạch hoặc góc giữa $B$ và pháp tuyến thay đổi. Vì vậy phương án $A$ đúng; các phương án còn lại sai.
}
\end{ex}

% ===== Câu 182 =====
% ID: L12C3B16-41-DS
% Mức: TH
\dangbt{Nhận biết hiện tượng cảm ứng điện từ và điều kiện xuất hiện dòng điện cảm ứng}
\begin{ex}
% ID: L12C3B16-41-DS
% Mức: TH
Xét Nhận biết hiện tượng cảm ứng điện từ và điều kiện xuất hiện dòng điện cảm ứng (Tình huống 1). Hãy xác định tính đúng, sai của bốn phát biểu sau.
\choiceTF
{\True Dòng điện cảm ứng xuất hiện trong mạch kín khi từ thông qua mạch biến thiên.}
{Chỉ cần đặt mạch kín đứng yên trong từ trường đều không đổi là luôn có dòng điện cảm ứng.}
{\True Khi giải cần xác định đúng dữ kiện, phương chiều, điều kiện áp dụng và đơn vị.}
{Kết quả không phụ thuộc dữ kiện và có thể bỏ qua điều kiện.}
\loigiai{
\begin{itemchoice}
\item (Đúng) Dòng điện cảm ứng xuất hiện trong mạch kín khi từ thông qua mạch biến thiên. (Vì): òng điện cảm ứng xuất hiện trong mạch kín khi từ thông qua mạch biến thiên. b) Sai: Chỉ cần đặt mạch kín đứng yên trong từ trường đều không đổi là luôn có dòng điện cảm ứng. c) Đúng: Khi giải cần xác định đúng dữ kiện, phương chiều, điều kiện áp dụng và đơn vị. d) Sai: Kết quả không phụ thuộc dữ kiện và có thể bỏ qua điều kiện. Kiến thức cốt lõi: Dòng điện cảm ứng xuất hiện trong mạch kín khi từ thông qua mạch biến thiên
\item (Sai) Chỉ cần đặt mạch kín đứng yên trong từ trường đều không đổi là luôn có dòng điện cảm ứng.
\item (Đúng) Khi giải cần xác định đúng dữ kiện, phương chiều, điều kiện áp dụng và đơn vị.
\item (Sai) Kết quả không phụ thuộc dữ kiện và có thể bỏ qua điều kiện.
\end{itemchoice}
}
\end{ex}

% ===== Câu 183 =====
% ID: L12C3B16-41-TL
% Mức: VD
\dangbt{Nhận biết từ thông và tính từ thông qua một diện tích}
\begin{bt}
% ID: L12C3B16-41-TL
% Mức: VD
Nêu một tình huống thực tiễn liên quan đến Nhận biết từ thông và tính từ thông qua một diện tích và giải thích bằng kiến thức Vật lí.
\loigiai{
Cần nêu được: Từ thông qua diện tích $S$ trong từ trường đều là Phi = BS cos(alpha), với alpha là góc giữa $B$ và pháp tuyến. Khi vận dụng phải xác định đúng dữ kiện, phương chiều nếu có, điều kiện áp dụng, hệ đơn vị và kiểm tra tính hợp lí. Nhận định “Từ thông luôn bằng BS sin(alpha), với alpha là góc giữa $B$ và pháp tuyến.” sai.
}
\end{bt}

% ===== Câu 184 =====
% ID: L12C3B16-41-TLN
% Mức: VD
\begin{ex}
% ID: L12C3B16-41-TLN
% Mức: VD
Trong bốn phát biểu về Nhận biết từ thông và tính từ thông qua một diện tích, có bao nhiêu phát biểu đúng? (Chỉ ghi một số nguyên.) (1) Từ thông qua diện tích $S$ trong từ trường đều là Phi = BS cos(alpha), với alpha là góc giữa $B$ và pháp tuyến. (2) Từ thông luôn bằng BS sin(alpha), với alpha là góc giữa $B$ và pháp tuyến. (3) Cần kiểm tra điều kiện áp dụng. (4) Có thể bỏ qua đơn vị.
\shortans{2}
\loigiai{
Đối chiếu với kiến thức: Từ thông qua diện tích $S$ trong từ trường đều là Phi = BS cos(alpha), với alpha là góc giữa $B$ và pháp tuyến. Có 2 phát biểu đúng.
}
\end{ex}

% ===== Câu 185 =====
% ID: L12C3B16-41-TN
% Mức: NB
\dangbt{Tính suất điện động cảm ứng theo định luật Faraday}
\begin{ex}
% ID: L12C3B16-41-TN
% Mức: NB
Phát biểu nào sau đây đúng về Tính suất điện động cảm ứng theo định luật Faraday?
\choice
{Suất điện động cảm ứng bằng $N$ Delta t/Delta Phi.}
{Có thể bỏ qua phương, chiều và đơn vị trong mọi trường hợp.}
{Kết quả không phụ thuộc dữ kiện và điều kiện áp dụng.}
{\True Độ lớn suất điện động cảm ứng trung bình của cuộn $N$ vòng là |e| = N|Delta Phi|/Delta t.}
\loigiai{
Độ lớn suất điện động cảm ứng trung bình của cuộn $N$ vòng là |e| = N|Delta Phi|/Delta t. Vì vậy phương án $D$ đúng; các phương án còn lại sai.
}
\end{ex}

% ===== Câu 186 =====
% ID: L12C3B16-42-DS
% Mức: TH
\dangbt{Nhận biết hiện tượng cảm ứng điện từ và điều kiện xuất hiện dòng điện cảm ứng}
\begin{ex}
% ID: L12C3B16-42-DS
% Mức: TH
Xét Nhận biết hiện tượng cảm ứng điện từ và điều kiện xuất hiện dòng điện cảm ứng (Tình huống 2). Hãy xác định tính đúng, sai của bốn phát biểu sau.
\choiceTF
{Chỉ cần đặt mạch kín đứng yên trong từ trường đều không đổi là luôn có dòng điện cảm ứng.}
{\True Dòng điện cảm ứng xuất hiện trong mạch kín khi từ thông qua mạch biến thiên.}
{Kết quả không phụ thuộc dữ kiện và có thể bỏ qua điều kiện.}
{\True Khi giải cần xác định đúng dữ kiện, phương chiều, điều kiện áp dụng và đơn vị.}
\loigiai{
\begin{itemchoice}
\item (Sai) Chỉ cần đặt mạch kín đứng yên trong từ trường đều không đổi là luôn có dòng điện cảm ứng. (Vì): Chỉ cần đặt mạch kín đứng yên trong từ trường đều không đổi là luôn có dòng điện cảm ứng. b) Đúng: Dòng điện cảm ứng xuất hiện trong mạch kín khi từ thông qua mạch biến thiên. c) Sai: Kết quả không phụ thuộc dữ kiện và có thể bỏ qua điều kiện. d) Đúng: Khi giải cần xác định đúng dữ kiện, phương chiều, điều kiện áp dụng và đơn vị. Kiến thức cốt lõi: Dòng điện cảm ứng xuất hiện trong mạch kín khi từ thông qua mạch biến thiên
\item (Đúng) Dòng điện cảm ứng xuất hiện trong mạch kín khi từ thông qua mạch biến thiên.
\item (Sai) Kết quả không phụ thuộc dữ kiện và có thể bỏ qua điều kiện.
\item (Đúng) Khi giải cần xác định đúng dữ kiện, phương chiều, điều kiện áp dụng và đơn vị.
\end{itemchoice}
}
\end{ex}

% ===== Câu 187 =====
% ID: L12C3B16-42-TL
% Mức: VDC
\dangbt{Nhận biết từ thông và tính từ thông qua một diện tích}
\begin{bt}
% ID: L12C3B16-42-TL
% Mức: VDC
So sánh nhận định đúng “Từ thông qua diện tích $S$ trong từ trường đều là Phi = BS cos(alpha), với alpha là góc giữa $B$ và pháp tuyến.” với nhận định sai “Từ thông luôn bằng BS sin(alpha), với alpha là góc giữa $B$ và pháp tuyến.”; chỉ rõ căn cứ Vật lí.
\loigiai{
Cần nêu được: Từ thông qua diện tích $S$ trong từ trường đều là Phi = BS cos(alpha), với alpha là góc giữa $B$ và pháp tuyến. Khi vận dụng phải xác định đúng dữ kiện, phương chiều nếu có, điều kiện áp dụng, hệ đơn vị và kiểm tra tính hợp lí. Nhận định “Từ thông luôn bằng BS sin(alpha), với alpha là góc giữa $B$ và pháp tuyến.” sai.
}
\end{bt}

% ===== Câu 188 =====
% ID: L12C3B16-42-TLN
% Mức: VDC
\begin{ex}
% ID: L12C3B16-42-TLN
% Mức: VDC
Trong bốn phát biểu về Nhận biết từ thông và tính từ thông qua một diện tích, có bao nhiêu phát biểu đúng? (Chỉ ghi một số nguyên.) (1) Khi giải cần xác định đúng dữ kiện, phương chiều, điều kiện áp dụng và đơn vị. (2) Từ thông luôn bằng BS sin(alpha), với alpha là góc giữa $B$ và pháp tuyến. (3) Phải dùng đúng định luật cho tình huống. (4) Từ thông qua diện tích $S$ trong từ trường đều là Phi = BS cos(alpha), với alpha là góc giữa $B$ và pháp tuyến.
\shortans{3}
\loigiai{
Đối chiếu với kiến thức: Từ thông qua diện tích $S$ trong từ trường đều là Phi = BS cos(alpha), với alpha là góc giữa $B$ và pháp tuyến. Có 3 phát biểu đúng.
}
\end{ex}

% ===== Câu 189 =====
% ID: L12C3B16-42-TN
% Mức: NB
\dangbt{Tính suất điện động cảm ứng theo định luật Faraday}
\begin{ex}
% ID: L12C3B16-42-TN
% Mức: NB
Trong nội dung Tính suất điện động cảm ứng theo định luật Faraday?
\choice
{\True Độ lớn suất điện động cảm ứng trung bình của cuộn $N$ vòng là |e| = N|Delta Phi|/Delta t.}
{Suất điện động cảm ứng bằng $N$ Delta t/Delta Phi.}
{Có thể bỏ qua phương, chiều và đơn vị trong mọi trường hợp.}
{Kết quả không phụ thuộc dữ kiện và điều kiện áp dụng.}
\loigiai{
Độ lớn suất điện động cảm ứng trung bình của cuộn $N$ vòng là |e| = N|Delta Phi|/Delta t. Vì vậy phương án $A$ đúng; các phương án còn lại sai.
}
\end{ex}

% ===== Câu 190 =====
% ID: L12C3B16-43-DS
% Mức: VD
\dangbt{Nhận biết hiện tượng cảm ứng điện từ và điều kiện xuất hiện dòng điện cảm ứng}
\begin{ex}
% ID: L12C3B16-43-DS
% Mức: VD
Xét Nhận biết hiện tượng cảm ứng điện từ và điều kiện xuất hiện dòng điện cảm ứng (Tình huống 3). Hãy xác định tính đúng, sai của bốn phát biểu sau.
\choiceTF
{\True Khi giải cần xác định đúng dữ kiện, phương chiều, điều kiện áp dụng và đơn vị.}
{Kết quả không phụ thuộc dữ kiện và có thể bỏ qua điều kiện.}
{\True Dòng điện cảm ứng xuất hiện trong mạch kín khi từ thông qua mạch biến thiên.}
{Chỉ cần đặt mạch kín đứng yên trong từ trường đều không đổi là luôn có dòng điện cảm ứng.}
\loigiai{
\begin{itemchoice}
\item (Đúng) Khi giải cần xác định đúng dữ kiện, phương chiều, điều kiện áp dụng và đơn vị. (Vì): Khi giải cần xác định đúng dữ kiện, phương chiều, điều kiện áp dụng và đơn vị. b) Sai: Kết quả không phụ thuộc dữ kiện và có thể bỏ qua điều kiện. c) Đúng: Dòng điện cảm ứng xuất hiện trong mạch kín khi từ thông qua mạch biến thiên. d) Sai: Chỉ cần đặt mạch kín đứng yên trong từ trường đều không đổi là luôn có dòng điện cảm ứng. Kiến thức cốt lõi: Dòng điện cảm ứng xuất hiện trong mạch kín khi từ thông qua mạch biến thiên
\item (Sai) Kết quả không phụ thuộc dữ kiện và có thể bỏ qua điều kiện.
\item (Đúng) Dòng điện cảm ứng xuất hiện trong mạch kín khi từ thông qua mạch biến thiên.
\item (Sai) Chỉ cần đặt mạch kín đứng yên trong từ trường đều không đổi là luôn có dòng điện cảm ứng.
\end{itemchoice}
}
\end{ex}

% ===== Câu 191 =====
% ID: L12C3B16-43-TL
% Mức: TH
\dangbt{Phân tích sự biến thiên từ thông qua khung dây}
\begin{bt}
% ID: L12C3B16-43-TL
% Mức: TH
Trình bày kiến thức cốt lõi của Phân tích sự biến thiên từ thông qua khung dây và nêu điều kiện áp dụng.
\loigiai{
Cần nêu được: Từ thông biến thiên khi $B$, diện tích mạch hoặc góc giữa $B$ và pháp tuyến thay đổi. Khi vận dụng phải xác định đúng dữ kiện, phương chiều nếu có, điều kiện áp dụng, hệ đơn vị và kiểm tra tính hợp lí. Nhận định “Từ thông không đổi trong mọi chuyển động của khung dây.” sai.
}
\end{bt}

% ===== Câu 192 =====
% ID: L12C3B16-43-TLN
% Mức: TH
\begin{ex}
% ID: L12C3B16-43-TLN
% Mức: TH
Trong bốn phát biểu về Phân tích sự biến thiên từ thông qua khung dây, có bao nhiêu phát biểu đúng? (Chỉ ghi một số nguyên.) (1) Từ thông biến thiên khi $B$, diện tích mạch hoặc góc giữa $B$ và pháp tuyến thay đổi. (2) Từ thông không đổi trong mọi chuyển động của khung dây. (3) Cần kiểm tra điều kiện áp dụng. (4) Có thể bỏ qua đơn vị.
\shortans{2}
\loigiai{
Đối chiếu với kiến thức: Từ thông biến thiên khi $B$, diện tích mạch hoặc góc giữa $B$ và pháp tuyến thay đổi. Có 2 phát biểu đúng.
}
\end{ex}

% ===== Câu 193 =====
% ID: L12C3B16-43-TN
% Mức: NB
\dangbt{Tính suất điện động cảm ứng theo định luật Faraday}
\begin{ex}
% ID: L12C3B16-43-TN
% Mức: NB
Tình huống 3: chọn kết luận chính xác về Tính suất điện động cảm ứng theo định luật Faraday?
\choice
{Kết quả không phụ thuộc dữ kiện và điều kiện áp dụng.}
{\True Độ lớn suất điện động cảm ứng trung bình của cuộn $N$ vòng là |e| = N|Delta Phi|/Delta t.}
{Suất điện động cảm ứng bằng $N$ Delta t/Delta Phi.}
{Có thể bỏ qua phương, chiều và đơn vị trong mọi trường hợp.}
\loigiai{
Độ lớn suất điện động cảm ứng trung bình của cuộn $N$ vòng là |e| = N|Delta Phi|/Delta t. Vì vậy phương án $B$ đúng; các phương án còn lại sai.
}
\end{ex}

% ===== Câu 194 =====
% ID: L12C3B16-44-DS
% Mức: VD
\dangbt{Nhận biết hiện tượng cảm ứng điện từ và điều kiện xuất hiện dòng điện cảm ứng}
\begin{ex}
% ID: L12C3B16-44-DS
% Mức: VD
Xét Nhận biết hiện tượng cảm ứng điện từ và điều kiện xuất hiện dòng điện cảm ứng (Tình huống 4). Hãy xác định tính đúng, sai của bốn phát biểu sau.
\choiceTF
{Kết quả không phụ thuộc dữ kiện và có thể bỏ qua điều kiện.}
{\True Khi giải cần xác định đúng dữ kiện, phương chiều, điều kiện áp dụng và đơn vị.}
{Chỉ cần đặt mạch kín đứng yên trong từ trường đều không đổi là luôn có dòng điện cảm ứng.}
{\True Dòng điện cảm ứng xuất hiện trong mạch kín khi từ thông qua mạch biến thiên.}
\loigiai{
\begin{itemchoice}
\item (Sai) Kết quả không phụ thuộc dữ kiện và có thể bỏ qua điều kiện. (Vì): Kết quả không phụ thuộc dữ kiện và có thể bỏ qua điều kiện. b) Đúng: Khi giải cần xác định đúng dữ kiện, phương chiều, điều kiện áp dụng và đơn vị. c) Sai: Chỉ cần đặt mạch kín đứng yên trong từ trường đều không đổi là luôn có dòng điện cảm ứng. d) Đúng: Dòng điện cảm ứng xuất hiện trong mạch kín khi từ thông qua mạch biến thiên. Kiến thức cốt lõi: Dòng điện cảm ứng xuất hiện trong mạch kín khi từ thông qua mạch biến thiên
\item (Đúng) Khi giải cần xác định đúng dữ kiện, phương chiều, điều kiện áp dụng và đơn vị.
\item (Sai) Chỉ cần đặt mạch kín đứng yên trong từ trường đều không đổi là luôn có dòng điện cảm ứng.
\item (Đúng) Dòng điện cảm ứng xuất hiện trong mạch kín khi từ thông qua mạch biến thiên.
\end{itemchoice}
}
\end{ex}

% ===== Câu 195 =====
% ID: L12C3B16-44-TL
% Mức: TH
\dangbt{Phân tích sự biến thiên từ thông qua khung dây}
\begin{bt}
% ID: L12C3B16-44-TL
% Mức: TH
Giải thích vì sao nhận định sau đúng: “Từ thông biến thiên khi $B$, diện tích mạch hoặc góc giữa $B$ và pháp tuyến thay đổi.”
\loigiai{
Cần nêu được: Từ thông biến thiên khi $B$, diện tích mạch hoặc góc giữa $B$ và pháp tuyến thay đổi. Khi vận dụng phải xác định đúng dữ kiện, phương chiều nếu có, điều kiện áp dụng, hệ đơn vị và kiểm tra tính hợp lí. Nhận định “Từ thông không đổi trong mọi chuyển động của khung dây.” sai.
}
\end{bt}

% ===== Câu 196 =====
% ID: L12C3B16-44-TLN
% Mức: TH
\begin{ex}
% ID: L12C3B16-44-TLN
% Mức: TH
Trong bốn phát biểu về Phân tích sự biến thiên từ thông qua khung dây, có bao nhiêu phát biểu đúng? (Chỉ ghi một số nguyên.) (1) Từ thông không đổi trong mọi chuyển động của khung dây. (2) Từ thông biến thiên khi $B$, diện tích mạch hoặc góc giữa $B$ và pháp tuyến thay đổi. (3) Kết quả phải phù hợp ý nghĩa vật lí. (4) Mọi đại lượng đều là vô hướng.
\shortans{1}
\loigiai{
Đối chiếu với kiến thức: Từ thông biến thiên khi $B$, diện tích mạch hoặc góc giữa $B$ và pháp tuyến thay đổi. Có 1 phát biểu đúng.
}
\end{ex}

% ===== Câu 197 =====
% ID: L12C3B16-44-TN
% Mức: TH
\dangbt{Tính suất điện động cảm ứng theo định luật Faraday}
\begin{ex}
% ID: L12C3B16-44-TN
% Mức: TH
Nội dung nào mô tả đúng nhất Tính suất điện động cảm ứng theo định luật Faraday?
\choice
{Có thể bỏ qua phương, chiều và đơn vị trong mọi trường hợp.}
{Kết quả không phụ thuộc dữ kiện và điều kiện áp dụng.}
{\True Độ lớn suất điện động cảm ứng trung bình của cuộn $N$ vòng là |e| = N|Delta Phi|/Delta t.}
{Suất điện động cảm ứng bằng $N$ Delta t/Delta Phi.}
\loigiai{
Độ lớn suất điện động cảm ứng trung bình của cuộn $N$ vòng là |e| = N|Delta Phi|/Delta t. Vì vậy phương án $C$ đúng; các phương án còn lại sai.
}
\end{ex}

% ===== Câu 198 =====
% ID: L12C3B16-45-DS
% Mức: VD
\dangbt{Nhận biết hiện tượng cảm ứng điện từ và điều kiện xuất hiện dòng điện cảm ứng}
\begin{ex}
% ID: L12C3B16-45-DS
% Mức: VD
Xét Nhận biết hiện tượng cảm ứng điện từ và điều kiện xuất hiện dòng điện cảm ứng (Tình huống 5). Hãy xác định tính đúng, sai của bốn phát biểu sau.
\choiceTF
{\True Dòng điện cảm ứng xuất hiện trong mạch kín khi từ thông qua mạch biến thiên.}
{Kết quả không phụ thuộc dữ kiện và có thể bỏ qua điều kiện.}
{Chỉ cần đặt mạch kín đứng yên trong từ trường đều không đổi là luôn có dòng điện cảm ứng.}
{\True Khi giải cần xác định đúng dữ kiện, phương chiều, điều kiện áp dụng và đơn vị.}
\loigiai{
\begin{itemchoice}
\item (Đúng) Dòng điện cảm ứng xuất hiện trong mạch kín khi từ thông qua mạch biến thiên. (Vì): òng điện cảm ứng xuất hiện trong mạch kín khi từ thông qua mạch biến thiên. b) Sai: Kết quả không phụ thuộc dữ kiện và có thể bỏ qua điều kiện. c) Sai: Chỉ cần đặt mạch kín đứng yên trong từ trường đều không đổi là luôn có dòng điện cảm ứng. d) Đúng: Khi giải cần xác định đúng dữ kiện, phương chiều, điều kiện áp dụng và đơn vị. Kiến thức cốt lõi: Dòng điện cảm ứng xuất hiện trong mạch kín khi từ thông qua mạch biến thiên
\item (Sai) Kết quả không phụ thuộc dữ kiện và có thể bỏ qua điều kiện.
\item (Sai) Chỉ cần đặt mạch kín đứng yên trong từ trường đều không đổi là luôn có dòng điện cảm ứng.
\item (Đúng) Khi giải cần xác định đúng dữ kiện, phương chiều, điều kiện áp dụng và đơn vị.
\end{itemchoice}
}
\end{ex}

% ===== Câu 199 =====
% ID: L12C3B16-45-TL
% Mức: VD
\dangbt{Phân tích sự biến thiên từ thông qua khung dây}
\begin{bt}
% ID: L12C3B16-45-TL
% Mức: VD
Phân tích sai lầm trong nhận định: “Từ thông không đổi trong mọi chuyển động của khung dây.”
\loigiai{
Cần nêu được: Từ thông biến thiên khi $B$, diện tích mạch hoặc góc giữa $B$ và pháp tuyến thay đổi. Khi vận dụng phải xác định đúng dữ kiện, phương chiều nếu có, điều kiện áp dụng, hệ đơn vị và kiểm tra tính hợp lí. Nhận định “Từ thông không đổi trong mọi chuyển động của khung dây.” sai.
}
\end{bt}

% ===== Câu 200 =====
% ID: L12C3B16-45-TLN
% Mức: TH
\begin{ex}
% ID: L12C3B16-45-TLN
% Mức: TH
Trong bốn phát biểu về Phân tích sự biến thiên từ thông qua khung dây, có bao nhiêu phát biểu đúng? (Chỉ ghi một số nguyên.) (1) Từ thông biến thiên khi $B$, diện tích mạch hoặc góc giữa $B$ và pháp tuyến thay đổi. (2) Từ thông không đổi trong mọi chuyển động của khung dây. (3) Cần kiểm tra điều kiện áp dụng. (4) Có thể bỏ qua đơn vị.
\shortans{2}
\loigiai{
Đối chiếu với kiến thức: Từ thông biến thiên khi $B$, diện tích mạch hoặc góc giữa $B$ và pháp tuyến thay đổi. Có 2 phát biểu đúng.
}
\end{ex}

% ===== Câu 201 =====
% ID: L12C3B16-45-TN
% Mức: TH
\dangbt{Tính suất điện động cảm ứng theo định luật Faraday}
\begin{ex}
% ID: L12C3B16-45-TN
% Mức: TH
Khi phân tích Tính suất điện động cảm ứng theo định luật Faraday?
\choice
{Suất điện động cảm ứng bằng $N$ Delta t/Delta Phi.}
{Có thể bỏ qua phương, chiều và đơn vị trong mọi trường hợp.}
{Kết quả không phụ thuộc dữ kiện và điều kiện áp dụng.}
{\True Độ lớn suất điện động cảm ứng trung bình của cuộn $N$ vòng là |e| = N|Delta Phi|/Delta t.}
\loigiai{
Độ lớn suất điện động cảm ứng trung bình của cuộn $N$ vòng là |e| = N|Delta Phi|/Delta t. Vì vậy phương án $D$ đúng; các phương án còn lại sai.
}
\end{ex}

% ===== Câu 202 =====
% ID: L12C3B16-46-DS
% Mức: TH
\dangbt{Xác định chiều dòng điện cảm ứng bằng định luật Lenz}
\begin{ex}
% ID: L12C3B16-46-DS
% Mức: TH
Xét Xác định chiều dòng điện cảm ứng bằng định luật Lenz (Tình huống 1). Hãy xác định tính đúng, sai của bốn phát biểu sau.
\choiceTF
{\True Dòng điện cảm ứng có chiều sao cho từ trường do nó sinh ra chống lại sự biến thiên từ thông đã gây ra nó.}
{Dòng điện cảm ứng luôn làm tăng biến thiên từ thông ban đầu.}
{\True Khi giải cần xác định đúng dữ kiện, phương chiều, điều kiện áp dụng và đơn vị.}
{Kết quả không phụ thuộc dữ kiện và có thể bỏ qua điều kiện.}
\loigiai{
\begin{itemchoice}
\item (Đúng) Dòng điện cảm ứng có chiều sao cho từ trường do nó sinh ra chống lại sự biến thiên từ thông đã gây ra nó. (Vì): òng điện cảm ứng có chiều sao cho từ trường do nó sinh ra chống lại sự biến thiên từ thông đã gây ra nó. b) Sai: Dòng điện cảm ứng luôn làm tăng biến thiên từ thông ban đầu. c) Đúng: Khi giải cần xác định đúng dữ kiện, phương chiều, điều kiện áp dụng và đơn vị. d) Sai: Kết quả không phụ thuộc dữ kiện và có thể bỏ qua điều kiện. Kiến thức cốt lõi: Dòng điện cảm ứng có chiều sao cho từ trường do nó sinh ra chống lại sự biến thiên từ thông đã gây ra nó
\item (Sai) Dòng điện cảm ứng luôn làm tăng biến thiên từ thông ban đầu.
\item (Đúng) Khi giải cần xác định đúng dữ kiện, phương chiều, điều kiện áp dụng và đơn vị.
\item (Sai) Kết quả không phụ thuộc dữ kiện và có thể bỏ qua điều kiện.
\end{itemchoice}
}
\end{ex}

% ===== Câu 203 =====
% ID: L12C3B16-46-TL
% Mức: VD
\dangbt{Phân tích sự biến thiên từ thông qua khung dây}
\begin{bt}
% ID: L12C3B16-46-TL
% Mức: VD
Đề xuất các bước giải một bài toán thuộc dạng Phân tích sự biến thiên từ thông qua khung dây.
\loigiai{
Cần nêu được: Từ thông biến thiên khi $B$, diện tích mạch hoặc góc giữa $B$ và pháp tuyến thay đổi. Khi vận dụng phải xác định đúng dữ kiện, phương chiều nếu có, điều kiện áp dụng, hệ đơn vị và kiểm tra tính hợp lí. Nhận định “Từ thông không đổi trong mọi chuyển động của khung dây.” sai.
}
\end{bt}

% ===== Câu 204 =====
% ID: L12C3B16-46-TLN
% Mức: VD
\begin{ex}
% ID: L12C3B16-46-TLN
% Mức: VD
Trong bốn phát biểu về Phân tích sự biến thiên từ thông qua khung dây, có bao nhiêu phát biểu đúng? (Chỉ ghi một số nguyên.) (1) Từ thông không đổi trong mọi chuyển động của khung dây. (2) Từ thông biến thiên khi $B$, diện tích mạch hoặc góc giữa $B$ và pháp tuyến thay đổi. (3) Kết quả phải phù hợp ý nghĩa vật lí. (4) Mọi đại lượng đều là vô hướng.
\shortans{2}
\loigiai{
Đối chiếu với kiến thức: Từ thông biến thiên khi $B$, diện tích mạch hoặc góc giữa $B$ và pháp tuyến thay đổi. Có 2 phát biểu đúng.
}
\end{ex}

% ===== Câu 205 =====
% ID: L12C3B16-46-TN
% Mức: TH
\dangbt{Tính suất điện động cảm ứng theo định luật Faraday}
\begin{ex}
% ID: L12C3B16-46-TN
% Mức: TH
Một học sinh ôn tập Tính suất điện động cảm ứng theo định luật Faraday?
\choice
{\True Độ lớn suất điện động cảm ứng trung bình của cuộn $N$ vòng là |e| = N|Delta Phi|/Delta t.}
{Suất điện động cảm ứng bằng $N$ Delta t/Delta Phi.}
{Có thể bỏ qua phương, chiều và đơn vị trong mọi trường hợp.}
{Kết quả không phụ thuộc dữ kiện và điều kiện áp dụng.}
\loigiai{
Độ lớn suất điện động cảm ứng trung bình của cuộn $N$ vòng là |e| = N|Delta Phi|/Delta t. Vì vậy phương án $A$ đúng; các phương án còn lại sai.
}
\end{ex}

% ===== Câu 206 =====
% ID: L12C3B16-47-DS
% Mức: TH
\dangbt{Xác định chiều dòng điện cảm ứng bằng định luật Lenz}
\begin{ex}
% ID: L12C3B16-47-DS
% Mức: TH
Xét Xác định chiều dòng điện cảm ứng bằng định luật Lenz (Tình huống 2). Hãy xác định tính đúng, sai của bốn phát biểu sau.
\choiceTF
{Dòng điện cảm ứng luôn làm tăng biến thiên từ thông ban đầu.}
{\True Dòng điện cảm ứng có chiều sao cho từ trường do nó sinh ra chống lại sự biến thiên từ thông đã gây ra nó.}
{Kết quả không phụ thuộc dữ kiện và có thể bỏ qua điều kiện.}
{\True Khi giải cần xác định đúng dữ kiện, phương chiều, điều kiện áp dụng và đơn vị.}
\loigiai{
\begin{itemchoice}
\item (Sai) Dòng điện cảm ứng luôn làm tăng biến thiên từ thông ban đầu. (Vì): Dòng điện cảm ứng luôn làm tăng biến thiên từ thông ban đầu. b) Đúng: Dòng điện cảm ứng có chiều sao cho từ trường do nó sinh ra chống lại sự biến thiên từ thông đã gây ra nó. c) Sai: Kết quả không phụ thuộc dữ kiện và có thể bỏ qua điều kiện. d) Đúng: Khi giải cần xác định đúng dữ kiện, phương chiều, điều kiện áp dụng và đơn vị. Kiến thức cốt lõi: Dòng điện cảm ứng có chiều sao cho từ trường do nó sinh ra chống lại sự biến thiên từ thông đã gây ra nó
\item (Đúng) Dòng điện cảm ứng có chiều sao cho từ trường do nó sinh ra chống lại sự biến thiên từ thông đã gây ra nó.
\item (Sai) Kết quả không phụ thuộc dữ kiện và có thể bỏ qua điều kiện.
\item (Đúng) Khi giải cần xác định đúng dữ kiện, phương chiều, điều kiện áp dụng và đơn vị.
\end{itemchoice}
}
\end{ex}

% ===== Câu 207 =====
% ID: L12C3B16-47-TL
% Mức: VD
\dangbt{Phân tích sự biến thiên từ thông qua khung dây}
\begin{bt}
% ID: L12C3B16-47-TL
% Mức: VD
Nêu một tình huống thực tiễn liên quan đến Phân tích sự biến thiên từ thông qua khung dây và giải thích bằng kiến thức Vật lí.
\loigiai{
Cần nêu được: Từ thông biến thiên khi $B$, diện tích mạch hoặc góc giữa $B$ và pháp tuyến thay đổi. Khi vận dụng phải xác định đúng dữ kiện, phương chiều nếu có, điều kiện áp dụng, hệ đơn vị và kiểm tra tính hợp lí. Nhận định “Từ thông không đổi trong mọi chuyển động của khung dây.” sai.
}
\end{bt}

% ===== Câu 208 =====
% ID: L12C3B16-47-TLN
% Mức: VD
\begin{ex}
% ID: L12C3B16-47-TLN
% Mức: VD
Trong bốn phát biểu về Phân tích sự biến thiên từ thông qua khung dây, có bao nhiêu phát biểu đúng? (Chỉ ghi một số nguyên.) (1) Từ thông biến thiên khi $B$, diện tích mạch hoặc góc giữa $B$ và pháp tuyến thay đổi. (2) Từ thông không đổi trong mọi chuyển động của khung dây. (3) Cần kiểm tra điều kiện áp dụng. (4) Có thể bỏ qua đơn vị.
\shortans{2}
\loigiai{
Đối chiếu với kiến thức: Từ thông biến thiên khi $B$, diện tích mạch hoặc góc giữa $B$ và pháp tuyến thay đổi. Có 2 phát biểu đúng.
}
\end{ex}

% ===== Câu 209 =====
% ID: L12C3B16-47-TN
% Mức: TH
\dangbt{Tính suất điện động cảm ứng theo định luật Faraday}
\begin{ex}
% ID: L12C3B16-47-TN
% Mức: TH
Chọn phát biểu không trái với kiến thức về Tính suất điện động cảm ứng theo định luật Faraday?
\choice
{Kết quả không phụ thuộc dữ kiện và điều kiện áp dụng.}
{\True Độ lớn suất điện động cảm ứng trung bình của cuộn $N$ vòng là |e| = N|Delta Phi|/Delta t.}
{Suất điện động cảm ứng bằng $N$ Delta t/Delta Phi.}
{Có thể bỏ qua phương, chiều và đơn vị trong mọi trường hợp.}
\loigiai{
Độ lớn suất điện động cảm ứng trung bình của cuộn $N$ vòng là |e| = N|Delta Phi|/Delta t. Vì vậy phương án $B$ đúng; các phương án còn lại sai.
}
\end{ex}

% ===== Câu 210 =====
% ID: L12C3B16-48-DS
% Mức: VD
\dangbt{Xác định chiều dòng điện cảm ứng bằng định luật Lenz}
\begin{ex}
% ID: L12C3B16-48-DS
% Mức: VD
Xét Xác định chiều dòng điện cảm ứng bằng định luật Lenz (Tình huống 3). Hãy xác định tính đúng, sai của bốn phát biểu sau.
\choiceTF
{\True Khi giải cần xác định đúng dữ kiện, phương chiều, điều kiện áp dụng và đơn vị.}
{Kết quả không phụ thuộc dữ kiện và có thể bỏ qua điều kiện.}
{\True Dòng điện cảm ứng có chiều sao cho từ trường do nó sinh ra chống lại sự biến thiên từ thông đã gây ra nó.}
{Dòng điện cảm ứng luôn làm tăng biến thiên từ thông ban đầu.}
\loigiai{
\begin{itemchoice}
\item (Đúng) Khi giải cần xác định đúng dữ kiện, phương chiều, điều kiện áp dụng và đơn vị. (Vì): Khi giải cần xác định đúng dữ kiện, phương chiều, điều kiện áp dụng và đơn vị. b) Sai: Kết quả không phụ thuộc dữ kiện và có thể bỏ qua điều kiện. c) Đúng: Dòng điện cảm ứng có chiều sao cho từ trường do nó sinh ra chống lại sự biến thiên từ thông đã gây ra nó. d) Sai: Dòng điện cảm ứng luôn làm tăng biến thiên từ thông ban đầu. Kiến thức cốt lõi: Dòng điện cảm ứng có chiều sao cho từ trường do nó sinh ra chống lại sự biến thiên từ thông đã gây ra nó
\item (Sai) Kết quả không phụ thuộc dữ kiện và có thể bỏ qua điều kiện.
\item (Đúng) Dòng điện cảm ứng có chiều sao cho từ trường do nó sinh ra chống lại sự biến thiên từ thông đã gây ra nó.
\item (Sai) Dòng điện cảm ứng luôn làm tăng biến thiên từ thông ban đầu.
\end{itemchoice}
}
\end{ex}

% ===== Câu 211 =====
% ID: L12C3B16-48-TL
% Mức: VDC
\dangbt{Phân tích sự biến thiên từ thông qua khung dây}
\begin{bt}
% ID: L12C3B16-48-TL
% Mức: VDC
So sánh nhận định đúng “Từ thông biến thiên khi $B$, diện tích mạch hoặc góc giữa $B$ và pháp tuyến thay đổi.” với nhận định sai “Từ thông không đổi trong mọi chuyển động của khung dây.”; chỉ rõ căn cứ Vật lí.
\loigiai{
Cần nêu được: Từ thông biến thiên khi $B$, diện tích mạch hoặc góc giữa $B$ và pháp tuyến thay đổi. Khi vận dụng phải xác định đúng dữ kiện, phương chiều nếu có, điều kiện áp dụng, hệ đơn vị và kiểm tra tính hợp lí. Nhận định “Từ thông không đổi trong mọi chuyển động của khung dây.” sai.
}
\end{bt}

% ===== Câu 212 =====
% ID: L12C3B16-48-TLN
% Mức: VDC
\begin{ex}
% ID: L12C3B16-48-TLN
% Mức: VDC
Trong bốn phát biểu về Phân tích sự biến thiên từ thông qua khung dây, có bao nhiêu phát biểu đúng? (Chỉ ghi một số nguyên.) (1) Khi giải cần xác định đúng dữ kiện, phương chiều, điều kiện áp dụng và đơn vị. (2) Từ thông không đổi trong mọi chuyển động của khung dây. (3) Phải dùng đúng định luật cho tình huống. (4) Từ thông biến thiên khi $B$, diện tích mạch hoặc góc giữa $B$ và pháp tuyến thay đổi.
\shortans{3}
\loigiai{
Đối chiếu với kiến thức: Từ thông biến thiên khi $B$, diện tích mạch hoặc góc giữa $B$ và pháp tuyến thay đổi. Có 3 phát biểu đúng.
}
\end{ex}

% ===== Câu 213 =====
% ID: L12C3B16-48-TN
% Mức: VD
\dangbt{Tính suất điện động cảm ứng theo định luật Faraday}
\begin{ex}
% ID: L12C3B16-48-TN
% Mức: VD
Cơ sở Vật lí đúng dùng để giải Tính suất điện động cảm ứng theo định luật Faraday?
\choice
{Có thể bỏ qua phương, chiều và đơn vị trong mọi trường hợp.}
{Kết quả không phụ thuộc dữ kiện và điều kiện áp dụng.}
{\True Độ lớn suất điện động cảm ứng trung bình của cuộn $N$ vòng là |e| = N|Delta Phi|/Delta t.}
{Suất điện động cảm ứng bằng $N$ Delta t/Delta Phi.}
\loigiai{
Độ lớn suất điện động cảm ứng trung bình của cuộn $N$ vòng là |e| = N|Delta Phi|/Delta t. Vì vậy phương án $C$ đúng; các phương án còn lại sai.
}
\end{ex}

% ===== Câu 214 =====
% ID: L12C3B16-49-DS
% Mức: VD
\dangbt{Xác định chiều dòng điện cảm ứng bằng định luật Lenz}
\begin{ex}
% ID: L12C3B16-49-DS
% Mức: VD
Xét Xác định chiều dòng điện cảm ứng bằng định luật Lenz (Tình huống 4). Hãy xác định tính đúng, sai của bốn phát biểu sau.
\choiceTF
{Kết quả không phụ thuộc dữ kiện và có thể bỏ qua điều kiện.}
{\True Khi giải cần xác định đúng dữ kiện, phương chiều, điều kiện áp dụng và đơn vị.}
{Dòng điện cảm ứng luôn làm tăng biến thiên từ thông ban đầu.}
{\True Dòng điện cảm ứng có chiều sao cho từ trường do nó sinh ra chống lại sự biến thiên từ thông đã gây ra nó.}
\loigiai{
\begin{itemchoice}
\item (Sai) Kết quả không phụ thuộc dữ kiện và có thể bỏ qua điều kiện. (Vì): Kết quả không phụ thuộc dữ kiện và có thể bỏ qua điều kiện. b) Đúng: Khi giải cần xác định đúng dữ kiện, phương chiều, điều kiện áp dụng và đơn vị. c) Sai: Dòng điện cảm ứng luôn làm tăng biến thiên từ thông ban đầu. d) Đúng: Dòng điện cảm ứng có chiều sao cho từ trường do nó sinh ra chống lại sự biến thiên từ thông đã gây ra nó. Kiến thức cốt lõi: Dòng điện cảm ứng có chiều sao cho từ trường do nó sinh ra chống lại sự biến thiên từ thông đã gây ra nó
\item (Đúng) Khi giải cần xác định đúng dữ kiện, phương chiều, điều kiện áp dụng và đơn vị.
\item (Sai) Dòng điện cảm ứng luôn làm tăng biến thiên từ thông ban đầu.
\item (Đúng) Dòng điện cảm ứng có chiều sao cho từ trường do nó sinh ra chống lại sự biến thiên từ thông đã gây ra nó.
\end{itemchoice}
}
\end{ex}

% ===== Câu 215 =====
% ID: L12C3B16-49-TL
% Mức: TH
\dangbt{Nhận biết hiện tượng cảm ứng điện từ và điều kiện xuất hiện dòng điện cảm ứng}
\begin{bt}
% ID: L12C3B16-49-TL
% Mức: TH
Trình bày kiến thức cốt lõi của Nhận biết hiện tượng cảm ứng điện từ và điều kiện xuất hiện dòng điện cảm ứng và nêu điều kiện áp dụng.
\loigiai{
Cần nêu được: Dòng điện cảm ứng xuất hiện trong mạch kín khi từ thông qua mạch biến thiên. Khi vận dụng phải xác định đúng dữ kiện, phương chiều nếu có, điều kiện áp dụng, hệ đơn vị và kiểm tra tính hợp lí. Nhận định “Chỉ cần đặt mạch kín đứng yên trong từ trường đều không đổi là luôn có dòng điện cảm ứng.” sai.
}
\end{bt}

% ===== Câu 216 =====
% ID: L12C3B16-49-TLN
% Mức: TH
\begin{ex}
% ID: L12C3B16-49-TLN
% Mức: TH
Trong bốn phát biểu về Nhận biết hiện tượng cảm ứng điện từ và điều kiện xuất hiện dòng điện cảm ứng, có bao nhiêu phát biểu đúng? (Chỉ ghi một số nguyên.) (1) Dòng điện cảm ứng xuất hiện trong mạch kín khi từ thông qua mạch biến thiên. (2) Chỉ cần đặt mạch kín đứng yên trong từ trường đều không đổi là luôn có dòng điện cảm ứng. (3) Cần kiểm tra điều kiện áp dụng. (4) Có thể bỏ qua đơn vị.
\shortans{2}
\loigiai{
Đối chiếu với kiến thức: Dòng điện cảm ứng xuất hiện trong mạch kín khi từ thông qua mạch biến thiên. Có 2 phát biểu đúng.
}
\end{ex}

% ===== Câu 217 =====
% ID: L12C3B16-49-TN
% Mức: VD
\dangbt{Tính suất điện động cảm ứng theo định luật Faraday}
\begin{ex}
% ID: L12C3B16-49-TN
% Mức: VD
Trong một bài thực hành về Tính suất điện động cảm ứng theo định luật Faraday?
\choice
{Suất điện động cảm ứng bằng $N$ Delta t/Delta Phi.}
{Có thể bỏ qua phương, chiều và đơn vị trong mọi trường hợp.}
{Kết quả không phụ thuộc dữ kiện và điều kiện áp dụng.}
{\True Độ lớn suất điện động cảm ứng trung bình của cuộn $N$ vòng là |e| = N|Delta Phi|/Delta t.}
\loigiai{
Độ lớn suất điện động cảm ứng trung bình của cuộn $N$ vòng là |e| = N|Delta Phi|/Delta t. Vì vậy phương án $D$ đúng; các phương án còn lại sai.
}
\end{ex}

% ===== Câu 218 =====
% ID: L12C3B16-50-DS
% Mức: VD
\dangbt{Xác định chiều dòng điện cảm ứng bằng định luật Lenz}
\begin{ex}
% ID: L12C3B16-50-DS
% Mức: VD
Xét Xác định chiều dòng điện cảm ứng bằng định luật Lenz (Tình huống 5). Hãy xác định tính đúng, sai của bốn phát biểu sau.
\choiceTF
{\True Dòng điện cảm ứng có chiều sao cho từ trường do nó sinh ra chống lại sự biến thiên từ thông đã gây ra nó.}
{Kết quả không phụ thuộc dữ kiện và có thể bỏ qua điều kiện.}
{Dòng điện cảm ứng luôn làm tăng biến thiên từ thông ban đầu.}
{\True Khi giải cần xác định đúng dữ kiện, phương chiều, điều kiện áp dụng và đơn vị.}
\loigiai{
\begin{itemchoice}
\item (Đúng) Dòng điện cảm ứng có chiều sao cho từ trường do nó sinh ra chống lại sự biến thiên từ thông đã gây ra nó. (Vì): òng điện cảm ứng có chiều sao cho từ trường do nó sinh ra chống lại sự biến thiên từ thông đã gây ra nó. b) Sai: Kết quả không phụ thuộc dữ kiện và có thể bỏ qua điều kiện. c) Sai: Dòng điện cảm ứng luôn làm tăng biến thiên từ thông ban đầu. d) Đúng: Khi giải cần xác định đúng dữ kiện, phương chiều, điều kiện áp dụng và đơn vị. Kiến thức cốt lõi: Dòng điện cảm ứng có chiều sao cho từ trường do nó sinh ra chống lại sự biến thiên từ thông đã gây ra nó
\item (Sai) Kết quả không phụ thuộc dữ kiện và có thể bỏ qua điều kiện.
\item (Sai) Dòng điện cảm ứng luôn làm tăng biến thiên từ thông ban đầu.
\item (Đúng) Khi giải cần xác định đúng dữ kiện, phương chiều, điều kiện áp dụng và đơn vị.
\end{itemchoice}
}
\end{ex}

% ===== Câu 219 =====
% ID: L12C3B16-50-TL
% Mức: TH
\dangbt{Nhận biết hiện tượng cảm ứng điện từ và điều kiện xuất hiện dòng điện cảm ứng}
\begin{bt}
% ID: L12C3B16-50-TL
% Mức: TH
Giải thích vì sao nhận định sau đúng: “Dòng điện cảm ứng xuất hiện trong mạch kín khi từ thông qua mạch biến thiên.”
\loigiai{
Cần nêu được: Dòng điện cảm ứng xuất hiện trong mạch kín khi từ thông qua mạch biến thiên. Khi vận dụng phải xác định đúng dữ kiện, phương chiều nếu có, điều kiện áp dụng, hệ đơn vị và kiểm tra tính hợp lí. Nhận định “Chỉ cần đặt mạch kín đứng yên trong từ trường đều không đổi là luôn có dòng điện cảm ứng.” sai.
}
\end{bt}

% ===== Câu 220 =====
% ID: L12C3B16-50-TLN
% Mức: TH
\begin{ex}
% ID: L12C3B16-50-TLN
% Mức: TH
Trong bốn phát biểu về Nhận biết hiện tượng cảm ứng điện từ và điều kiện xuất hiện dòng điện cảm ứng, có bao nhiêu phát biểu đúng? (Chỉ ghi một số nguyên.) (1) Chỉ cần đặt mạch kín đứng yên trong từ trường đều không đổi là luôn có dòng điện cảm ứng. (2) Dòng điện cảm ứng xuất hiện trong mạch kín khi từ thông qua mạch biến thiên. (3) Kết quả phải phù hợp ý nghĩa vật lí. (4) Mọi đại lượng đều là vô hướng.
\shortans{1}
\loigiai{
Đối chiếu với kiến thức: Dòng điện cảm ứng xuất hiện trong mạch kín khi từ thông qua mạch biến thiên. Có 1 phát biểu đúng.
}
\end{ex}

% ===== Câu 221 =====
% ID: L12C3B16-50-TN
% Mức: VD
\dangbt{Tính suất điện động cảm ứng theo định luật Faraday}
\begin{ex}
% ID: L12C3B16-50-TN
% Mức: VD
Kết luận đúng khi vận dụng Tính suất điện động cảm ứng theo định luật Faraday?
\choice
{\True Độ lớn suất điện động cảm ứng trung bình của cuộn $N$ vòng là |e| = N|Delta Phi|/Delta t.}
{Suất điện động cảm ứng bằng $N$ Delta t/Delta Phi.}
{Có thể bỏ qua phương, chiều và đơn vị trong mọi trường hợp.}
{Kết quả không phụ thuộc dữ kiện và điều kiện áp dụng.}
\loigiai{
Độ lớn suất điện động cảm ứng trung bình của cuộn $N$ vòng là |e| = N|Delta Phi|/Delta t. Vì vậy phương án $A$ đúng; các phương án còn lại sai.
}
\end{ex}

% ===== Câu 222 =====
% ID: L12C3B16-51-DS
% Mức: TH
\begin{ex}
% ID: L12C3B16-51-DS
% Mức: TH
Xét Tính suất điện động cảm ứng theo định luật Faraday (Tình huống 1). Hãy xác định tính đúng, sai của bốn phát biểu sau.
\choiceTF
{\True Độ lớn suất điện động cảm ứng trung bình của cuộn $N$ vòng là |e| = N|Delta Phi|/Delta t.}
{Suất điện động cảm ứng bằng $N$ Delta t/Delta Phi.}
{\True Khi giải cần xác định đúng dữ kiện, phương chiều, điều kiện áp dụng và đơn vị.}
{Kết quả không phụ thuộc dữ kiện và có thể bỏ qua điều kiện.}
\loigiai{
\begin{itemchoice}
\item (Đúng) Độ lớn suất điện động cảm ứng trung bình của cuộn $N$ vòng là |e| = N|Delta Phi|/Delta t. (Vì): ộ lớn suất điện động cảm ứng trung bình của cuộn $N$ vòng là |e| = N|Delta Phi|/Delta t. b) Sai: Suất điện động cảm ứng bằng $N$ Delta t/Delta Phi. c) Đúng: Khi giải cần xác định đúng dữ kiện, phương chiều, điều kiện áp dụng và đơn vị. d) Sai: Kết quả không phụ thuộc dữ kiện và có thể bỏ qua điều kiện. Kiến thức cốt lõi: Độ lớn suất điện động cảm ứng trung bình của cuộn $N$ vòng là |e| = N|Delta Phi|/Delta t
\item (Sai) Suất điện động cảm ứng bằng $N$ Delta t/Delta Phi.
\item (Đúng) Khi giải cần xác định đúng dữ kiện, phương chiều, điều kiện áp dụng và đơn vị.
\item (Sai) Kết quả không phụ thuộc dữ kiện và có thể bỏ qua điều kiện.
\end{itemchoice}
}
\end{ex}

% ===== Câu 223 =====
% ID: L12C3B16-51-TL
% Mức: VD
\dangbt{Nhận biết hiện tượng cảm ứng điện từ và điều kiện xuất hiện dòng điện cảm ứng}
\begin{bt}
% ID: L12C3B16-51-TL
% Mức: VD
Phân tích sai lầm trong nhận định: “Chỉ cần đặt mạch kín đứng yên trong từ trường đều không đổi là luôn có dòng điện cảm ứng.”
\loigiai{
Cần nêu được: Dòng điện cảm ứng xuất hiện trong mạch kín khi từ thông qua mạch biến thiên. Khi vận dụng phải xác định đúng dữ kiện, phương chiều nếu có, điều kiện áp dụng, hệ đơn vị và kiểm tra tính hợp lí. Nhận định “Chỉ cần đặt mạch kín đứng yên trong từ trường đều không đổi là luôn có dòng điện cảm ứng.” sai.
}
\end{bt}

% ===== Câu 224 =====
% ID: L12C3B16-51-TLN
% Mức: TH
\begin{ex}
% ID: L12C3B16-51-TLN
% Mức: TH
Trong bốn phát biểu về Nhận biết hiện tượng cảm ứng điện từ và điều kiện xuất hiện dòng điện cảm ứng, có bao nhiêu phát biểu đúng? (Chỉ ghi một số nguyên.) (1) Dòng điện cảm ứng xuất hiện trong mạch kín khi từ thông qua mạch biến thiên. (2) Chỉ cần đặt mạch kín đứng yên trong từ trường đều không đổi là luôn có dòng điện cảm ứng. (3) Cần kiểm tra điều kiện áp dụng. (4) Có thể bỏ qua đơn vị.
\shortans{2}
\loigiai{
Đối chiếu với kiến thức: Dòng điện cảm ứng xuất hiện trong mạch kín khi từ thông qua mạch biến thiên. Có 2 phát biểu đúng.
}
\end{ex}

% ===== Câu 225 =====
% ID: L12C3B16-51-TN
% Mức: NB
\dangbt{Xác định chiều dòng điện cảm ứng bằng định luật Lenz}
\begin{ex}
% ID: L12C3B16-51-TN
% Mức: NB
Phát biểu nào sau đây đúng về Xác định chiều dòng điện cảm ứng bằng định luật Lenz?
\choice
{Dòng điện cảm ứng luôn làm tăng biến thiên từ thông ban đầu.}
{Có thể bỏ qua phương, chiều và đơn vị trong mọi trường hợp.}
{Kết quả không phụ thuộc dữ kiện và điều kiện áp dụng.}
{\True Dòng điện cảm ứng có chiều sao cho từ trường do nó sinh ra chống lại sự biến thiên từ thông đã gây ra nó.}
\loigiai{
Dòng điện cảm ứng có chiều sao cho từ trường do nó sinh ra chống lại sự biến thiên từ thông đã gây ra nó. Vì vậy phương án $D$ đúng; các phương án còn lại sai.
}
\end{ex}

% ===== Câu 226 =====
% ID: L12C3B16-52-DS
% Mức: TH
\dangbt{Tính suất điện động cảm ứng theo định luật Faraday}
\begin{ex}
% ID: L12C3B16-52-DS
% Mức: TH
Xét Tính suất điện động cảm ứng theo định luật Faraday (Tình huống 2). Hãy xác định tính đúng, sai của bốn phát biểu sau.
\choiceTF
{Suất điện động cảm ứng bằng $N$ Delta t/Delta Phi.}
{\True Độ lớn suất điện động cảm ứng trung bình của cuộn $N$ vòng là |e| = N|Delta Phi|/Delta t.}
{Kết quả không phụ thuộc dữ kiện và có thể bỏ qua điều kiện.}
{\True Khi giải cần xác định đúng dữ kiện, phương chiều, điều kiện áp dụng và đơn vị.}
\loigiai{
\begin{itemchoice}
\item (Sai) Suất điện động cảm ứng bằng $N$ Delta t/Delta Phi. (Vì): uất điện động cảm ứng bằng $N$ Delta t/Delta Phi. b) Đúng: Độ lớn suất điện động cảm ứng trung bình của cuộn $N$ vòng là |e| = N|Delta Phi|/Delta t. c) Sai: Kết quả không phụ thuộc dữ kiện và có thể bỏ qua điều kiện. d) Đúng: Khi giải cần xác định đúng dữ kiện, phương chiều, điều kiện áp dụng và đơn vị. Kiến thức cốt lõi: Độ lớn suất điện động cảm ứng trung bình của cuộn $N$ vòng là |e| = N|Delta Phi|/Delta t
\item (Đúng) Độ lớn suất điện động cảm ứng trung bình của cuộn $N$ vòng là |e| = N|Delta Phi|/Delta t.
\item (Sai) Kết quả không phụ thuộc dữ kiện và có thể bỏ qua điều kiện.
\item (Đúng) Khi giải cần xác định đúng dữ kiện, phương chiều, điều kiện áp dụng và đơn vị.
\end{itemchoice}
}
\end{ex}

% ===== Câu 227 =====
% ID: L12C3B16-52-TL
% Mức: VD
\dangbt{Nhận biết hiện tượng cảm ứng điện từ và điều kiện xuất hiện dòng điện cảm ứng}
\begin{bt}
% ID: L12C3B16-52-TL
% Mức: VD
Đề xuất các bước giải một bài toán thuộc dạng Nhận biết hiện tượng cảm ứng điện từ và điều kiện xuất hiện dòng điện cảm ứng.
\loigiai{
Cần nêu được: Dòng điện cảm ứng xuất hiện trong mạch kín khi từ thông qua mạch biến thiên. Khi vận dụng phải xác định đúng dữ kiện, phương chiều nếu có, điều kiện áp dụng, hệ đơn vị và kiểm tra tính hợp lí. Nhận định “Chỉ cần đặt mạch kín đứng yên trong từ trường đều không đổi là luôn có dòng điện cảm ứng.” sai.
}
\end{bt}

% ===== Câu 228 =====
% ID: L12C3B16-52-TLN
% Mức: VD
\begin{ex}
% ID: L12C3B16-52-TLN
% Mức: VD
Trong bốn phát biểu về Nhận biết hiện tượng cảm ứng điện từ và điều kiện xuất hiện dòng điện cảm ứng, có bao nhiêu phát biểu đúng? (Chỉ ghi một số nguyên.) (1) Chỉ cần đặt mạch kín đứng yên trong từ trường đều không đổi là luôn có dòng điện cảm ứng. (2) Dòng điện cảm ứng xuất hiện trong mạch kín khi từ thông qua mạch biến thiên. (3) Kết quả phải phù hợp ý nghĩa vật lí. (4) Mọi đại lượng đều là vô hướng.
\shortans{2}
\loigiai{
Đối chiếu với kiến thức: Dòng điện cảm ứng xuất hiện trong mạch kín khi từ thông qua mạch biến thiên. Có 2 phát biểu đúng.
}
\end{ex}

% ===== Câu 229 =====
% ID: L12C3B16-52-TN
% Mức: NB
\dangbt{Xác định chiều dòng điện cảm ứng bằng định luật Lenz}
\begin{ex}
% ID: L12C3B16-52-TN
% Mức: NB
Trong nội dung Xác định chiều dòng điện cảm ứng bằng định luật Lenz?
\choice
{\True Dòng điện cảm ứng có chiều sao cho từ trường do nó sinh ra chống lại sự biến thiên từ thông đã gây ra nó.}
{Dòng điện cảm ứng luôn làm tăng biến thiên từ thông ban đầu.}
{Có thể bỏ qua phương, chiều và đơn vị trong mọi trường hợp.}
{Kết quả không phụ thuộc dữ kiện và điều kiện áp dụng.}
\loigiai{
Dòng điện cảm ứng có chiều sao cho từ trường do nó sinh ra chống lại sự biến thiên từ thông đã gây ra nó. Vì vậy phương án $A$ đúng; các phương án còn lại sai.
}
\end{ex}

% ===== Câu 230 =====
% ID: L12C3B16-53-DS
% Mức: VD
\dangbt{Tính suất điện động cảm ứng theo định luật Faraday}
\begin{ex}
% ID: L12C3B16-53-DS
% Mức: VD
Xét Tính suất điện động cảm ứng theo định luật Faraday (Tình huống 3). Hãy xác định tính đúng, sai của bốn phát biểu sau.
\choiceTF
{\True Khi giải cần xác định đúng dữ kiện, phương chiều, điều kiện áp dụng và đơn vị.}
{Kết quả không phụ thuộc dữ kiện và có thể bỏ qua điều kiện.}
{\True Độ lớn suất điện động cảm ứng trung bình của cuộn $N$ vòng là |e| = N|Delta Phi|/Delta t.}
{Suất điện động cảm ứng bằng $N$ Delta t/Delta Phi.}
\loigiai{
\begin{itemchoice}
\item (Đúng) Khi giải cần xác định đúng dữ kiện, phương chiều, điều kiện áp dụng và đơn vị. (Vì): Khi giải cần xác định đúng dữ kiện, phương chiều, điều kiện áp dụng và đơn vị. b) Sai: Kết quả không phụ thuộc dữ kiện và có thể bỏ qua điều kiện. c) Đúng: Độ lớn suất điện động cảm ứng trung bình của cuộn $N$ vòng là |e| = N|Delta Phi|/Delta t. d) Sai: Suất điện động cảm ứng bằng $N$ Delta t/Delta Phi. Kiến thức cốt lõi: Độ lớn suất điện động cảm ứng trung bình của cuộn $N$ vòng là |e| = N|Delta Phi|/Delta t
\item (Sai) Kết quả không phụ thuộc dữ kiện và có thể bỏ qua điều kiện.
\item (Đúng) Độ lớn suất điện động cảm ứng trung bình của cuộn $N$ vòng là |e| = N|Delta Phi|/Delta t.
\item (Sai) Suất điện động cảm ứng bằng $N$ Delta t/Delta Phi.
\end{itemchoice}
}
\end{ex}

% ===== Câu 231 =====
% ID: L12C3B16-53-TL
% Mức: VD
\dangbt{Nhận biết hiện tượng cảm ứng điện từ và điều kiện xuất hiện dòng điện cảm ứng}
\begin{bt}
% ID: L12C3B16-53-TL
% Mức: VD
Nêu một tình huống thực tiễn liên quan đến Nhận biết hiện tượng cảm ứng điện từ và điều kiện xuất hiện dòng điện cảm ứng và giải thích bằng kiến thức Vật lí.
\loigiai{
Cần nêu được: Dòng điện cảm ứng xuất hiện trong mạch kín khi từ thông qua mạch biến thiên. Khi vận dụng phải xác định đúng dữ kiện, phương chiều nếu có, điều kiện áp dụng, hệ đơn vị và kiểm tra tính hợp lí. Nhận định “Chỉ cần đặt mạch kín đứng yên trong từ trường đều không đổi là luôn có dòng điện cảm ứng.” sai.
}
\end{bt}

% ===== Câu 232 =====
% ID: L12C3B16-53-TLN
% Mức: VD
\begin{ex}
% ID: L12C3B16-53-TLN
% Mức: VD
Trong bốn phát biểu về Nhận biết hiện tượng cảm ứng điện từ và điều kiện xuất hiện dòng điện cảm ứng, có bao nhiêu phát biểu đúng? (Chỉ ghi một số nguyên.) (1) Dòng điện cảm ứng xuất hiện trong mạch kín khi từ thông qua mạch biến thiên. (2) Chỉ cần đặt mạch kín đứng yên trong từ trường đều không đổi là luôn có dòng điện cảm ứng. (3) Cần kiểm tra điều kiện áp dụng. (4) Có thể bỏ qua đơn vị.
\shortans{2}
\loigiai{
Đối chiếu với kiến thức: Dòng điện cảm ứng xuất hiện trong mạch kín khi từ thông qua mạch biến thiên. Có 2 phát biểu đúng.
}
\end{ex}

% ===== Câu 233 =====
% ID: L12C3B16-53-TN
% Mức: NB
\dangbt{Xác định chiều dòng điện cảm ứng bằng định luật Lenz}
\begin{ex}
% ID: L12C3B16-53-TN
% Mức: NB
Tình huống 3: chọn kết luận chính xác về Xác định chiều dòng điện cảm ứng bằng định luật Lenz?
\choice
{Kết quả không phụ thuộc dữ kiện và điều kiện áp dụng.}
{\True Dòng điện cảm ứng có chiều sao cho từ trường do nó sinh ra chống lại sự biến thiên từ thông đã gây ra nó.}
{Dòng điện cảm ứng luôn làm tăng biến thiên từ thông ban đầu.}
{Có thể bỏ qua phương, chiều và đơn vị trong mọi trường hợp.}
\loigiai{
Dòng điện cảm ứng có chiều sao cho từ trường do nó sinh ra chống lại sự biến thiên từ thông đã gây ra nó. Vì vậy phương án $B$ đúng; các phương án còn lại sai.
}
\end{ex}

% ===== Câu 234 =====
% ID: L12C3B16-54-DS
% Mức: VD
\dangbt{Tính suất điện động cảm ứng theo định luật Faraday}
\begin{ex}
% ID: L12C3B16-54-DS
% Mức: VD
Xét Tính suất điện động cảm ứng theo định luật Faraday (Tình huống 4). Hãy xác định tính đúng, sai của bốn phát biểu sau.
\choiceTF
{Kết quả không phụ thuộc dữ kiện và có thể bỏ qua điều kiện.}
{\True Khi giải cần xác định đúng dữ kiện, phương chiều, điều kiện áp dụng và đơn vị.}
{Suất điện động cảm ứng bằng $N$ Delta t/Delta Phi.}
{\True Độ lớn suất điện động cảm ứng trung bình của cuộn $N$ vòng là |e| = N|Delta Phi|/Delta t.}
\loigiai{
\begin{itemchoice}
\item (Sai) Kết quả không phụ thuộc dữ kiện và có thể bỏ qua điều kiện. (Vì): Kết quả không phụ thuộc dữ kiện và có thể bỏ qua điều kiện. b) Đúng: Khi giải cần xác định đúng dữ kiện, phương chiều, điều kiện áp dụng và đơn vị. c) Sai: Suất điện động cảm ứng bằng $N$ Delta t/Delta Phi. d) Đúng: Độ lớn suất điện động cảm ứng trung bình của cuộn $N$ vòng là |e| = N|Delta Phi|/Delta t. Kiến thức cốt lõi: Độ lớn suất điện động cảm ứng trung bình của cuộn $N$ vòng là |e| = N|Delta Phi|/Delta t
\item (Đúng) Khi giải cần xác định đúng dữ kiện, phương chiều, điều kiện áp dụng và đơn vị.
\item (Sai) Suất điện động cảm ứng bằng $N$ Delta t/Delta Phi.
\item (Đúng) Độ lớn suất điện động cảm ứng trung bình của cuộn $N$ vòng là |e| = N|Delta Phi|/Delta t.
\end{itemchoice}
}
\end{ex}

% ===== Câu 235 =====
% ID: L12C3B16-54-TL
% Mức: VDC
\dangbt{Nhận biết hiện tượng cảm ứng điện từ và điều kiện xuất hiện dòng điện cảm ứng}
\begin{bt}
% ID: L12C3B16-54-TL
% Mức: VDC
So sánh nhận định đúng “Dòng điện cảm ứng xuất hiện trong mạch kín khi từ thông qua mạch biến thiên.” với nhận định sai “Chỉ cần đặt mạch kín đứng yên trong từ trường đều không đổi là luôn có dòng điện cảm ứng.”; chỉ rõ căn cứ Vật lí.
\loigiai{
Cần nêu được: Dòng điện cảm ứng xuất hiện trong mạch kín khi từ thông qua mạch biến thiên. Khi vận dụng phải xác định đúng dữ kiện, phương chiều nếu có, điều kiện áp dụng, hệ đơn vị và kiểm tra tính hợp lí. Nhận định “Chỉ cần đặt mạch kín đứng yên trong từ trường đều không đổi là luôn có dòng điện cảm ứng.” sai.
}
\end{bt}

% ===== Câu 236 =====
% ID: L12C3B16-54-TLN
% Mức: VDC
\begin{ex}
% ID: L12C3B16-54-TLN
% Mức: VDC
Trong bốn phát biểu về Nhận biết hiện tượng cảm ứng điện từ và điều kiện xuất hiện dòng điện cảm ứng, có bao nhiêu phát biểu đúng? (Chỉ ghi một số nguyên.) (1) Khi giải cần xác định đúng dữ kiện, phương chiều, điều kiện áp dụng và đơn vị. (2) Chỉ cần đặt mạch kín đứng yên trong từ trường đều không đổi là luôn có dòng điện cảm ứng. (3) Phải dùng đúng định luật cho tình huống. (4) Dòng điện cảm ứng xuất hiện trong mạch kín khi từ thông qua mạch biến thiên.
\shortans{3}
\loigiai{
Đối chiếu với kiến thức: Dòng điện cảm ứng xuất hiện trong mạch kín khi từ thông qua mạch biến thiên. Có 3 phát biểu đúng.
}
\end{ex}

% ===== Câu 237 =====
% ID: L12C3B16-54-TN
% Mức: TH
\dangbt{Xác định chiều dòng điện cảm ứng bằng định luật Lenz}
\begin{ex}
% ID: L12C3B16-54-TN
% Mức: TH
Nội dung nào mô tả đúng nhất Xác định chiều dòng điện cảm ứng bằng định luật Lenz?
\choice
{Có thể bỏ qua phương, chiều và đơn vị trong mọi trường hợp.}
{Kết quả không phụ thuộc dữ kiện và điều kiện áp dụng.}
{\True Dòng điện cảm ứng có chiều sao cho từ trường do nó sinh ra chống lại sự biến thiên từ thông đã gây ra nó.}
{Dòng điện cảm ứng luôn làm tăng biến thiên từ thông ban đầu.}
\loigiai{
Dòng điện cảm ứng có chiều sao cho từ trường do nó sinh ra chống lại sự biến thiên từ thông đã gây ra nó. Vì vậy phương án $C$ đúng; các phương án còn lại sai.
}
\end{ex}

% ===== Câu 238 =====
% ID: L12C3B16-55-DS
% Mức: VD
\dangbt{Tính suất điện động cảm ứng theo định luật Faraday}
\begin{ex}
% ID: L12C3B16-55-DS
% Mức: VD
Xét Tính suất điện động cảm ứng theo định luật Faraday (Tình huống 5). Hãy xác định tính đúng, sai của bốn phát biểu sau.
\choiceTF
{\True Độ lớn suất điện động cảm ứng trung bình của cuộn $N$ vòng là |e| = N|Delta Phi|/Delta t.}
{Kết quả không phụ thuộc dữ kiện và có thể bỏ qua điều kiện.}
{Suất điện động cảm ứng bằng $N$ Delta t/Delta Phi.}
{\True Khi giải cần xác định đúng dữ kiện, phương chiều, điều kiện áp dụng và đơn vị.}
\loigiai{
\begin{itemchoice}
\item (Đúng) Độ lớn suất điện động cảm ứng trung bình của cuộn $N$ vòng là |e| = N|Delta Phi|/Delta t. (Vì): ộ lớn suất điện động cảm ứng trung bình của cuộn $N$ vòng là |e| = N|Delta Phi|/Delta t. b) Sai: Kết quả không phụ thuộc dữ kiện và có thể bỏ qua điều kiện. c) Sai: Suất điện động cảm ứng bằng $N$ Delta t/Delta Phi. d) Đúng: Khi giải cần xác định đúng dữ kiện, phương chiều, điều kiện áp dụng và đơn vị. Kiến thức cốt lõi: Độ lớn suất điện động cảm ứng trung bình của cuộn $N$ vòng là |e| = N|Delta Phi|/Delta t
\item (Sai) Kết quả không phụ thuộc dữ kiện và có thể bỏ qua điều kiện.
\item (Sai) Suất điện động cảm ứng bằng $N$ Delta t/Delta Phi.
\item (Đúng) Khi giải cần xác định đúng dữ kiện, phương chiều, điều kiện áp dụng và đơn vị.
\end{itemchoice}
}
\end{ex}

% ===== Câu 239 =====
% ID: L12C3B16-55-TL
% Mức: TH
\dangbt{Xác định chiều dòng điện cảm ứng bằng định luật Lenz}
\begin{bt}
% ID: L12C3B16-55-TL
% Mức: TH
Trình bày kiến thức cốt lõi của Xác định chiều dòng điện cảm ứng bằng định luật Lenz và nêu điều kiện áp dụng.
\loigiai{
Cần nêu được: Dòng điện cảm ứng có chiều sao cho từ trường do nó sinh ra chống lại sự biến thiên từ thông đã gây ra nó. Khi vận dụng phải xác định đúng dữ kiện, phương chiều nếu có, điều kiện áp dụng, hệ đơn vị và kiểm tra tính hợp lí. Nhận định “Dòng điện cảm ứng luôn làm tăng biến thiên từ thông ban đầu.” sai.
}
\end{bt}

% ===== Câu 240 =====
% ID: L12C3B16-55-TLN
% Mức: TH
\begin{ex}
% ID: L12C3B16-55-TLN
% Mức: TH
Trong bốn phát biểu về Xác định chiều dòng điện cảm ứng bằng định luật Lenz, có bao nhiêu phát biểu đúng? (Chỉ ghi một số nguyên.) (1) Dòng điện cảm ứng có chiều sao cho từ trường do nó sinh ra chống lại sự biến thiên từ thông đã gây ra nó. (2) Dòng điện cảm ứng luôn làm tăng biến thiên từ thông ban đầu. (3) Cần kiểm tra điều kiện áp dụng. (4) Có thể bỏ qua đơn vị.
\shortans{2}
\loigiai{
Đối chiếu với kiến thức: Dòng điện cảm ứng có chiều sao cho từ trường do nó sinh ra chống lại sự biến thiên từ thông đã gây ra nó. Có 2 phát biểu đúng.
}
\end{ex}

% ===== Câu 241 =====
% ID: L12C3B16-55-TN
% Mức: TH
\begin{ex}
% ID: L12C3B16-55-TN
% Mức: TH
Khi phân tích Xác định chiều dòng điện cảm ứng bằng định luật Lenz?
\choice
{Dòng điện cảm ứng luôn làm tăng biến thiên từ thông ban đầu.}
{Có thể bỏ qua phương, chiều và đơn vị trong mọi trường hợp.}
{Kết quả không phụ thuộc dữ kiện và điều kiện áp dụng.}
{\True Dòng điện cảm ứng có chiều sao cho từ trường do nó sinh ra chống lại sự biến thiên từ thông đã gây ra nó.}
\loigiai{
Dòng điện cảm ứng có chiều sao cho từ trường do nó sinh ra chống lại sự biến thiên từ thông đã gây ra nó. Vì vậy phương án $D$ đúng; các phương án còn lại sai.
}
\end{ex}

% ===== Câu 242 =====
% ID: L12C3B16-56-DS
% Mức: TH
\dangbt{Khai thác đồ thị từ thông – thời gian và bài toán cảm ứng điện từ tổng hợp}
\begin{ex}
% ID: L12C3B16-56-DS
% Mức: TH
Xét Khai thác đồ thị từ thông – thời gian và bài toán cảm ứng điện từ tổng hợp (Tình huống 1). Hãy xác định tính đúng, sai của bốn phát biểu sau.
\choiceTF
{\True Độ lớn suất điện động cảm ứng bằng độ lớn hệ số góc của đồ thị liên kết từ thông theo thời gian.}
{Đồ thị từ thông nằm ngang ứng với suất điện động cảm ứng khác không.}
{\True Khi giải cần xác định đúng dữ kiện, phương chiều, điều kiện áp dụng và đơn vị.}
{Kết quả không phụ thuộc dữ kiện và có thể bỏ qua điều kiện.}
\loigiai{
\begin{itemchoice}
\item (Đúng) Độ lớn suất điện động cảm ứng bằng độ lớn hệ số góc của đồ thị liên kết từ thông theo thời gian. (Vì): ộ lớn suất điện động cảm ứng bằng độ lớn hệ số góc của đồ thị liên kết từ thông theo thời gian. b) Sai: Đồ thị từ thông nằm ngang ứng với suất điện động cảm ứng khác không. c) Đúng: Khi giải cần xác định đúng dữ kiện, phương chiều, điều kiện áp dụng và đơn vị. d) Sai: Kết quả không phụ thuộc dữ kiện và có thể bỏ qua điều kiện. Kiến thức cốt lõi: Độ lớn suất điện động cảm ứng bằng độ lớn hệ số góc của đồ thị liên kết từ thông theo thời gian
\item (Sai) Đồ thị từ thông nằm ngang ứng với suất điện động cảm ứng khác không.
\item (Đúng) Khi giải cần xác định đúng dữ kiện, phương chiều, điều kiện áp dụng và đơn vị.
\item (Sai) Kết quả không phụ thuộc dữ kiện và có thể bỏ qua điều kiện.
\end{itemchoice}
}
\end{ex}

% ===== Câu 243 =====
% ID: L12C3B16-56-TL
% Mức: TH
\dangbt{Xác định chiều dòng điện cảm ứng bằng định luật Lenz}
\begin{bt}
% ID: L12C3B16-56-TL
% Mức: TH
Giải thích vì sao nhận định sau đúng: “Dòng điện cảm ứng có chiều sao cho từ trường do nó sinh ra chống lại sự biến thiên từ thông đã gây ra nó.”
\loigiai{
Cần nêu được: Dòng điện cảm ứng có chiều sao cho từ trường do nó sinh ra chống lại sự biến thiên từ thông đã gây ra nó. Khi vận dụng phải xác định đúng dữ kiện, phương chiều nếu có, điều kiện áp dụng, hệ đơn vị và kiểm tra tính hợp lí. Nhận định “Dòng điện cảm ứng luôn làm tăng biến thiên từ thông ban đầu.” sai.
}
\end{bt}

% ===== Câu 244 =====
% ID: L12C3B16-56-TLN
% Mức: TH
\begin{ex}
% ID: L12C3B16-56-TLN
% Mức: TH
Trong bốn phát biểu về Xác định chiều dòng điện cảm ứng bằng định luật Lenz, có bao nhiêu phát biểu đúng? (Chỉ ghi một số nguyên.) (1) Dòng điện cảm ứng luôn làm tăng biến thiên từ thông ban đầu. (2) Dòng điện cảm ứng có chiều sao cho từ trường do nó sinh ra chống lại sự biến thiên từ thông đã gây ra nó. (3) Kết quả phải phù hợp ý nghĩa vật lí. (4) Mọi đại lượng đều là vô hướng.
\shortans{1}
\loigiai{
Đối chiếu với kiến thức: Dòng điện cảm ứng có chiều sao cho từ trường do nó sinh ra chống lại sự biến thiên từ thông đã gây ra nó. Có 1 phát biểu đúng.
}
\end{ex}

% ===== Câu 245 =====
% ID: L12C3B16-56-TN
% Mức: TH
\begin{ex}
% ID: L12C3B16-56-TN
% Mức: TH
Một học sinh ôn tập Xác định chiều dòng điện cảm ứng bằng định luật Lenz?
\choice
{\True Dòng điện cảm ứng có chiều sao cho từ trường do nó sinh ra chống lại sự biến thiên từ thông đã gây ra nó.}
{Dòng điện cảm ứng luôn làm tăng biến thiên từ thông ban đầu.}
{Có thể bỏ qua phương, chiều và đơn vị trong mọi trường hợp.}
{Kết quả không phụ thuộc dữ kiện và điều kiện áp dụng.}
\loigiai{
Dòng điện cảm ứng có chiều sao cho từ trường do nó sinh ra chống lại sự biến thiên từ thông đã gây ra nó. Vì vậy phương án $A$ đúng; các phương án còn lại sai.
}
\end{ex}

% ===== Câu 246 =====
% ID: L12C3B16-57-DS
% Mức: TH
\dangbt{Khai thác đồ thị từ thông – thời gian và bài toán cảm ứng điện từ tổng hợp}
\begin{ex}
% ID: L12C3B16-57-DS
% Mức: TH
Xét Khai thác đồ thị từ thông – thời gian và bài toán cảm ứng điện từ tổng hợp (Tình huống 2). Hãy xác định tính đúng, sai của bốn phát biểu sau.
\choiceTF
{Đồ thị từ thông nằm ngang ứng với suất điện động cảm ứng khác không.}
{\True Độ lớn suất điện động cảm ứng bằng độ lớn hệ số góc của đồ thị liên kết từ thông theo thời gian.}
{Kết quả không phụ thuộc dữ kiện và có thể bỏ qua điều kiện.}
{\True Khi giải cần xác định đúng dữ kiện, phương chiều, điều kiện áp dụng và đơn vị.}
\loigiai{
\begin{itemchoice}
\item (Sai) Đồ thị từ thông nằm ngang ứng với suất điện động cảm ứng khác không. (Vì): Đồ thị từ thông nằm ngang ứng với suất điện động cảm ứng khác không. b) Đúng: Độ lớn suất điện động cảm ứng bằng độ lớn hệ số góc của đồ thị liên kết từ thông theo thời gian. c) Sai: Kết quả không phụ thuộc dữ kiện và có thể bỏ qua điều kiện. d) Đúng: Khi giải cần xác định đúng dữ kiện, phương chiều, điều kiện áp dụng và đơn vị. Kiến thức cốt lõi: Độ lớn suất điện động cảm ứng bằng độ lớn hệ số góc của đồ thị liên kết từ thông theo thời gian
\item (Đúng) Độ lớn suất điện động cảm ứng bằng độ lớn hệ số góc của đồ thị liên kết từ thông theo thời gian.
\item (Sai) Kết quả không phụ thuộc dữ kiện và có thể bỏ qua điều kiện.
\item (Đúng) Khi giải cần xác định đúng dữ kiện, phương chiều, điều kiện áp dụng và đơn vị.
\end{itemchoice}
}
\end{ex}

% ===== Câu 247 =====
% ID: L12C3B16-57-TL
% Mức: VD
\dangbt{Xác định chiều dòng điện cảm ứng bằng định luật Lenz}
\begin{bt}
% ID: L12C3B16-57-TL
% Mức: VD
Phân tích sai lầm trong nhận định: “Dòng điện cảm ứng luôn làm tăng biến thiên từ thông ban đầu.”
\loigiai{
Cần nêu được: Dòng điện cảm ứng có chiều sao cho từ trường do nó sinh ra chống lại sự biến thiên từ thông đã gây ra nó. Khi vận dụng phải xác định đúng dữ kiện, phương chiều nếu có, điều kiện áp dụng, hệ đơn vị và kiểm tra tính hợp lí. Nhận định “Dòng điện cảm ứng luôn làm tăng biến thiên từ thông ban đầu.” sai.
}
\end{bt}

% ===== Câu 248 =====
% ID: L12C3B16-57-TLN
% Mức: TH
\begin{ex}
% ID: L12C3B16-57-TLN
% Mức: TH
Trong bốn phát biểu về Xác định chiều dòng điện cảm ứng bằng định luật Lenz, có bao nhiêu phát biểu đúng? (Chỉ ghi một số nguyên.) (1) Dòng điện cảm ứng có chiều sao cho từ trường do nó sinh ra chống lại sự biến thiên từ thông đã gây ra nó. (2) Dòng điện cảm ứng luôn làm tăng biến thiên từ thông ban đầu. (3) Cần kiểm tra điều kiện áp dụng. (4) Có thể bỏ qua đơn vị.
\shortans{2}
\loigiai{
Đối chiếu với kiến thức: Dòng điện cảm ứng có chiều sao cho từ trường do nó sinh ra chống lại sự biến thiên từ thông đã gây ra nó. Có 2 phát biểu đúng.
}
\end{ex}

% ===== Câu 249 =====
% ID: L12C3B16-57-TN
% Mức: TH
\begin{ex}
% ID: L12C3B16-57-TN
% Mức: TH
Chọn phát biểu không trái với kiến thức về Xác định chiều dòng điện cảm ứng bằng định luật Lenz?
\choice
{Kết quả không phụ thuộc dữ kiện và điều kiện áp dụng.}
{\True Dòng điện cảm ứng có chiều sao cho từ trường do nó sinh ra chống lại sự biến thiên từ thông đã gây ra nó.}
{Dòng điện cảm ứng luôn làm tăng biến thiên từ thông ban đầu.}
{Có thể bỏ qua phương, chiều và đơn vị trong mọi trường hợp.}
\loigiai{
Dòng điện cảm ứng có chiều sao cho từ trường do nó sinh ra chống lại sự biến thiên từ thông đã gây ra nó. Vì vậy phương án $B$ đúng; các phương án còn lại sai.
}
\end{ex}

% ===== Câu 250 =====
% ID: L12C3B16-58-DS
% Mức: VD
\dangbt{Khai thác đồ thị từ thông – thời gian và bài toán cảm ứng điện từ tổng hợp}
\begin{ex}
% ID: L12C3B16-58-DS
% Mức: VD
Xét Khai thác đồ thị từ thông – thời gian và bài toán cảm ứng điện từ tổng hợp (Tình huống 3). Hãy xác định tính đúng, sai của bốn phát biểu sau.
\choiceTF
{\True Khi giải cần xác định đúng dữ kiện, phương chiều, điều kiện áp dụng và đơn vị.}
{Kết quả không phụ thuộc dữ kiện và có thể bỏ qua điều kiện.}
{\True Độ lớn suất điện động cảm ứng bằng độ lớn hệ số góc của đồ thị liên kết từ thông theo thời gian.}
{Đồ thị từ thông nằm ngang ứng với suất điện động cảm ứng khác không.}
\loigiai{
\begin{itemchoice}
\item (Đúng) Khi giải cần xác định đúng dữ kiện, phương chiều, điều kiện áp dụng và đơn vị. (Vì): Khi giải cần xác định đúng dữ kiện, phương chiều, điều kiện áp dụng và đơn vị. b) Sai: Kết quả không phụ thuộc dữ kiện và có thể bỏ qua điều kiện. c) Đúng: Độ lớn suất điện động cảm ứng bằng độ lớn hệ số góc của đồ thị liên kết từ thông theo thời gian. d) Sai: Đồ thị từ thông nằm ngang ứng với suất điện động cảm ứng khác không. Kiến thức cốt lõi: Độ lớn suất điện động cảm ứng bằng độ lớn hệ số góc của đồ thị liên kết từ thông theo thời gian
\item (Sai) Kết quả không phụ thuộc dữ kiện và có thể bỏ qua điều kiện.
\item (Đúng) Độ lớn suất điện động cảm ứng bằng độ lớn hệ số góc của đồ thị liên kết từ thông theo thời gian.
\item (Sai) Đồ thị từ thông nằm ngang ứng với suất điện động cảm ứng khác không.
\end{itemchoice}
}
\end{ex}

% ===== Câu 251 =====
% ID: L12C3B16-58-TL
% Mức: VD
\dangbt{Xác định chiều dòng điện cảm ứng bằng định luật Lenz}
\begin{bt}
% ID: L12C3B16-58-TL
% Mức: VD
Đề xuất các bước giải một bài toán thuộc dạng Xác định chiều dòng điện cảm ứng bằng định luật Lenz.
\loigiai{
Cần nêu được: Dòng điện cảm ứng có chiều sao cho từ trường do nó sinh ra chống lại sự biến thiên từ thông đã gây ra nó. Khi vận dụng phải xác định đúng dữ kiện, phương chiều nếu có, điều kiện áp dụng, hệ đơn vị và kiểm tra tính hợp lí. Nhận định “Dòng điện cảm ứng luôn làm tăng biến thiên từ thông ban đầu.” sai.
}
\end{bt}

% ===== Câu 252 =====
% ID: L12C3B16-58-TLN
% Mức: VD
\begin{ex}
% ID: L12C3B16-58-TLN
% Mức: VD
Trong bốn phát biểu về Xác định chiều dòng điện cảm ứng bằng định luật Lenz, có bao nhiêu phát biểu đúng? (Chỉ ghi một số nguyên.) (1) Dòng điện cảm ứng luôn làm tăng biến thiên từ thông ban đầu. (2) Dòng điện cảm ứng có chiều sao cho từ trường do nó sinh ra chống lại sự biến thiên từ thông đã gây ra nó. (3) Kết quả phải phù hợp ý nghĩa vật lí. (4) Mọi đại lượng đều là vô hướng.
\shortans{2}
\loigiai{
Đối chiếu với kiến thức: Dòng điện cảm ứng có chiều sao cho từ trường do nó sinh ra chống lại sự biến thiên từ thông đã gây ra nó. Có 2 phát biểu đúng.
}
\end{ex}

% ===== Câu 253 =====
% ID: L12C3B16-58-TN
% Mức: VD
\begin{ex}
% ID: L12C3B16-58-TN
% Mức: VD
Cơ sở Vật lí đúng dùng để giải Xác định chiều dòng điện cảm ứng bằng định luật Lenz?
\choice
{Có thể bỏ qua phương, chiều và đơn vị trong mọi trường hợp.}
{Kết quả không phụ thuộc dữ kiện và điều kiện áp dụng.}
{\True Dòng điện cảm ứng có chiều sao cho từ trường do nó sinh ra chống lại sự biến thiên từ thông đã gây ra nó.}
{Dòng điện cảm ứng luôn làm tăng biến thiên từ thông ban đầu.}
\loigiai{
Dòng điện cảm ứng có chiều sao cho từ trường do nó sinh ra chống lại sự biến thiên từ thông đã gây ra nó. Vì vậy phương án $C$ đúng; các phương án còn lại sai.
}
\end{ex}

% ===== Câu 254 =====
% ID: L12C3B16-59-DS
% Mức: VD
\dangbt{Khai thác đồ thị từ thông – thời gian và bài toán cảm ứng điện từ tổng hợp}
\begin{ex}
% ID: L12C3B16-59-DS
% Mức: VD
Xét Khai thác đồ thị từ thông – thời gian và bài toán cảm ứng điện từ tổng hợp (Tình huống 4). Hãy xác định tính đúng, sai của bốn phát biểu sau.
\choiceTF
{Kết quả không phụ thuộc dữ kiện và có thể bỏ qua điều kiện.}
{\True Khi giải cần xác định đúng dữ kiện, phương chiều, điều kiện áp dụng và đơn vị.}
{Đồ thị từ thông nằm ngang ứng với suất điện động cảm ứng khác không.}
{\True Độ lớn suất điện động cảm ứng bằng độ lớn hệ số góc của đồ thị liên kết từ thông theo thời gian.}
\loigiai{
\begin{itemchoice}
\item (Sai) Kết quả không phụ thuộc dữ kiện và có thể bỏ qua điều kiện. (Vì): Kết quả không phụ thuộc dữ kiện và có thể bỏ qua điều kiện. b) Đúng: Khi giải cần xác định đúng dữ kiện, phương chiều, điều kiện áp dụng và đơn vị. c) Sai: Đồ thị từ thông nằm ngang ứng với suất điện động cảm ứng khác không. d) Đúng: Độ lớn suất điện động cảm ứng bằng độ lớn hệ số góc của đồ thị liên kết từ thông theo thời gian. Kiến thức cốt lõi: Độ lớn suất điện động cảm ứng bằng độ lớn hệ số góc của đồ thị liên kết từ thông theo thời gian
\item (Đúng) Khi giải cần xác định đúng dữ kiện, phương chiều, điều kiện áp dụng và đơn vị.
\item (Sai) Đồ thị từ thông nằm ngang ứng với suất điện động cảm ứng khác không.
\item (Đúng) Độ lớn suất điện động cảm ứng bằng độ lớn hệ số góc của đồ thị liên kết từ thông theo thời gian.
\end{itemchoice}
}
\end{ex}

% ===== Câu 255 =====
% ID: L12C3B16-59-TL
% Mức: VD
\dangbt{Xác định chiều dòng điện cảm ứng bằng định luật Lenz}
\begin{bt}
% ID: L12C3B16-59-TL
% Mức: VD
Nêu một tình huống thực tiễn liên quan đến Xác định chiều dòng điện cảm ứng bằng định luật Lenz và giải thích bằng kiến thức Vật lí.
\loigiai{
Cần nêu được: Dòng điện cảm ứng có chiều sao cho từ trường do nó sinh ra chống lại sự biến thiên từ thông đã gây ra nó. Khi vận dụng phải xác định đúng dữ kiện, phương chiều nếu có, điều kiện áp dụng, hệ đơn vị và kiểm tra tính hợp lí. Nhận định “Dòng điện cảm ứng luôn làm tăng biến thiên từ thông ban đầu.” sai.
}
\end{bt}

% ===== Câu 256 =====
% ID: L12C3B16-59-TLN
% Mức: VD
\begin{ex}
% ID: L12C3B16-59-TLN
% Mức: VD
Trong bốn phát biểu về Xác định chiều dòng điện cảm ứng bằng định luật Lenz, có bao nhiêu phát biểu đúng? (Chỉ ghi một số nguyên.) (1) Dòng điện cảm ứng có chiều sao cho từ trường do nó sinh ra chống lại sự biến thiên từ thông đã gây ra nó. (2) Dòng điện cảm ứng luôn làm tăng biến thiên từ thông ban đầu. (3) Cần kiểm tra điều kiện áp dụng. (4) Có thể bỏ qua đơn vị.
\shortans{2}
\loigiai{
Đối chiếu với kiến thức: Dòng điện cảm ứng có chiều sao cho từ trường do nó sinh ra chống lại sự biến thiên từ thông đã gây ra nó. Có 2 phát biểu đúng.
}
\end{ex}

% ===== Câu 257 =====
% ID: L12C3B16-59-TN
% Mức: VD
\begin{ex}
% ID: L12C3B16-59-TN
% Mức: VD
Trong một bài thực hành về Xác định chiều dòng điện cảm ứng bằng định luật Lenz?
\choice
{Dòng điện cảm ứng luôn làm tăng biến thiên từ thông ban đầu.}
{Có thể bỏ qua phương, chiều và đơn vị trong mọi trường hợp.}
{Kết quả không phụ thuộc dữ kiện và điều kiện áp dụng.}
{\True Dòng điện cảm ứng có chiều sao cho từ trường do nó sinh ra chống lại sự biến thiên từ thông đã gây ra nó.}
\loigiai{
Dòng điện cảm ứng có chiều sao cho từ trường do nó sinh ra chống lại sự biến thiên từ thông đã gây ra nó. Vì vậy phương án $D$ đúng; các phương án còn lại sai.
}
\end{ex}

% ===== Câu 258 =====
% ID: L12C3B16-60-DS
% Mức: VD
\dangbt{Khai thác đồ thị từ thông – thời gian và bài toán cảm ứng điện từ tổng hợp}
\begin{ex}
% ID: L12C3B16-60-DS
% Mức: VD
Xét Khai thác đồ thị từ thông – thời gian và bài toán cảm ứng điện từ tổng hợp (Tình huống 5). Hãy xác định tính đúng, sai của bốn phát biểu sau.
\choiceTF
{\True Độ lớn suất điện động cảm ứng bằng độ lớn hệ số góc của đồ thị liên kết từ thông theo thời gian.}
{Kết quả không phụ thuộc dữ kiện và có thể bỏ qua điều kiện.}
{Đồ thị từ thông nằm ngang ứng với suất điện động cảm ứng khác không.}
{\True Khi giải cần xác định đúng dữ kiện, phương chiều, điều kiện áp dụng và đơn vị.}
\loigiai{
\begin{itemchoice}
\item (Đúng) Độ lớn suất điện động cảm ứng bằng độ lớn hệ số góc của đồ thị liên kết từ thông theo thời gian. (Vì): ộ lớn suất điện động cảm ứng bằng độ lớn hệ số góc của đồ thị liên kết từ thông theo thời gian. b) Sai: Kết quả không phụ thuộc dữ kiện và có thể bỏ qua điều kiện. c) Sai: Đồ thị từ thông nằm ngang ứng với suất điện động cảm ứng khác không. d) Đúng: Khi giải cần xác định đúng dữ kiện, phương chiều, điều kiện áp dụng và đơn vị. Kiến thức cốt lõi: Độ lớn suất điện động cảm ứng bằng độ lớn hệ số góc của đồ thị liên kết từ thông theo thời gian
\item (Sai) Kết quả không phụ thuộc dữ kiện và có thể bỏ qua điều kiện.
\item (Sai) Đồ thị từ thông nằm ngang ứng với suất điện động cảm ứng khác không.
\item (Đúng) Khi giải cần xác định đúng dữ kiện, phương chiều, điều kiện áp dụng và đơn vị.
\end{itemchoice}
}
\end{ex}

% ===== Câu 259 =====
% ID: L12C3B16-60-TL
% Mức: VDC
\dangbt{Xác định chiều dòng điện cảm ứng bằng định luật Lenz}
\begin{bt}
% ID: L12C3B16-60-TL
% Mức: VDC
So sánh nhận định đúng “Dòng điện cảm ứng có chiều sao cho từ trường do nó sinh ra chống lại sự biến thiên từ thông đã gây ra nó.” với nhận định sai “Dòng điện cảm ứng luôn làm tăng biến thiên từ thông ban đầu.”; chỉ rõ căn cứ Vật lí.
\loigiai{
Cần nêu được: Dòng điện cảm ứng có chiều sao cho từ trường do nó sinh ra chống lại sự biến thiên từ thông đã gây ra nó. Khi vận dụng phải xác định đúng dữ kiện, phương chiều nếu có, điều kiện áp dụng, hệ đơn vị và kiểm tra tính hợp lí. Nhận định “Dòng điện cảm ứng luôn làm tăng biến thiên từ thông ban đầu.” sai.
}
\end{bt}

% ===== Câu 260 =====
% ID: L12C3B16-60-TLN
% Mức: VDC
\begin{ex}
% ID: L12C3B16-60-TLN
% Mức: VDC
Trong bốn phát biểu về Xác định chiều dòng điện cảm ứng bằng định luật Lenz, có bao nhiêu phát biểu đúng? (Chỉ ghi một số nguyên.) (1) Khi giải cần xác định đúng dữ kiện, phương chiều, điều kiện áp dụng và đơn vị. (2) Dòng điện cảm ứng luôn làm tăng biến thiên từ thông ban đầu. (3) Phải dùng đúng định luật cho tình huống. (4) Dòng điện cảm ứng có chiều sao cho từ trường do nó sinh ra chống lại sự biến thiên từ thông đã gây ra nó.
\shortans{3}
\loigiai{
Đối chiếu với kiến thức: Dòng điện cảm ứng có chiều sao cho từ trường do nó sinh ra chống lại sự biến thiên từ thông đã gây ra nó. Có 3 phát biểu đúng.
}
\end{ex}

% ===== Câu 261 =====
% ID: L12C3B16-60-TN
% Mức: VD
\begin{ex}
% ID: L12C3B16-60-TN
% Mức: VD
Kết luận đúng khi vận dụng Xác định chiều dòng điện cảm ứng bằng định luật Lenz?
\choice
{\True Dòng điện cảm ứng có chiều sao cho từ trường do nó sinh ra chống lại sự biến thiên từ thông đã gây ra nó.}
{Dòng điện cảm ứng luôn làm tăng biến thiên từ thông ban đầu.}
{Có thể bỏ qua phương, chiều và đơn vị trong mọi trường hợp.}
{Kết quả không phụ thuộc dữ kiện và điều kiện áp dụng.}
\loigiai{
Dòng điện cảm ứng có chiều sao cho từ trường do nó sinh ra chống lại sự biến thiên từ thông đã gây ra nó. Vì vậy phương án $A$ đúng; các phương án còn lại sai.
}
\end{ex}

% ===== Câu 262 =====
% ID: L12C3B16-61-TL
% Mức: TH
\dangbt{Tính suất điện động cảm ứng theo định luật Faraday}
\begin{bt}
% ID: L12C3B16-61-TL
% Mức: TH
Trình bày kiến thức cốt lõi của Tính suất điện động cảm ứng theo định luật Faraday và nêu điều kiện áp dụng.
\loigiai{
Cần nêu được: Độ lớn suất điện động cảm ứng trung bình của cuộn $N$ vòng là |e| = N|Delta Phi|/Delta t. Khi vận dụng phải xác định đúng dữ kiện, phương chiều nếu có, điều kiện áp dụng, hệ đơn vị và kiểm tra tính hợp lí. Nhận định “Suất điện động cảm ứng bằng $N$ Delta t/Delta Phi.” sai.
}
\end{bt}

% ===== Câu 263 =====
% ID: L12C3B16-61-TLN
% Mức: TH
\begin{ex}
% ID: L12C3B16-61-TLN
% Mức: TH
Trong bốn phát biểu về Tính suất điện động cảm ứng theo định luật Faraday, có bao nhiêu phát biểu đúng? (Chỉ ghi một số nguyên.) (1) Độ lớn suất điện động cảm ứng trung bình của cuộn $N$ vòng là |e| = N|Delta Phi|/Delta t. (2) Suất điện động cảm ứng bằng $N$ Delta t/Delta Phi. (3) Cần kiểm tra điều kiện áp dụng. (4) Có thể bỏ qua đơn vị.
\shortans{2}
\loigiai{
Đối chiếu với kiến thức: Độ lớn suất điện động cảm ứng trung bình của cuộn $N$ vòng là |e| = N|Delta Phi|/Delta t. Có 2 phát biểu đúng.
}
\end{ex}

% ===== Câu 264 =====
% ID: L12C3B16-61-TN
% Mức: NB
\dangbt{Nhận biết từ thông và tính từ thông qua một diện tích}
\begin{ex}
% ID: L12C3B16-61-TN
% Mức: NB
Phát biểu nào sau đây đúng về Nhận biết từ thông và tính từ thông qua một diện tích?
\choice
{Từ thông luôn bằng BS sin(alpha), với alpha là góc giữa $B$ và pháp tuyến.}
{Có thể bỏ qua phương, chiều và đơn vị trong mọi trường hợp.}
{Kết quả không phụ thuộc dữ kiện và điều kiện áp dụng.}
{\True Từ thông qua diện tích $S$ trong từ trường đều là Phi = BS cos(alpha), với alpha là góc giữa $B$ và pháp tuyến.}
\loigiai{
Từ thông qua diện tích $S$ trong từ trường đều là Phi = BS cos(alpha), với alpha là góc giữa $B$ và pháp tuyến. Vì vậy phương án $D$ đúng; các phương án còn lại sai.
}
\end{ex}

% ===== Câu 265 =====
% ID: L12C3B16-62-TL
% Mức: TH
\dangbt{Tính suất điện động cảm ứng theo định luật Faraday}
\begin{bt}
% ID: L12C3B16-62-TL
% Mức: TH
Giải thích vì sao nhận định sau đúng: “Độ lớn suất điện động cảm ứng trung bình của cuộn $N$ vòng là |e| = N|Delta Phi|/Delta t.”
\loigiai{
Cần nêu được: Độ lớn suất điện động cảm ứng trung bình của cuộn $N$ vòng là |e| = N|Delta Phi|/Delta t. Khi vận dụng phải xác định đúng dữ kiện, phương chiều nếu có, điều kiện áp dụng, hệ đơn vị và kiểm tra tính hợp lí. Nhận định “Suất điện động cảm ứng bằng $N$ Delta t/Delta Phi.” sai.
}
\end{bt}

% ===== Câu 266 =====
% ID: L12C3B16-62-TLN
% Mức: TH
\begin{ex}
% ID: L12C3B16-62-TLN
% Mức: TH
Trong bốn phát biểu về Tính suất điện động cảm ứng theo định luật Faraday, có bao nhiêu phát biểu đúng? (Chỉ ghi một số nguyên.) (1) Suất điện động cảm ứng bằng $N$ Delta t/Delta Phi. (2) Độ lớn suất điện động cảm ứng trung bình của cuộn $N$ vòng là |e| = N|Delta Phi|/Delta t. (3) Kết quả phải phù hợp ý nghĩa vật lí. (4) Mọi đại lượng đều là vô hướng.
\shortans{1}
\loigiai{
Đối chiếu với kiến thức: Độ lớn suất điện động cảm ứng trung bình của cuộn $N$ vòng là |e| = N|Delta Phi|/Delta t. Có 1 phát biểu đúng.
}
\end{ex}

% ===== Câu 267 =====
% ID: L12C3B16-62-TN
% Mức: NB
\dangbt{Nhận biết từ thông và tính từ thông qua một diện tích}
\begin{ex}
% ID: L12C3B16-62-TN
% Mức: NB
Trong nội dung Nhận biết từ thông và tính từ thông qua một diện tích?
\choice
{\True Từ thông qua diện tích $S$ trong từ trường đều là Phi = BS cos(alpha), với alpha là góc giữa $B$ và pháp tuyến.}
{Từ thông luôn bằng BS sin(alpha), với alpha là góc giữa $B$ và pháp tuyến.}
{Có thể bỏ qua phương, chiều và đơn vị trong mọi trường hợp.}
{Kết quả không phụ thuộc dữ kiện và điều kiện áp dụng.}
\loigiai{
Từ thông qua diện tích $S$ trong từ trường đều là Phi = BS cos(alpha), với alpha là góc giữa $B$ và pháp tuyến. Vì vậy phương án $A$ đúng; các phương án còn lại sai.
}
\end{ex}

% ===== Câu 268 =====
% ID: L12C3B16-63-TL
% Mức: VD
\dangbt{Tính suất điện động cảm ứng theo định luật Faraday}
\begin{bt}
% ID: L12C3B16-63-TL
% Mức: VD
Phân tích sai lầm trong nhận định: “Suất điện động cảm ứng bằng $N$ Delta t/Delta Phi.”
\loigiai{
Cần nêu được: Độ lớn suất điện động cảm ứng trung bình của cuộn $N$ vòng là |e| = N|Delta Phi|/Delta t. Khi vận dụng phải xác định đúng dữ kiện, phương chiều nếu có, điều kiện áp dụng, hệ đơn vị và kiểm tra tính hợp lí. Nhận định “Suất điện động cảm ứng bằng $N$ Delta t/Delta Phi.” sai.
}
\end{bt}

% ===== Câu 269 =====
% ID: L12C3B16-63-TLN
% Mức: TH
\begin{ex}
% ID: L12C3B16-63-TLN
% Mức: TH
Trong bốn phát biểu về Tính suất điện động cảm ứng theo định luật Faraday, có bao nhiêu phát biểu đúng? (Chỉ ghi một số nguyên.) (1) Độ lớn suất điện động cảm ứng trung bình của cuộn $N$ vòng là |e| = N|Delta Phi|/Delta t. (2) Suất điện động cảm ứng bằng $N$ Delta t/Delta Phi. (3) Cần kiểm tra điều kiện áp dụng. (4) Có thể bỏ qua đơn vị.
\shortans{2}
\loigiai{
Đối chiếu với kiến thức: Độ lớn suất điện động cảm ứng trung bình của cuộn $N$ vòng là |e| = N|Delta Phi|/Delta t. Có 2 phát biểu đúng.
}
\end{ex}

% ===== Câu 270 =====
% ID: L12C3B16-63-TN
% Mức: NB
\dangbt{Nhận biết từ thông và tính từ thông qua một diện tích}
\begin{ex}
% ID: L12C3B16-63-TN
% Mức: NB
Tình huống 3: chọn kết luận chính xác về Nhận biết từ thông và tính từ thông qua một diện tích?
\choice
{Kết quả không phụ thuộc dữ kiện và điều kiện áp dụng.}
{\True Từ thông qua diện tích $S$ trong từ trường đều là Phi = BS cos(alpha), với alpha là góc giữa $B$ và pháp tuyến.}
{Từ thông luôn bằng BS sin(alpha), với alpha là góc giữa $B$ và pháp tuyến.}
{Có thể bỏ qua phương, chiều và đơn vị trong mọi trường hợp.}
\loigiai{
Từ thông qua diện tích $S$ trong từ trường đều là Phi = BS cos(alpha), với alpha là góc giữa $B$ và pháp tuyến. Vì vậy phương án $B$ đúng; các phương án còn lại sai.
}
\end{ex}

% ===== Câu 271 =====
% ID: L12C3B16-64-TL
% Mức: VD
\dangbt{Tính suất điện động cảm ứng theo định luật Faraday}
\begin{bt}
% ID: L12C3B16-64-TL
% Mức: VD
Đề xuất các bước giải một bài toán thuộc dạng Tính suất điện động cảm ứng theo định luật Faraday.
\loigiai{
Cần nêu được: Độ lớn suất điện động cảm ứng trung bình của cuộn $N$ vòng là |e| = N|Delta Phi|/Delta t. Khi vận dụng phải xác định đúng dữ kiện, phương chiều nếu có, điều kiện áp dụng, hệ đơn vị và kiểm tra tính hợp lí. Nhận định “Suất điện động cảm ứng bằng $N$ Delta t/Delta Phi.” sai.
}
\end{bt}

% ===== Câu 272 =====
% ID: L12C3B16-64-TLN
% Mức: VD
\begin{ex}
% ID: L12C3B16-64-TLN
% Mức: VD
Trong bốn phát biểu về Tính suất điện động cảm ứng theo định luật Faraday, có bao nhiêu phát biểu đúng? (Chỉ ghi một số nguyên.) (1) Suất điện động cảm ứng bằng $N$ Delta t/Delta Phi. (2) Độ lớn suất điện động cảm ứng trung bình của cuộn $N$ vòng là |e| = N|Delta Phi|/Delta t. (3) Kết quả phải phù hợp ý nghĩa vật lí. (4) Mọi đại lượng đều là vô hướng.
\shortans{2}
\loigiai{
Đối chiếu với kiến thức: Độ lớn suất điện động cảm ứng trung bình của cuộn $N$ vòng là |e| = N|Delta Phi|/Delta t. Có 2 phát biểu đúng.
}
\end{ex}

% ===== Câu 273 =====
% ID: L12C3B16-64-TN
% Mức: TH
\dangbt{Nhận biết từ thông và tính từ thông qua một diện tích}
\begin{ex}
% ID: L12C3B16-64-TN
% Mức: TH
Nội dung nào mô tả đúng nhất Nhận biết từ thông và tính từ thông qua một diện tích?
\choice
{Có thể bỏ qua phương, chiều và đơn vị trong mọi trường hợp.}
{Kết quả không phụ thuộc dữ kiện và điều kiện áp dụng.}
{\True Từ thông qua diện tích $S$ trong từ trường đều là Phi = BS cos(alpha), với alpha là góc giữa $B$ và pháp tuyến.}
{Từ thông luôn bằng BS sin(alpha), với alpha là góc giữa $B$ và pháp tuyến.}
\loigiai{
Từ thông qua diện tích $S$ trong từ trường đều là Phi = BS cos(alpha), với alpha là góc giữa $B$ và pháp tuyến. Vì vậy phương án $C$ đúng; các phương án còn lại sai.
}
\end{ex}

% ===== Câu 274 =====
% ID: L12C3B16-65-TL
% Mức: VD
\dangbt{Tính suất điện động cảm ứng theo định luật Faraday}
\begin{bt}
% ID: L12C3B16-65-TL
% Mức: VD
Nêu một tình huống thực tiễn liên quan đến Tính suất điện động cảm ứng theo định luật Faraday và giải thích bằng kiến thức Vật lí.
\loigiai{
Cần nêu được: Độ lớn suất điện động cảm ứng trung bình của cuộn $N$ vòng là |e| = N|Delta Phi|/Delta t. Khi vận dụng phải xác định đúng dữ kiện, phương chiều nếu có, điều kiện áp dụng, hệ đơn vị và kiểm tra tính hợp lí. Nhận định “Suất điện động cảm ứng bằng $N$ Delta t/Delta Phi.” sai.
}
\end{bt}

% ===== Câu 275 =====
% ID: L12C3B16-65-TLN
% Mức: VD
\begin{ex}
% ID: L12C3B16-65-TLN
% Mức: VD
Trong bốn phát biểu về Tính suất điện động cảm ứng theo định luật Faraday, có bao nhiêu phát biểu đúng? (Chỉ ghi một số nguyên.) (1) Độ lớn suất điện động cảm ứng trung bình của cuộn $N$ vòng là |e| = N|Delta Phi|/Delta t. (2) Suất điện động cảm ứng bằng $N$ Delta t/Delta Phi. (3) Cần kiểm tra điều kiện áp dụng. (4) Có thể bỏ qua đơn vị.
\shortans{2}
\loigiai{
Đối chiếu với kiến thức: Độ lớn suất điện động cảm ứng trung bình của cuộn $N$ vòng là |e| = N|Delta Phi|/Delta t. Có 2 phát biểu đúng.
}
\end{ex}

% ===== Câu 276 =====
% ID: L12C3B16-65-TN
% Mức: TH
\dangbt{Nhận biết từ thông và tính từ thông qua một diện tích}
\begin{ex}
% ID: L12C3B16-65-TN
% Mức: TH
Khi phân tích Nhận biết từ thông và tính từ thông qua một diện tích?
\choice
{Từ thông luôn bằng BS sin(alpha), với alpha là góc giữa $B$ và pháp tuyến.}
{Có thể bỏ qua phương, chiều và đơn vị trong mọi trường hợp.}
{Kết quả không phụ thuộc dữ kiện và điều kiện áp dụng.}
{\True Từ thông qua diện tích $S$ trong từ trường đều là Phi = BS cos(alpha), với alpha là góc giữa $B$ và pháp tuyến.}
\loigiai{
Từ thông qua diện tích $S$ trong từ trường đều là Phi = BS cos(alpha), với alpha là góc giữa $B$ và pháp tuyến. Vì vậy phương án $D$ đúng; các phương án còn lại sai.
}
\end{ex}

% ===== Câu 277 =====
% ID: L12C3B16-66-TL
% Mức: VDC
\dangbt{Tính suất điện động cảm ứng theo định luật Faraday}
\begin{bt}
% ID: L12C3B16-66-TL
% Mức: VDC
So sánh nhận định đúng “Độ lớn suất điện động cảm ứng trung bình của cuộn $N$ vòng là |e| = N|Delta Phi|/Delta t.” với nhận định sai “Suất điện động cảm ứng bằng $N$ Delta t/Delta Phi.”; chỉ rõ căn cứ Vật lí.
\loigiai{
Cần nêu được: Độ lớn suất điện động cảm ứng trung bình của cuộn $N$ vòng là |e| = N|Delta Phi|/Delta t. Khi vận dụng phải xác định đúng dữ kiện, phương chiều nếu có, điều kiện áp dụng, hệ đơn vị và kiểm tra tính hợp lí. Nhận định “Suất điện động cảm ứng bằng $N$ Delta t/Delta Phi.” sai.
}
\end{bt}

% ===== Câu 278 =====
% ID: L12C3B16-66-TLN
% Mức: VDC
\begin{ex}
% ID: L12C3B16-66-TLN
% Mức: VDC
Trong bốn phát biểu về Tính suất điện động cảm ứng theo định luật Faraday, có bao nhiêu phát biểu đúng? (Chỉ ghi một số nguyên.) (1) Khi giải cần xác định đúng dữ kiện, phương chiều, điều kiện áp dụng và đơn vị. (2) Suất điện động cảm ứng bằng $N$ Delta t/Delta Phi. (3) Phải dùng đúng định luật cho tình huống. (4) Độ lớn suất điện động cảm ứng trung bình của cuộn $N$ vòng là |e| = N|Delta Phi|/Delta t.
\shortans{3}
\loigiai{
Đối chiếu với kiến thức: Độ lớn suất điện động cảm ứng trung bình của cuộn $N$ vòng là |e| = N|Delta Phi|/Delta t. Có 3 phát biểu đúng.
}
\end{ex}

% ===== Câu 279 =====
% ID: L12C3B16-66-TN
% Mức: TH
\dangbt{Nhận biết từ thông và tính từ thông qua một diện tích}
\begin{ex}
% ID: L12C3B16-66-TN
% Mức: TH
Một học sinh ôn tập Nhận biết từ thông và tính từ thông qua một diện tích?
\choice
{\True Từ thông qua diện tích $S$ trong từ trường đều là Phi = BS cos(alpha), với alpha là góc giữa $B$ và pháp tuyến.}
{Từ thông luôn bằng BS sin(alpha), với alpha là góc giữa $B$ và pháp tuyến.}
{Có thể bỏ qua phương, chiều và đơn vị trong mọi trường hợp.}
{Kết quả không phụ thuộc dữ kiện và điều kiện áp dụng.}
\loigiai{
Từ thông qua diện tích $S$ trong từ trường đều là Phi = BS cos(alpha), với alpha là góc giữa $B$ và pháp tuyến. Vì vậy phương án $A$ đúng; các phương án còn lại sai.
}
\end{ex}

% ===== Câu 280 =====
% ID: L12C3B16-67-TL
% Mức: TH
\dangbt{Khai thác đồ thị từ thông – thời gian và bài toán cảm ứng điện từ tổng hợp}
\begin{bt}
% ID: L12C3B16-67-TL
% Mức: TH
Trình bày kiến thức cốt lõi của Khai thác đồ thị từ thông – thời gian và bài toán cảm ứng điện từ tổng hợp và nêu điều kiện áp dụng.
\loigiai{
Cần nêu được: Độ lớn suất điện động cảm ứng bằng độ lớn hệ số góc của đồ thị liên kết từ thông theo thời gian. Khi vận dụng phải xác định đúng dữ kiện, phương chiều nếu có, điều kiện áp dụng, hệ đơn vị và kiểm tra tính hợp lí. Nhận định “Đồ thị từ thông nằm ngang ứng với suất điện động cảm ứng khác không.” sai.
}
\end{bt}

% ===== Câu 281 =====
% ID: L12C3B16-67-TLN
% Mức: TH
\begin{ex}
% ID: L12C3B16-67-TLN
% Mức: TH
Trong bốn phát biểu về Khai thác đồ thị từ thông – thời gian và bài toán cảm ứng điện từ tổng hợp, có bao nhiêu phát biểu đúng? (Chỉ ghi một số nguyên.) (1) Độ lớn suất điện động cảm ứng bằng độ lớn hệ số góc của đồ thị liên kết từ thông theo thời gian. (2) Đồ thị từ thông nằm ngang ứng với suất điện động cảm ứng khác không. (3) Cần kiểm tra điều kiện áp dụng. (4) Có thể bỏ qua đơn vị.
\shortans{2}
\loigiai{
Đối chiếu với kiến thức: Độ lớn suất điện động cảm ứng bằng độ lớn hệ số góc của đồ thị liên kết từ thông theo thời gian. Có 2 phát biểu đúng.
}
\end{ex}

% ===== Câu 282 =====
% ID: L12C3B16-67-TN
% Mức: TH
\dangbt{Nhận biết từ thông và tính từ thông qua một diện tích}
\begin{ex}
% ID: L12C3B16-67-TN
% Mức: TH
Chọn phát biểu không trái với kiến thức về Nhận biết từ thông và tính từ thông qua một diện tích?
\choice
{Kết quả không phụ thuộc dữ kiện và điều kiện áp dụng.}
{\True Từ thông qua diện tích $S$ trong từ trường đều là Phi = BS cos(alpha), với alpha là góc giữa $B$ và pháp tuyến.}
{Từ thông luôn bằng BS sin(alpha), với alpha là góc giữa $B$ và pháp tuyến.}
{Có thể bỏ qua phương, chiều và đơn vị trong mọi trường hợp.}
\loigiai{
Từ thông qua diện tích $S$ trong từ trường đều là Phi = BS cos(alpha), với alpha là góc giữa $B$ và pháp tuyến. Vì vậy phương án $B$ đúng; các phương án còn lại sai.
}
\end{ex}

% ===== Câu 283 =====
% ID: L12C3B16-68-TL
% Mức: TH
\dangbt{Khai thác đồ thị từ thông – thời gian và bài toán cảm ứng điện từ tổng hợp}
\begin{bt}
% ID: L12C3B16-68-TL
% Mức: TH
Giải thích vì sao nhận định sau đúng: “Độ lớn suất điện động cảm ứng bằng độ lớn hệ số góc của đồ thị liên kết từ thông theo thời gian.”
\loigiai{
Cần nêu được: Độ lớn suất điện động cảm ứng bằng độ lớn hệ số góc của đồ thị liên kết từ thông theo thời gian. Khi vận dụng phải xác định đúng dữ kiện, phương chiều nếu có, điều kiện áp dụng, hệ đơn vị và kiểm tra tính hợp lí. Nhận định “Đồ thị từ thông nằm ngang ứng với suất điện động cảm ứng khác không.” sai.
}
\end{bt}

% ===== Câu 284 =====
% ID: L12C3B16-68-TLN
% Mức: TH
\begin{ex}
% ID: L12C3B16-68-TLN
% Mức: TH
Trong bốn phát biểu về Khai thác đồ thị từ thông – thời gian và bài toán cảm ứng điện từ tổng hợp, có bao nhiêu phát biểu đúng? (Chỉ ghi một số nguyên.) (1) Đồ thị từ thông nằm ngang ứng với suất điện động cảm ứng khác không. (2) Độ lớn suất điện động cảm ứng bằng độ lớn hệ số góc của đồ thị liên kết từ thông theo thời gian. (3) Kết quả phải phù hợp ý nghĩa vật lí. (4) Mọi đại lượng đều là vô hướng.
\shortans{1}
\loigiai{
Đối chiếu với kiến thức: Độ lớn suất điện động cảm ứng bằng độ lớn hệ số góc của đồ thị liên kết từ thông theo thời gian. Có 1 phát biểu đúng.
}
\end{ex}

% ===== Câu 285 =====
% ID: L12C3B16-68-TN
% Mức: VD
\dangbt{Nhận biết từ thông và tính từ thông qua một diện tích}
\begin{ex}
% ID: L12C3B16-68-TN
% Mức: VD
Cơ sở Vật lí đúng dùng để giải Nhận biết từ thông và tính từ thông qua một diện tích?
\choice
{Có thể bỏ qua phương, chiều và đơn vị trong mọi trường hợp.}
{Kết quả không phụ thuộc dữ kiện và điều kiện áp dụng.}
{\True Từ thông qua diện tích $S$ trong từ trường đều là Phi = BS cos(alpha), với alpha là góc giữa $B$ và pháp tuyến.}
{Từ thông luôn bằng BS sin(alpha), với alpha là góc giữa $B$ và pháp tuyến.}
\loigiai{
Từ thông qua diện tích $S$ trong từ trường đều là Phi = BS cos(alpha), với alpha là góc giữa $B$ và pháp tuyến. Vì vậy phương án $C$ đúng; các phương án còn lại sai.
}
\end{ex}

% ===== Câu 286 =====
% ID: L12C3B16-69-TL
% Mức: VD
\dangbt{Khai thác đồ thị từ thông – thời gian và bài toán cảm ứng điện từ tổng hợp}
\begin{bt}
% ID: L12C3B16-69-TL
% Mức: VD
Phân tích sai lầm trong nhận định: “Đồ thị từ thông nằm ngang ứng với suất điện động cảm ứng khác không.”
\loigiai{
Cần nêu được: Độ lớn suất điện động cảm ứng bằng độ lớn hệ số góc của đồ thị liên kết từ thông theo thời gian. Khi vận dụng phải xác định đúng dữ kiện, phương chiều nếu có, điều kiện áp dụng, hệ đơn vị và kiểm tra tính hợp lí. Nhận định “Đồ thị từ thông nằm ngang ứng với suất điện động cảm ứng khác không.” sai.
}
\end{bt}

% ===== Câu 287 =====
% ID: L12C3B16-69-TLN
% Mức: TH
\begin{ex}
% ID: L12C3B16-69-TLN
% Mức: TH
Trong bốn phát biểu về Khai thác đồ thị từ thông – thời gian và bài toán cảm ứng điện từ tổng hợp, có bao nhiêu phát biểu đúng? (Chỉ ghi một số nguyên.) (1) Độ lớn suất điện động cảm ứng bằng độ lớn hệ số góc của đồ thị liên kết từ thông theo thời gian. (2) Đồ thị từ thông nằm ngang ứng với suất điện động cảm ứng khác không. (3) Cần kiểm tra điều kiện áp dụng. (4) Có thể bỏ qua đơn vị.
\shortans{2}
\loigiai{
Đối chiếu với kiến thức: Độ lớn suất điện động cảm ứng bằng độ lớn hệ số góc của đồ thị liên kết từ thông theo thời gian. Có 2 phát biểu đúng.
}
\end{ex}

% ===== Câu 288 =====
% ID: L12C3B16-69-TN
% Mức: VD
\dangbt{Nhận biết từ thông và tính từ thông qua một diện tích}
\begin{ex}
% ID: L12C3B16-69-TN
% Mức: VD
Trong một bài thực hành về Nhận biết từ thông và tính từ thông qua một diện tích?
\choice
{Từ thông luôn bằng BS sin(alpha), với alpha là góc giữa $B$ và pháp tuyến.}
{Có thể bỏ qua phương, chiều và đơn vị trong mọi trường hợp.}
{Kết quả không phụ thuộc dữ kiện và điều kiện áp dụng.}
{\True Từ thông qua diện tích $S$ trong từ trường đều là Phi = BS cos(alpha), với alpha là góc giữa $B$ và pháp tuyến.}
\loigiai{
Từ thông qua diện tích $S$ trong từ trường đều là Phi = BS cos(alpha), với alpha là góc giữa $B$ và pháp tuyến. Vì vậy phương án $D$ đúng; các phương án còn lại sai.
}
\end{ex}

% ===== Câu 289 =====
% ID: L12C3B16-70-TL
% Mức: VD
\dangbt{Khai thác đồ thị từ thông – thời gian và bài toán cảm ứng điện từ tổng hợp}
\begin{bt}
% ID: L12C3B16-70-TL
% Mức: VD
Đề xuất các bước giải một bài toán thuộc dạng Khai thác đồ thị từ thông – thời gian và bài toán cảm ứng điện từ tổng hợp.
\loigiai{
Cần nêu được: Độ lớn suất điện động cảm ứng bằng độ lớn hệ số góc của đồ thị liên kết từ thông theo thời gian. Khi vận dụng phải xác định đúng dữ kiện, phương chiều nếu có, điều kiện áp dụng, hệ đơn vị và kiểm tra tính hợp lí. Nhận định “Đồ thị từ thông nằm ngang ứng với suất điện động cảm ứng khác không.” sai.
}
\end{bt}

% ===== Câu 290 =====
% ID: L12C3B16-70-TLN
% Mức: VD
\begin{ex}
% ID: L12C3B16-70-TLN
% Mức: VD
Trong bốn phát biểu về Khai thác đồ thị từ thông – thời gian và bài toán cảm ứng điện từ tổng hợp, có bao nhiêu phát biểu đúng? (Chỉ ghi một số nguyên.) (1) Đồ thị từ thông nằm ngang ứng với suất điện động cảm ứng khác không. (2) Độ lớn suất điện động cảm ứng bằng độ lớn hệ số góc của đồ thị liên kết từ thông theo thời gian. (3) Kết quả phải phù hợp ý nghĩa vật lí. (4) Mọi đại lượng đều là vô hướng.
\shortans{2}
\loigiai{
Đối chiếu với kiến thức: Độ lớn suất điện động cảm ứng bằng độ lớn hệ số góc của đồ thị liên kết từ thông theo thời gian. Có 2 phát biểu đúng.
}
\end{ex}

% ===== Câu 291 =====
% ID: L12C3B16-70-TN
% Mức: VD
\dangbt{Nhận biết từ thông và tính từ thông qua một diện tích}
\begin{ex}
% ID: L12C3B16-70-TN
% Mức: VD
Kết luận đúng khi vận dụng Nhận biết từ thông và tính từ thông qua một diện tích?
\choice
{\True Từ thông qua diện tích $S$ trong từ trường đều là Phi = BS cos(alpha), với alpha là góc giữa $B$ và pháp tuyến.}
{Từ thông luôn bằng BS sin(alpha), với alpha là góc giữa $B$ và pháp tuyến.}
{Có thể bỏ qua phương, chiều và đơn vị trong mọi trường hợp.}
{Kết quả không phụ thuộc dữ kiện và điều kiện áp dụng.}
\loigiai{
Từ thông qua diện tích $S$ trong từ trường đều là Phi = BS cos(alpha), với alpha là góc giữa $B$ và pháp tuyến. Vì vậy phương án $A$ đúng; các phương án còn lại sai.
}
\end{ex}

% ===== Câu 292 =====
% ID: L12C3B16-71-TL
% Mức: VD
\dangbt{Khai thác đồ thị từ thông – thời gian và bài toán cảm ứng điện từ tổng hợp}
\begin{bt}
% ID: L12C3B16-71-TL
% Mức: VD
Nêu một tình huống thực tiễn liên quan đến Khai thác đồ thị từ thông – thời gian và bài toán cảm ứng điện từ tổng hợp và giải thích bằng kiến thức Vật lí.
\loigiai{
Cần nêu được: Độ lớn suất điện động cảm ứng bằng độ lớn hệ số góc của đồ thị liên kết từ thông theo thời gian. Khi vận dụng phải xác định đúng dữ kiện, phương chiều nếu có, điều kiện áp dụng, hệ đơn vị và kiểm tra tính hợp lí. Nhận định “Đồ thị từ thông nằm ngang ứng với suất điện động cảm ứng khác không.” sai.
}
\end{bt}

% ===== Câu 293 =====
% ID: L12C3B16-71-TLN
% Mức: VD
\begin{ex}
% ID: L12C3B16-71-TLN
% Mức: VD
Trong bốn phát biểu về Khai thác đồ thị từ thông – thời gian và bài toán cảm ứng điện từ tổng hợp, có bao nhiêu phát biểu đúng? (Chỉ ghi một số nguyên.) (1) Độ lớn suất điện động cảm ứng bằng độ lớn hệ số góc của đồ thị liên kết từ thông theo thời gian. (2) Đồ thị từ thông nằm ngang ứng với suất điện động cảm ứng khác không. (3) Cần kiểm tra điều kiện áp dụng. (4) Có thể bỏ qua đơn vị.
\shortans{2}
\loigiai{
Đối chiếu với kiến thức: Độ lớn suất điện động cảm ứng bằng độ lớn hệ số góc của đồ thị liên kết từ thông theo thời gian. Có 2 phát biểu đúng.
}
\end{ex}

% ===== Câu 294 =====
% ID: L12C3B16-71-TN
% Mức: NB
\dangbt{Phân tích sự biến thiên từ thông qua khung dây}
\begin{ex}
% ID: L12C3B16-71-TN
% Mức: NB
Phát biểu nào sau đây đúng về Phân tích sự biến thiên từ thông qua khung dây?
\choice
{Từ thông không đổi trong mọi chuyển động của khung dây.}
{Có thể bỏ qua phương, chiều và đơn vị trong mọi trường hợp.}
{Kết quả không phụ thuộc dữ kiện và điều kiện áp dụng.}
{\True Từ thông biến thiên khi $B$, diện tích mạch hoặc góc giữa $B$ và pháp tuyến thay đổi.}
\loigiai{
Từ thông biến thiên khi $B$, diện tích mạch hoặc góc giữa $B$ và pháp tuyến thay đổi. Vì vậy phương án $D$ đúng; các phương án còn lại sai.
}
\end{ex}

% ===== Câu 295 =====
% ID: L12C3B16-72-TL
% Mức: VDC
\dangbt{Khai thác đồ thị từ thông – thời gian và bài toán cảm ứng điện từ tổng hợp}
\begin{bt}
% ID: L12C3B16-72-TL
% Mức: VDC
So sánh nhận định đúng “Độ lớn suất điện động cảm ứng bằng độ lớn hệ số góc của đồ thị liên kết từ thông theo thời gian.” với nhận định sai “Đồ thị từ thông nằm ngang ứng với suất điện động cảm ứng khác không.”; chỉ rõ căn cứ Vật lí.
\loigiai{
Cần nêu được: Độ lớn suất điện động cảm ứng bằng độ lớn hệ số góc của đồ thị liên kết từ thông theo thời gian. Khi vận dụng phải xác định đúng dữ kiện, phương chiều nếu có, điều kiện áp dụng, hệ đơn vị và kiểm tra tính hợp lí. Nhận định “Đồ thị từ thông nằm ngang ứng với suất điện động cảm ứng khác không.” sai.
}
\end{bt}

% ===== Câu 296 =====
% ID: L12C3B16-72-TLN
% Mức: VDC
\begin{ex}
% ID: L12C3B16-72-TLN
% Mức: VDC
Trong bốn phát biểu về Khai thác đồ thị từ thông – thời gian và bài toán cảm ứng điện từ tổng hợp, có bao nhiêu phát biểu đúng? (Chỉ ghi một số nguyên.) (1) Khi giải cần xác định đúng dữ kiện, phương chiều, điều kiện áp dụng và đơn vị. (2) Đồ thị từ thông nằm ngang ứng với suất điện động cảm ứng khác không. (3) Phải dùng đúng định luật cho tình huống. (4) Độ lớn suất điện động cảm ứng bằng độ lớn hệ số góc của đồ thị liên kết từ thông theo thời gian.
\shortans{3}
\loigiai{
Đối chiếu với kiến thức: Độ lớn suất điện động cảm ứng bằng độ lớn hệ số góc của đồ thị liên kết từ thông theo thời gian. Có 3 phát biểu đúng.
}
\end{ex}

% ===== Câu 297 =====
% ID: L12C3B16-72-TN
% Mức: NB
\dangbt{Phân tích sự biến thiên từ thông qua khung dây}
\begin{ex}
% ID: L12C3B16-72-TN
% Mức: NB
Trong nội dung Phân tích sự biến thiên từ thông qua khung dây?
\choice
{\True Từ thông biến thiên khi $B$, diện tích mạch hoặc góc giữa $B$ và pháp tuyến thay đổi.}
{Từ thông không đổi trong mọi chuyển động của khung dây.}
{Có thể bỏ qua phương, chiều và đơn vị trong mọi trường hợp.}
{Kết quả không phụ thuộc dữ kiện và điều kiện áp dụng.}
\loigiai{
Từ thông biến thiên khi $B$, diện tích mạch hoặc góc giữa $B$ và pháp tuyến thay đổi. Vì vậy phương án $A$ đúng; các phương án còn lại sai.
}
\end{ex}

% ===== Câu 298 =====
% ID: L12C3B16-73-TN
% Mức: NB
\begin{ex}
% ID: L12C3B16-73-TN
% Mức: NB
Tình huống 3: chọn kết luận chính xác về Phân tích sự biến thiên từ thông qua khung dây?
\choice
{Kết quả không phụ thuộc dữ kiện và điều kiện áp dụng.}
{\True Từ thông biến thiên khi $B$, diện tích mạch hoặc góc giữa $B$ và pháp tuyến thay đổi.}
{Từ thông không đổi trong mọi chuyển động của khung dây.}
{Có thể bỏ qua phương, chiều và đơn vị trong mọi trường hợp.}
\loigiai{
Từ thông biến thiên khi $B$, diện tích mạch hoặc góc giữa $B$ và pháp tuyến thay đổi. Vì vậy phương án $B$ đúng; các phương án còn lại sai.
}
\end{ex}

% ===== Câu 299 =====
% ID: L12C3B16-74-TN
% Mức: TH
\begin{ex}
% ID: L12C3B16-74-TN
% Mức: TH
Nội dung nào mô tả đúng nhất Phân tích sự biến thiên từ thông qua khung dây?
\choice
{Có thể bỏ qua phương, chiều và đơn vị trong mọi trường hợp.}
{Kết quả không phụ thuộc dữ kiện và điều kiện áp dụng.}
{\True Từ thông biến thiên khi $B$, diện tích mạch hoặc góc giữa $B$ và pháp tuyến thay đổi.}
{Từ thông không đổi trong mọi chuyển động của khung dây.}
\loigiai{
Từ thông biến thiên khi $B$, diện tích mạch hoặc góc giữa $B$ và pháp tuyến thay đổi. Vì vậy phương án $C$ đúng; các phương án còn lại sai.
}
\end{ex}

% ===== Câu 300 =====
% ID: L12C3B16-75-TN
% Mức: TH
\begin{ex}
% ID: L12C3B16-75-TN
% Mức: TH
Khi phân tích Phân tích sự biến thiên từ thông qua khung dây?
\choice
{Từ thông không đổi trong mọi chuyển động của khung dây.}
{Có thể bỏ qua phương, chiều và đơn vị trong mọi trường hợp.}
{Kết quả không phụ thuộc dữ kiện và điều kiện áp dụng.}
{\True Từ thông biến thiên khi $B$, diện tích mạch hoặc góc giữa $B$ và pháp tuyến thay đổi.}
\loigiai{
Từ thông biến thiên khi $B$, diện tích mạch hoặc góc giữa $B$ và pháp tuyến thay đổi. Vì vậy phương án $D$ đúng; các phương án còn lại sai.
}
\end{ex}

% ===== Câu 301 =====
% ID: L12C3B16-76-TN
% Mức: TH
\begin{ex}
% ID: L12C3B16-76-TN
% Mức: TH
Một học sinh ôn tập Phân tích sự biến thiên từ thông qua khung dây?
\choice
{\True Từ thông biến thiên khi $B$, diện tích mạch hoặc góc giữa $B$ và pháp tuyến thay đổi.}
{Từ thông không đổi trong mọi chuyển động của khung dây.}
{Có thể bỏ qua phương, chiều và đơn vị trong mọi trường hợp.}
{Kết quả không phụ thuộc dữ kiện và điều kiện áp dụng.}
\loigiai{
Từ thông biến thiên khi $B$, diện tích mạch hoặc góc giữa $B$ và pháp tuyến thay đổi. Vì vậy phương án $A$ đúng; các phương án còn lại sai.
}
\end{ex}

% ===== Câu 302 =====
% ID: L12C3B16-77-TN
% Mức: TH
\begin{ex}
% ID: L12C3B16-77-TN
% Mức: TH
Chọn phát biểu không trái với kiến thức về Phân tích sự biến thiên từ thông qua khung dây?
\choice
{Kết quả không phụ thuộc dữ kiện và điều kiện áp dụng.}
{\True Từ thông biến thiên khi $B$, diện tích mạch hoặc góc giữa $B$ và pháp tuyến thay đổi.}
{Từ thông không đổi trong mọi chuyển động của khung dây.}
{Có thể bỏ qua phương, chiều và đơn vị trong mọi trường hợp.}
\loigiai{
Từ thông biến thiên khi $B$, diện tích mạch hoặc góc giữa $B$ và pháp tuyến thay đổi. Vì vậy phương án $B$ đúng; các phương án còn lại sai.
}
\end{ex}

% ===== Câu 303 =====
% ID: L12C3B16-78-TN
% Mức: VD
\begin{ex}
% ID: L12C3B16-78-TN
% Mức: VD
Cơ sở Vật lí đúng dùng để giải Phân tích sự biến thiên từ thông qua khung dây?
\choice
{Có thể bỏ qua phương, chiều và đơn vị trong mọi trường hợp.}
{Kết quả không phụ thuộc dữ kiện và điều kiện áp dụng.}
{\True Từ thông biến thiên khi $B$, diện tích mạch hoặc góc giữa $B$ và pháp tuyến thay đổi.}
{Từ thông không đổi trong mọi chuyển động của khung dây.}
\loigiai{
Từ thông biến thiên khi $B$, diện tích mạch hoặc góc giữa $B$ và pháp tuyến thay đổi. Vì vậy phương án $C$ đúng; các phương án còn lại sai.
}
\end{ex}

% ===== Câu 304 =====
% ID: L12C3B16-79-TN
% Mức: VD
\begin{ex}
% ID: L12C3B16-79-TN
% Mức: VD
Trong một bài thực hành về Phân tích sự biến thiên từ thông qua khung dây?
\choice
{Từ thông không đổi trong mọi chuyển động của khung dây.}
{Có thể bỏ qua phương, chiều và đơn vị trong mọi trường hợp.}
{Kết quả không phụ thuộc dữ kiện và điều kiện áp dụng.}
{\True Từ thông biến thiên khi $B$, diện tích mạch hoặc góc giữa $B$ và pháp tuyến thay đổi.}
\loigiai{
Từ thông biến thiên khi $B$, diện tích mạch hoặc góc giữa $B$ và pháp tuyến thay đổi. Vì vậy phương án $D$ đúng; các phương án còn lại sai.
}
\end{ex}

% ===== Câu 305 =====
% ID: L12C3B16-80-TN
% Mức: VD
\begin{ex}
% ID: L12C3B16-80-TN
% Mức: VD
Kết luận đúng khi vận dụng Phân tích sự biến thiên từ thông qua khung dây?
\choice
{\True Từ thông biến thiên khi $B$, diện tích mạch hoặc góc giữa $B$ và pháp tuyến thay đổi.}
{Từ thông không đổi trong mọi chuyển động của khung dây.}
{Có thể bỏ qua phương, chiều và đơn vị trong mọi trường hợp.}
{Kết quả không phụ thuộc dữ kiện và điều kiện áp dụng.}
\loigiai{
Từ thông biến thiên khi $B$, diện tích mạch hoặc góc giữa $B$ và pháp tuyến thay đổi. Vì vậy phương án $A$ đúng; các phương án còn lại sai.
}
\end{ex}

% ===== Câu 306 =====
% ID: L12C3B16-81-TN
% Mức: NB
\dangbt{Nhận biết hiện tượng cảm ứng điện từ và điều kiện xuất hiện dòng điện cảm ứng}
\begin{ex}
% ID: L12C3B16-81-TN
% Mức: NB
Phát biểu nào sau đây đúng về Nhận biết hiện tượng cảm ứng điện từ và điều kiện xuất hiện dòng điện cảm ứng?
\choice
{Chỉ cần đặt mạch kín đứng yên trong từ trường đều không đổi là luôn có dòng điện cảm ứng.}
{Có thể bỏ qua phương, chiều và đơn vị trong mọi trường hợp.}
{Kết quả không phụ thuộc dữ kiện và điều kiện áp dụng.}
{\True Dòng điện cảm ứng xuất hiện trong mạch kín khi từ thông qua mạch biến thiên.}
\loigiai{
Dòng điện cảm ứng xuất hiện trong mạch kín khi từ thông qua mạch biến thiên. Vì vậy phương án $D$ đúng; các phương án còn lại sai.
}
\end{ex}

% ===== Câu 307 =====
% ID: L12C3B16-82-TN
% Mức: NB
\begin{ex}
% ID: L12C3B16-82-TN
% Mức: NB
Trong nội dung Nhận biết hiện tượng cảm ứng điện từ và điều kiện xuất hiện dòng điện cảm ứng?
\choice
{\True Dòng điện cảm ứng xuất hiện trong mạch kín khi từ thông qua mạch biến thiên.}
{Chỉ cần đặt mạch kín đứng yên trong từ trường đều không đổi là luôn có dòng điện cảm ứng.}
{Có thể bỏ qua phương, chiều và đơn vị trong mọi trường hợp.}
{Kết quả không phụ thuộc dữ kiện và điều kiện áp dụng.}
\loigiai{
Dòng điện cảm ứng xuất hiện trong mạch kín khi từ thông qua mạch biến thiên. Vì vậy phương án $A$ đúng; các phương án còn lại sai.
}
\end{ex}

% ===== Câu 308 =====
% ID: L12C3B16-83-TN
% Mức: NB
\begin{ex}
% ID: L12C3B16-83-TN
% Mức: NB
Tình huống 3: chọn kết luận chính xác về Nhận biết hiện tượng cảm ứng điện từ và điều kiện xuất hiện dòng điện cảm ứng?
\choice
{Kết quả không phụ thuộc dữ kiện và điều kiện áp dụng.}
{\True Dòng điện cảm ứng xuất hiện trong mạch kín khi từ thông qua mạch biến thiên.}
{Chỉ cần đặt mạch kín đứng yên trong từ trường đều không đổi là luôn có dòng điện cảm ứng.}
{Có thể bỏ qua phương, chiều và đơn vị trong mọi trường hợp.}
\loigiai{
Dòng điện cảm ứng xuất hiện trong mạch kín khi từ thông qua mạch biến thiên. Vì vậy phương án $B$ đúng; các phương án còn lại sai.
}
\end{ex}

% ===== Câu 309 =====
% ID: L12C3B16-84-TN
% Mức: TH
\begin{ex}
% ID: L12C3B16-84-TN
% Mức: TH
Nội dung nào mô tả đúng nhất Nhận biết hiện tượng cảm ứng điện từ và điều kiện xuất hiện dòng điện cảm ứng?
\choice
{Có thể bỏ qua phương, chiều và đơn vị trong mọi trường hợp.}
{Kết quả không phụ thuộc dữ kiện và điều kiện áp dụng.}
{\True Dòng điện cảm ứng xuất hiện trong mạch kín khi từ thông qua mạch biến thiên.}
{Chỉ cần đặt mạch kín đứng yên trong từ trường đều không đổi là luôn có dòng điện cảm ứng.}
\loigiai{
Dòng điện cảm ứng xuất hiện trong mạch kín khi từ thông qua mạch biến thiên. Vì vậy phương án $C$ đúng; các phương án còn lại sai.
}
\end{ex}

% ===== Câu 310 =====
% ID: L12C3B16-85-TN
% Mức: TH
\begin{ex}
% ID: L12C3B16-85-TN
% Mức: TH
Khi phân tích Nhận biết hiện tượng cảm ứng điện từ và điều kiện xuất hiện dòng điện cảm ứng?
\choice
{Chỉ cần đặt mạch kín đứng yên trong từ trường đều không đổi là luôn có dòng điện cảm ứng.}
{Có thể bỏ qua phương, chiều và đơn vị trong mọi trường hợp.}
{Kết quả không phụ thuộc dữ kiện và điều kiện áp dụng.}
{\True Dòng điện cảm ứng xuất hiện trong mạch kín khi từ thông qua mạch biến thiên.}
\loigiai{
Dòng điện cảm ứng xuất hiện trong mạch kín khi từ thông qua mạch biến thiên. Vì vậy phương án $D$ đúng; các phương án còn lại sai.
}
\end{ex}

% ===== Câu 311 =====
% ID: L12C3B16-86-TN
% Mức: TH
\begin{ex}
% ID: L12C3B16-86-TN
% Mức: TH
Một học sinh ôn tập Nhận biết hiện tượng cảm ứng điện từ và điều kiện xuất hiện dòng điện cảm ứng?
\choice
{\True Dòng điện cảm ứng xuất hiện trong mạch kín khi từ thông qua mạch biến thiên.}
{Chỉ cần đặt mạch kín đứng yên trong từ trường đều không đổi là luôn có dòng điện cảm ứng.}
{Có thể bỏ qua phương, chiều và đơn vị trong mọi trường hợp.}
{Kết quả không phụ thuộc dữ kiện và điều kiện áp dụng.}
\loigiai{
Dòng điện cảm ứng xuất hiện trong mạch kín khi từ thông qua mạch biến thiên. Vì vậy phương án $A$ đúng; các phương án còn lại sai.
}
\end{ex}

% ===== Câu 312 =====
% ID: L12C3B16-87-TN
% Mức: TH
\begin{ex}
% ID: L12C3B16-87-TN
% Mức: TH
Chọn phát biểu không trái với kiến thức về Nhận biết hiện tượng cảm ứng điện từ và điều kiện xuất hiện dòng điện cảm ứng?
\choice
{Kết quả không phụ thuộc dữ kiện và điều kiện áp dụng.}
{\True Dòng điện cảm ứng xuất hiện trong mạch kín khi từ thông qua mạch biến thiên.}
{Chỉ cần đặt mạch kín đứng yên trong từ trường đều không đổi là luôn có dòng điện cảm ứng.}
{Có thể bỏ qua phương, chiều và đơn vị trong mọi trường hợp.}
\loigiai{
Dòng điện cảm ứng xuất hiện trong mạch kín khi từ thông qua mạch biến thiên. Vì vậy phương án $B$ đúng; các phương án còn lại sai.
}
\end{ex}

% ===== Câu 313 =====
% ID: L12C3B16-88-TN
% Mức: VD
\begin{ex}
% ID: L12C3B16-88-TN
% Mức: VD
Cơ sở Vật lí đúng dùng để giải Nhận biết hiện tượng cảm ứng điện từ và điều kiện xuất hiện dòng điện cảm ứng?
\choice
{Có thể bỏ qua phương, chiều và đơn vị trong mọi trường hợp.}
{Kết quả không phụ thuộc dữ kiện và điều kiện áp dụng.}
{\True Dòng điện cảm ứng xuất hiện trong mạch kín khi từ thông qua mạch biến thiên.}
{Chỉ cần đặt mạch kín đứng yên trong từ trường đều không đổi là luôn có dòng điện cảm ứng.}
\loigiai{
Dòng điện cảm ứng xuất hiện trong mạch kín khi từ thông qua mạch biến thiên. Vì vậy phương án $C$ đúng; các phương án còn lại sai.
}
\end{ex}

% ===== Câu 314 =====
% ID: L12C3B16-89-TN
% Mức: VD
\begin{ex}
% ID: L12C3B16-89-TN
% Mức: VD
Trong một bài thực hành về Nhận biết hiện tượng cảm ứng điện từ và điều kiện xuất hiện dòng điện cảm ứng?
\choice
{Chỉ cần đặt mạch kín đứng yên trong từ trường đều không đổi là luôn có dòng điện cảm ứng.}
{Có thể bỏ qua phương, chiều và đơn vị trong mọi trường hợp.}
{Kết quả không phụ thuộc dữ kiện và điều kiện áp dụng.}
{\True Dòng điện cảm ứng xuất hiện trong mạch kín khi từ thông qua mạch biến thiên.}
\loigiai{
Dòng điện cảm ứng xuất hiện trong mạch kín khi từ thông qua mạch biến thiên. Vì vậy phương án $D$ đúng; các phương án còn lại sai.
}
\end{ex}

% ===== Câu 315 =====
% ID: L12C3B16-90-TN
% Mức: VD
\begin{ex}
% ID: L12C3B16-90-TN
% Mức: VD
Kết luận đúng khi vận dụng Nhận biết hiện tượng cảm ứng điện từ và điều kiện xuất hiện dòng điện cảm ứng?
\choice
{\True Dòng điện cảm ứng xuất hiện trong mạch kín khi từ thông qua mạch biến thiên.}
{Chỉ cần đặt mạch kín đứng yên trong từ trường đều không đổi là luôn có dòng điện cảm ứng.}
{Có thể bỏ qua phương, chiều và đơn vị trong mọi trường hợp.}
{Kết quả không phụ thuộc dữ kiện và điều kiện áp dụng.}
\loigiai{
Dòng điện cảm ứng xuất hiện trong mạch kín khi từ thông qua mạch biến thiên. Vì vậy phương án $A$ đúng; các phương án còn lại sai.
}
\end{ex}

% ===== Câu 316 =====
% ID: L12C3B16-91-TN
% Mức: NB
\dangbt{Xác định chiều dòng điện cảm ứng bằng định luật Lenz}
\begin{ex}
% ID: L12C3B16-91-TN
% Mức: NB
Phát biểu nào sau đây đúng về Xác định chiều dòng điện cảm ứng bằng định luật Lenz?
\choice
{Dòng điện cảm ứng luôn làm tăng biến thiên từ thông ban đầu.}
{Có thể bỏ qua phương, chiều và đơn vị trong mọi trường hợp.}
{Kết quả không phụ thuộc dữ kiện và điều kiện áp dụng.}
{\True Dòng điện cảm ứng có chiều sao cho từ trường do nó sinh ra chống lại sự biến thiên từ thông đã gây ra nó.}
\loigiai{
Dòng điện cảm ứng có chiều sao cho từ trường do nó sinh ra chống lại sự biến thiên từ thông đã gây ra nó. Vì vậy phương án $D$ đúng; các phương án còn lại sai.
}
\end{ex}

% ===== Câu 317 =====
% ID: L12C3B16-92-TN
% Mức: NB
\begin{ex}
% ID: L12C3B16-92-TN
% Mức: NB
Trong nội dung Xác định chiều dòng điện cảm ứng bằng định luật Lenz?
\choice
{\True Dòng điện cảm ứng có chiều sao cho từ trường do nó sinh ra chống lại sự biến thiên từ thông đã gây ra nó.}
{Dòng điện cảm ứng luôn làm tăng biến thiên từ thông ban đầu.}
{Có thể bỏ qua phương, chiều và đơn vị trong mọi trường hợp.}
{Kết quả không phụ thuộc dữ kiện và điều kiện áp dụng.}
\loigiai{
Dòng điện cảm ứng có chiều sao cho từ trường do nó sinh ra chống lại sự biến thiên từ thông đã gây ra nó. Vì vậy phương án $A$ đúng; các phương án còn lại sai.
}
\end{ex}

% ===== Câu 318 =====
% ID: L12C3B16-93-TN
% Mức: NB
\begin{ex}
% ID: L12C3B16-93-TN
% Mức: NB
Tình huống 3: chọn kết luận chính xác về Xác định chiều dòng điện cảm ứng bằng định luật Lenz?
\choice
{Kết quả không phụ thuộc dữ kiện và điều kiện áp dụng.}
{\True Dòng điện cảm ứng có chiều sao cho từ trường do nó sinh ra chống lại sự biến thiên từ thông đã gây ra nó.}
{Dòng điện cảm ứng luôn làm tăng biến thiên từ thông ban đầu.}
{Có thể bỏ qua phương, chiều và đơn vị trong mọi trường hợp.}
\loigiai{
Dòng điện cảm ứng có chiều sao cho từ trường do nó sinh ra chống lại sự biến thiên từ thông đã gây ra nó. Vì vậy phương án $B$ đúng; các phương án còn lại sai.
}
\end{ex}

% ===== Câu 319 =====
% ID: L12C3B16-94-TN
% Mức: TH
\begin{ex}
% ID: L12C3B16-94-TN
% Mức: TH
Nội dung nào mô tả đúng nhất Xác định chiều dòng điện cảm ứng bằng định luật Lenz?
\choice
{Có thể bỏ qua phương, chiều và đơn vị trong mọi trường hợp.}
{Kết quả không phụ thuộc dữ kiện và điều kiện áp dụng.}
{\True Dòng điện cảm ứng có chiều sao cho từ trường do nó sinh ra chống lại sự biến thiên từ thông đã gây ra nó.}
{Dòng điện cảm ứng luôn làm tăng biến thiên từ thông ban đầu.}
\loigiai{
Dòng điện cảm ứng có chiều sao cho từ trường do nó sinh ra chống lại sự biến thiên từ thông đã gây ra nó. Vì vậy phương án $C$ đúng; các phương án còn lại sai.
}
\end{ex}

% ===== Câu 320 =====
% ID: L12C3B16-95-TN
% Mức: TH
\begin{ex}
% ID: L12C3B16-95-TN
% Mức: TH
Khi phân tích Xác định chiều dòng điện cảm ứng bằng định luật Lenz?
\choice
{Dòng điện cảm ứng luôn làm tăng biến thiên từ thông ban đầu.}
{Có thể bỏ qua phương, chiều và đơn vị trong mọi trường hợp.}
{Kết quả không phụ thuộc dữ kiện và điều kiện áp dụng.}
{\True Dòng điện cảm ứng có chiều sao cho từ trường do nó sinh ra chống lại sự biến thiên từ thông đã gây ra nó.}
\loigiai{
Dòng điện cảm ứng có chiều sao cho từ trường do nó sinh ra chống lại sự biến thiên từ thông đã gây ra nó. Vì vậy phương án $D$ đúng; các phương án còn lại sai.
}
\end{ex}

% ===== Câu 321 =====
% ID: L12C3B16-96-TN
% Mức: TH
\begin{ex}
% ID: L12C3B16-96-TN
% Mức: TH
Một học sinh ôn tập Xác định chiều dòng điện cảm ứng bằng định luật Lenz?
\choice
{\True Dòng điện cảm ứng có chiều sao cho từ trường do nó sinh ra chống lại sự biến thiên từ thông đã gây ra nó.}
{Dòng điện cảm ứng luôn làm tăng biến thiên từ thông ban đầu.}
{Có thể bỏ qua phương, chiều và đơn vị trong mọi trường hợp.}
{Kết quả không phụ thuộc dữ kiện và điều kiện áp dụng.}
\loigiai{
Dòng điện cảm ứng có chiều sao cho từ trường do nó sinh ra chống lại sự biến thiên từ thông đã gây ra nó. Vì vậy phương án $A$ đúng; các phương án còn lại sai.
}
\end{ex}

% ===== Câu 322 =====
% ID: L12C3B16-97-TN
% Mức: TH
\begin{ex}
% ID: L12C3B16-97-TN
% Mức: TH
Chọn phát biểu không trái với kiến thức về Xác định chiều dòng điện cảm ứng bằng định luật Lenz?
\choice
{Kết quả không phụ thuộc dữ kiện và điều kiện áp dụng.}
{\True Dòng điện cảm ứng có chiều sao cho từ trường do nó sinh ra chống lại sự biến thiên từ thông đã gây ra nó.}
{Dòng điện cảm ứng luôn làm tăng biến thiên từ thông ban đầu.}
{Có thể bỏ qua phương, chiều và đơn vị trong mọi trường hợp.}
\loigiai{
Dòng điện cảm ứng có chiều sao cho từ trường do nó sinh ra chống lại sự biến thiên từ thông đã gây ra nó. Vì vậy phương án $B$ đúng; các phương án còn lại sai.
}
\end{ex}

% ===== Câu 323 =====
% ID: L12C3B16-98-TN
% Mức: VD
\begin{ex}
% ID: L12C3B16-98-TN
% Mức: VD
Cơ sở Vật lí đúng dùng để giải Xác định chiều dòng điện cảm ứng bằng định luật Lenz?
\choice
{Có thể bỏ qua phương, chiều và đơn vị trong mọi trường hợp.}
{Kết quả không phụ thuộc dữ kiện và điều kiện áp dụng.}
{\True Dòng điện cảm ứng có chiều sao cho từ trường do nó sinh ra chống lại sự biến thiên từ thông đã gây ra nó.}
{Dòng điện cảm ứng luôn làm tăng biến thiên từ thông ban đầu.}
\loigiai{
Dòng điện cảm ứng có chiều sao cho từ trường do nó sinh ra chống lại sự biến thiên từ thông đã gây ra nó. Vì vậy phương án $C$ đúng; các phương án còn lại sai.
}
\end{ex}

% ===== Câu 324 =====
% ID: L12C3B16-99-TN
% Mức: VD
\begin{ex}
% ID: L12C3B16-99-TN
% Mức: VD
Trong một bài thực hành về Xác định chiều dòng điện cảm ứng bằng định luật Lenz?
\choice
{Dòng điện cảm ứng luôn làm tăng biến thiên từ thông ban đầu.}
{Có thể bỏ qua phương, chiều và đơn vị trong mọi trường hợp.}
{Kết quả không phụ thuộc dữ kiện và điều kiện áp dụng.}
{\True Dòng điện cảm ứng có chiều sao cho từ trường do nó sinh ra chống lại sự biến thiên từ thông đã gây ra nó.}
\loigiai{
Dòng điện cảm ứng có chiều sao cho từ trường do nó sinh ra chống lại sự biến thiên từ thông đã gây ra nó. Vì vậy phương án $D$ đúng; các phương án còn lại sai.
}
\end{ex}
